\documentclass[11pt,letterpaper]{article}
\usepackage[margin=1in]{geometry}
\usepackage[T1]{fontenc}
\usepackage{lmodern,microtype,amsmath,amssymb,amsthm,mathtools}
\usepackage[hidelinks]{hyperref}
\setlength{\emergencystretch}{2em}
\numberwithin{equation}{section}
\newtheorem{theorem}{Theorem}[section]
\newtheorem{proposition}[theorem]{Proposition}
\newtheorem{lemma}[theorem]{Lemma}
\newtheorem{corollary}[theorem]{Corollary}
\theoremstyle{remark}
\newtheorem{remark}[theorem]{Remark}
\DeclareMathOperator{\Ric}{Ric}
\DeclareMathOperator{\Rm}{Rm}
\newcommand{\Eucl}{\delta_{\rm Euc}}
\title{Fixed-type descent for chiral Weyl energy on homotopy four-spheres}
\author{Edward B. Baker III}
\date{Working draft, 25 November 2026}

\begin{document}
\maketitle

\begin{abstract}
We study the infimum of chiral Weyl energy over positive Yamabe
classes on smooth homotopy four-spheres. A common positive Ricci-flow
time for the modified Yamabe representatives gives a universal
recognition gap for the conformal integral \(4e+3\int L^2dV\),
where \(L\) is the largest eigenvalue of either Weyl block. This
also yields a uniform gap above \(4\pi^2\) for the chiral infimum. Throughout the strict subcritical band a positive-Ricci
comparison scale carries absolutely continuous physical Weyl charge.
We locate one smoothly marked disk with positive mass on each side.
For a larger disk with prescribed exterior masses we bound the energy
that any subcritical round-seam repair must remove from the disk.
In the regime of vanishing spectral deficit, every complete smoothly
marked shrinking limit of the modified metrics with strong first-order
convergence is Euclidean. We give a
simultaneous replacement of finitely many marked standard disks
using only first-order comparison data and full-curvature bounds
on their collars. A subcritical path and product-end compactification
give a positive metric on the original smooth type. Its energy
decreases by at least the raw discarded charge minus an aggregate
collar error. For a supplied exterior with exact round product ends,
we give a separate necessary and sufficient boundary form for positive
Yamabe sign. We distinguish joint selection of physical volume and
chiral mass from geometric construction of the collars, and derive
the leading obstruction for controlled boundary-fixed normal-ball
fillings. A whole marked smooth limit leads to closed attainment or a
round restart. An exact evacuated round graft can repeat with unchanged
energy and Yamabe constant. Variational and near-Einstein tests restrict
the closed alternative under their respective hypotheses.
\end{abstract}

\section{The fixed-type variational problem}

Let \(M\) be a closed connected oriented smooth homotopy four-sphere,
and let all metrics be smooth positive-definite Riemannian metrics.
For each chirality, use the Weyl endomorphism on real two-forms and
define its local energy by
\[
 E_\pm(g;U)=\int_U |W^\pm_g|_{\rm HS}^2\,dV_g,\qquad
 E_\pm(g)=E_\pm(g;M).
\]
Our sign convention gives the unit round sphere the identity
curvature operator. Full tensor norms count both antisymmetric
index pairs, so
\(\lvert W\rvert_{\rm tens}^2=4(\lvert W^+\rvert_{\rm HS}^2+
\lvert W^-\rvert_{\rm HS}^2)\).
Signature gives \(E_+(g)=E_-(g)\) on \(M\), since \(\tau(M)=0\)
\cite{Hirzebruch1971Signature}. The two energy measures are separately
conformally invariant, since \(g\mapsto u^2g\) multiplies the Weyl
endomorphism by \(u^{-2}\) and volume by \(u^4\).

Use \(\Delta=\operatorname{div}\nabla\) and let
\[
 Q_g(v)=\int_M(6|dv|_g^2+R_gv^2)\,dV_g,\qquad
 Y([g])=\inf_{v\in H^1(M)\setminus\{0\}}\frac{Q_g(v)}{\|v\|_{4,g}^2}.
\]
Conformal covariance gives \(Q_{u^2g}(v)=Q_g(uv)\) and matches the
denominators. Every positive class has a smooth volume-one Yamabe
minimizer of scalar curvature \(Y([g])\), with \(Y([g])\le
S=8\sqrt6\,\pi\), the round value
\cite[\S\S2--3]{LeeParker1987Yamabe}\cite[Theorem 2]{Schoen1984Yamabe}.
In this paper the \emph{physical} metric is this representative of
the source class.

Set \(T=4\pi^2\) and \(I=I^+(M)=\inf\{E_+(g):Y([g])>0\}\).
A positive class of energy at most \(T\) identifies the standard
smooth type. Indeed, its full Weyl integral is \(8E_+\) and
\(\chi(M)=2\). Chang--Gursky--Yang's strict sphere theorem
gives \(S^4\) or \(\mathbb{RP}^4\), and simple connectivity excludes
the latter. The equality theorem on the remaining types gives the
Fubini--Study class or a manifold covered by \(S^3\times S^1\).
These are excluded by \(b_2(M)=0\) and \(\pi_1(M)=0\),
respectively \cite[Theorems A and B and Remark 2]{ChangGurskyYang2003}.
A nonstandard type with a positive class below \(2T\) therefore
satisfies \(T\le I<2T\). Each positive class has energy strictly
larger than \(T\), with \(I=T\) possible as an unattained endpoint
of the finite-class argument. Theorem~\ref{thm:endpoint-smoothing}
below excludes a uniform interval of such infima.

Assume such a positive entry class exists. Fix one sequence of
physical metrics \(g_i\) with \(e_i=E_+(g_i)\to I\) and
\(e_i\le 8\pi^2-\delta\) for some \(\delta>0\).
Use \(s_i=e_i-I\) for the minimizing excess, \(Y_i=R_{g_i}\) and
\(Z_g=\Ric_g-R_gg/4\). Chern--Gauss--Bonnet and the curvature
decomposition give, with \(\chi(M)=2\),
\begin{equation}\label{eq:budgets}
 \|Z_{g_i}\|_2^2=4e_i+\frac{Y_i^2}{12}-32\pi^2,\qquad
 \sqrt{48\delta}\le Y_i\le S,\qquad
 \int_M |\Rm_{g_i}|_{\rm tens}^2\,dV_{g_i}
   =16e_i+\frac{Y_i^2}{3}-64\pi^2 .
\end{equation}
For the Schouten tensor \(A_g=(\Ric_g-R_gg/6)/2\),
\begin{equation}\label{eq:schouten}
 4\int_M\sigma_2(g_i^{-1}A_{g_i})\,dV_{g_i}=16\pi^2-2e_i>0 .
\end{equation}
Here \(4\sigma_2(g^{-1}A_g)=R_g^2/24-|Z_g|^2/2\).

\begin{lemma}[Positive paths below the scalar-flat threshold]\label{lem:wall}
If \(g_s\), \(0\le s\le1\), is a smooth metric path on \(M\) with
\(Y([g_0])>0\) and \(E_+(g_s)\le b<8\pi^2\), then
\(Y([g_s])\ge\sqrt{48(8\pi^2-b)}\) for every \(s\).
\end{lemma}
\begin{proof}
The first eigenvalue of \(-6\Delta_{g_s}+R_{g_s}\) is continuous by
its compact Rayleigh quotient, and its sign is the Yamabe sign.
If it reached zero, its positive first eigenfunction would give a
scalar-flat metric in that conformal class. Gauss--Bonnet in this
metric gives \(E_+=8\pi^2+\|\Ric\|_2^2/4\), contradicting the
energy bound. In the volume-one Yamabe representative of each
positive class the first equality in \eqref{eq:budgets} proves
the quantitative claim.
\end{proof}

../01_fixed_type_variational_descent/manuscript/physical_volume.tex
\section{Where the Weyl energy can go at the recognition endpoint}
\label{sec:endpoint-charge}

Write \(\lambda_1^\pm\le\lambda_2^\pm\le z_\pm\) for the eigenvalues
of the real Weyl blocks, set \(L=\max(z_+,z_-)\), and write
\(d\lambda^\pm=|W^\pm|^2_{\rm HS}\,dV\). Use the physical
volume-one Yamabe metric for \(D=\int|Z|^2dV\). Define the
conformal invariant
\[
 B=\int_M\left\{\tfrac13\sum_{\pm}(\lambda_1^\pm-\lambda_2^\pm)^2
                   +\min(z_+^2,z_-^2)\right\}dV .
\]
For any positive class on a nonstandard smooth homotopy sphere,
\begin{equation}\label{eq:endpoint-deficit}
                 D+3B < 8(e-T),\qquad e=E_+(g).
\end{equation}
Indeed \(|W^\pm|^2=3z_\pm^2/2+(\lambda_1^\pm-\lambda_2^\pm)^2/2\)
by trace freeness. Hence \(B=4e/3-\int L^2dV\).
Chen and Zhu's strict and equality classifications give
\(Y^2<36\int L^2dV\) on the indicated smooth type
\cite[Theorems 1.4 and 1.6]{ChenZhu2014Conformal}.
For the strict classification simple connectivity excludes every
nontrivial connected-sum factor by van Kampen. At equality the other
possibilities are the projective plane over \(\mathbb C\) and a
universal cover conformal to a product of constant-curvature surfaces.
The cover here is the compact simply connected manifold itself.
Both surface factors would be compact simply connected surfaces, so
the product has nonzero second homology. The projective plane also
has nonzero second homology. Substitution of \eqref{eq:budgets} gives
\eqref{eq:endpoint-deficit}.

For the spectral and curvature estimates suppose \(e_i\to T\)
on hypothetical nonstandard smooth types, which may vary with \(i\).
The subsequent disk selection concerns the fixed-type case \(I=T\).
These are steps under a contradiction hypothesis excluded in
Section~\ref{sec:endpoint-smoothing}. Put
\(y_*=8\sqrt3\,\pi\). Then \(D_i+B_i\to0\),
\(Y_i\to y_*\), and \(6\|L_i\|_{2,g_i}\to y_*\).
On the set where \(z_+\ge z_-\) the negative energy density is
the sum of \(3z_-^2/2\) and half its lower eigengap squared.
On the complementary set interchange the signs. The combined
energy of these two smaller blocks is at most \(3B_i/2\).
The sets in this estimate may vary with \(i\).

\begin{theorem}[Finite disks carrying both chiral energies]
\label{thm:endpoint-disks}
For every such sequence there are smooth positive-Ricci metrics
\(k_i\in[g_i]\), with a constant scale specified below, and a
subsequence having the following property. Given \(\epsilon,\zeta,r_*>0\),
there are finitely many fixed radii \(r_j<r_*\) and, at every
sufficiently large index, pairwise disjoint parametrized smooth
closed standard disks \(B_{ij}\subset B_{k_i}(x_{ij},r_j)\) satisfying
\begin{align}
 \int_{B_{k_i}(x_{ij},r_j)} |\Rm_{g_i}|_{\rm tens}^2\,dV_{g_i}
 &<\zeta\quad\text{for each }j,\label{eq:endpoint-curv}\\
 (\lambda_i^++\lambda_i^-)
 \big(M\setminus\textstyle\bigcup_j B_{ij}\big)&<\epsilon .
 \label{eq:endpoint-capture}
\end{align}
The centers tend to Euclidean regular points in one pointed limit of
the comparison metrics.
\end{theorem}

\begin{proof}
Let \(q_i\) be a volume-one modified Yamabe representative with
\(R_{q_i}-6L_{q_i}=c_i\) constant, where \(c_i\) is the infimum of
the modified Yamabe quotient. Its positive conformal factor is
\(C^{2,\alpha}\) for \(0<\alpha<1\), sufficient for the scalar
equation and test functions below
\cite[Lemma 2.1 and proof of Theorem 1.4]{ChenZhu2014Conformal}.
By their Theorem 1.7, \(c_i\le0\) on this smooth type. Put
\(l_i=\|L_{q_i}\|_{2,q_i}=\|L_{g_i}\|_{2,g_i}\). The
ordinary Yamabe inequality and Cauchy--Schwarz give
\[
 Y_i\le\int R_{q_i}\,dV_{q_i}
      =c_i+6\int L_{q_i}\,dV_{q_i}\le c_i+6l_i.
\]
Thus \(c_i\to0\) and \(0\le6(l_i-\int L_{q_i}\,dV_{q_i})
\le6l_i-Y_i\to0\). Since the volume is one,
\(\|L_{q_i}-l_i\|_2^2=2l_i(l_i-\int L_{q_i}\,dV_{q_i})\).
In particular \(\|R_{q_i}-Y_i\|_{2,q_i}\to0\).
The difference
\[
 \tfrac23(d\lambda_i^++d\lambda_i^-)-L_{q_i}^2\,dV_{q_i}
\]
is a positive measure of mass \(B_i\). The square of \(L_{q_i}\)
converges in \(L^1(q_i)\) to \((y_*/6)^2=4T/3\).
Consequently, for \(\mu_i=dV_{q_i}\),
\begin{equation}\label{eq:charge-volume-tv}
 \left\|\frac{\lambda_i^++\lambda_i^-}{2T}-\mu_i\right\|_{\rm TV}
 \longrightarrow0.
\end{equation}
Here the norm is the mass of the variation measure. Both Weyl
measures are those of the original physical class on its actual sets.
The identity uses the auxiliary volume \(dV_{q_i}\).

The same representative has a further curvature bound. The equation
gives \(R_{q_i}/6\) within \(o(1)\) of \(L_{q_i}\) uniformly and
hence \(R_{q_i}\to y_*\) in \(L^2(q_i)\). Gauss--Bonnet gives
\(\int|Z_{q_i}|^2=4e_i+\int R_{q_i}^2/12-32\pi^2\to0\).
It follows from the curvature decomposition and
\eqref{eq:charge-volume-tv} that
\[
 \big\|\,|\Rm_{q_i}|_{\rm tens}^2-64\pi^2\,\big\|_{L^1(q_i)}
 \longrightarrow0.
\]
Indeed the four-index Weyl contribution tends to \(8T\) in this
norm and the scalar contribution to \(y_*^2/6\). The identities
for \(q_i\) follow by \(C^2\) approximation of its positive
conformal factor on each compact source. If smooth representatives
are needed, choose the approximations separately at each index so
that the errors in these integrals and in volume tend to zero.
These curvature conclusions concern the auxiliary representative;
\eqref{eq:endpoint-curv} below concerns the physical Yamabe metric.

Choose a smooth positive-Ricci comparison in the class as in
Section~\ref{sec:physical-volume}. Normalize its scale by the least
radius carrying half of \(\mu_i\), set this radius to one, and use an
attaining center \(p_i\). This is legitimate since the smooth
comparison distance spheres have zero \(\mu_i\)-measure.
We give the changes to the proof of Theorem~\ref{thm:volume-balls}.
On a cutoff \(\chi\) of support of \(\mu_i\)-measure at most one half,
the Yamabe inequality implies
\[
 \big[(1-2^{-1/2})Y_i-\|R_{q_i}-Y_i\|_{2,q_i}\big]
 \|\chi\|_{4,q_i}^2 \le 6\|d\chi\|_{2,q_i}^2.
\]
The bracket has a uniform positive floor. The conformally invariant
fourth-power gradient estimate in \eqref{eq:physical-small-support}
therefore applies with this floor. The logarithmic cutoffs and
packing argument give \eqref{eq:physical-power} for \(\mu_i\)
and \eqref{eq:physical-noncollapse} for \(k_i\). In particular
the auxiliary volumes are tight and uniformly small on small
comparison balls, with uniform noncollapse of comparison volume
on each bounded window.

Write \(q_i=w_i^2 k_i\). The scalar equation gives
\(-\Delta_{k_i}w_i\le V_iw_i\), where
\(V_i=\tfrac16(R_{q_i})_+ w_i^2\), since \(R_{k_i}>0\).
On any comparison ball \(U\) one has
\[
 6\|V_i\|_{2,k_i;U}\le Y_i\sqrt{\mu_i(U)}
                         +\|R_{q_i}-Y_i\|_{2,q_i}.
\]
On each fixed bounded window choose the radius sufficiently small,
then the index sufficiently large, to absorb the potential term at
\(p=4\) in \eqref{eq:physical-moser}. The resulting local
\(L^8(k_i)\) bound on \(w_i\) uses only
\(\int w_i^4dV_{k_i}=1\). Packing covers the window by uniformly
many such half-balls. Thus the auxiliary volume densities \(w_i^4\)
are bounded in \(L^2(k_i)\) there. Against compactly supported
continuous tests in a common pointed realization, Cauchy--Schwarz
and noncollapsed comparison-volume convergence give
\(\mu=f\mathcal H^4\) with \(f\in L^2_{\rm loc}\) in the limit.
Tightness gives \(\mu(X)=1\).

The original measures \(\nu_i=|\Rm_{g_i}|_{\rm tens}^2dV_{g_i}\)
have bounded total mass by \eqref{eq:budgets}. Pass to a further
subsequence with a local weak limit for these measures. Euclidean
regular points outside the atoms of this limit have full
\(\mu\)-measure. At \(\mu\)-almost every such center the density ratio
\(\mu(B_r(x))/\mathcal H^4(B_r(x))\) tends to \(f(x)\in(0,\infty)\)
as in
\eqref{eq:physical-ordered-limits}. Apply the finite disjoint selection,
shrink and smooth embedding in the proof of
Theorem~\ref{thm:volume-balls} at these centers. For any
\(\delta'>0\) the disks capture more than \(1-\delta'\) of \(\mu_i\),
with the prescribed bound on the original \(\nu_i\) in the containing
balls. Take \(2T\delta'<\epsilon\) and transfer
\eqref{eq:charge-volume-tv} on those same source sets.
\end{proof}

At this endpoint the selection removes the missing complement estimate
for these disks. To use them in Theorem~\ref{thm:simultaneous} one
needs the actual bicollar charts and their aggregate cost. For
example, if the uncovered positive charge is less than \(\epsilon<T\),
a cost \(c_i<T-\epsilon\) yields \(e_i+c_i<2T\) for all large
\(i\), and a positive retained class of energy less than \(T\).
A bound on the curvature integral of a containing ball alone gives
no first-order comparison on the parametrized boundary.

\section{Ricci flow at the recognition endpoint}\label{sec:endpoint-smoothing}

\begin{lemma}[Smoothing from a small curvature excess]\label{lem:endpoint-smoothing}
Let \(h_i\) be smooth Riemannian metrics on closed connected
four-manifolds, with volume one, scalar curvature at least \(-1\),
and a uniform Sobolev inequality
\begin{equation}\label{eq:smoothing-initial-sob}
 \|f\|_4^2\le S_0(\|df\|_2^2+\|f\|_2^2).
\end{equation}
Suppose for some fixed \(K_0>0\) that
\(\|(|\Rm_{h_i}|_{\rm tens}-K_0)_+\|_{2,h_i}\to0\).
The Ricci flows \(h_i(t)\) then exist for a common positive time
\(t_*\) at all sufficiently large indices. They have uniform
curvature bounds on each closed positive-time subinterval. If
\(\|\Ric_{h_i}-\lambda_0 h_i\|_{2,h_i}\to0\) for some \(\lambda_0>0\),
the time can be chosen so that, for
\(\lambda(t)=\lambda_0/(1-2\lambda_0t)\),
\begin{equation}\label{eq:smoothing-einstein-transport}
 \sup_{0\le t\le t_*}\|\Ric_{h_i(t)}-\lambda(t)h_i(t)\|_{2,h_i(t)}\to0,
 \qquad
 \operatorname{Vol}_{h_i(t)}\to (1-2\lambda_0t)^2
\end{equation}
uniformly in time. At any fixed \(0<t<t_*\), a subsequence converges
smoothly modulo diffeomorphisms to a closed Einstein metric with
Einstein constant \(\lambda(t)\).
\end{lemma}

\begin{proof}
Start each maximal smooth compact Ricci flow individually and restrict
estimates to a fixed time horizon at most one. The minimum scalar
curvature cannot decrease. Put \(Q_t(f)=\int(|df|^2+Rf^2/4)dV_t\).
For \(\|f\|_{2,h_i}=1\), Jensen's inequality and
\eqref{eq:smoothing-initial-sob} give
\[
 \int f^2\log f^2\,dV_0\le 2\log\|f\|_4^2
 \le 2\log\big[S_0(Q_0(f)+5/4)\big].
\]
For \(0<\sigma\le5\), the last bound is at most
\(\sigma Q_0(f)-2\log\sigma+C\), by the elementary upper bound of
the logarithm by its tangent. This is a uniform lower bound for
Perelman's \(\mu(h_i,\sigma/4)\). The monotonicity
\(\mu(h_i(t),\tau)\ge\mu(h_i,\tau+t)\) gives the log Sobolev inequality
\(\int f^2\log f^2\,dV_t\le\sigma Q_t(f)-2\log\sigma+C\) for
\(0<\sigma<1\) and \(\|f\|_{2,h_i(t)}=1\), uniformly on the
existing flows \cite[\S3]{Perelman2002}.
Apply Ye's conversion from logarithmic to ordinary Sobolev inequalities
with exponent~\(4\) and potential \(R/4\), whose negative part is bounded
\cite[Theorem~5.5, (5.17)--(5.18)]{YeSobolevFlow}. The result is
\begin{equation}\label{eq:smoothing-weighted-sob}
 \|f\|_{4,t}^2\le A Q_t(f)+B\|f\|_{2,t}^2
\end{equation}
with constants independent of the index and the initial pointwise
curvature maximum.

Write \(U=|\Rm|_{\rm tens}\). The curvature evolution gives
\((\partial_t-\Delta)U\le aU^2\) for a dimensional \(a>0\), in the
weak sense at zeros \cite[Theorem~7.1]{HamiltonSingularities}. Let
\(\phi'=a\phi^2\), \(\phi(0)=K_0\), and choose a common horizon
\(t_*\le 1\) on which \(\phi\le 2K_0\), also before the pole of
\(\lambda\) if the Einstein defect is to be propagated.
For \(v=(U-\phi)_+\),
\[
 (\partial_t-\Delta)v\le a(v^2+2\phi v),\qquad
 R_+\le c_4(\phi+v).
\]
By H\"older, the \(v\) contribution in \eqref{eq:smoothing-weighted-sob}
is bounded by \((Ac_4/4)\|v\|_2\|f\|_4^2\). As long as
\(\|v\|_2\le\delta\) for a sufficiently small uniform \(\delta>0\),
absorption gives \(\|f\|_4^2\le S(\|df\|_2^2+\|f\|_2^2)\) for
all \(f\), with uniform \(S\). Since \(\partial_t dV_t=-R\,dV_t\)
and \(R\ge-1\), testing the subsolution against \(2v\) gives
\[
 \frac{d}{dt}\|v\|_2^2+2\|dv\|_2^2
 \le 2a\int v^3\,dV_t+C\|v\|_2^2.
\]
Here \(\int v^3\le\|v\|_2\|v\|_4^2\). Shrinking \(\delta\)
absorbs the cubic term in the derivative estimate, and Gronwall
keeps \(\|v(t)\|_2\le e^{Ct_*}\|v(0)\|_2<\delta/2\) for all existing
times up to the horizon at every sufficiently large index.

The unweighted Sobolev inequality gives \(\operatorname{Vol}_t
B_t(x,r)\ge b r^4\) for \(0<r\le1\), uniformly. To see this, test
on a cutoff equal to one on the half-ball with Lipschitz constant
\(2/r\). For \(m(r)=r^{-4}\operatorname{Vol}_t B_t(x,r)\), the
inequality yields \(m(r)\ge c\sqrt{m(r/2)}\), and iteration uses
\(\lim_{s\downarrow0}m(s)=\omega_4\) for an individual smooth metric.
Take a maximum of \(sU(x,s)\) over space and a closed existing time
interval \([0,b_0]\) with \(b_0\le t_*\), and write \(K=U(x,t)\)
at a maximizing point. If \(tK\ge2\), then \(U\le2K\) on the
whole slab \([t-K^{-1},t]\). Rescale it to length one and curvature
at most two. In this normalization the maximum principle applied to
\(s|\nabla\Rm|^2+A_1|\Rm|^2\), where \(s\) is time from the start
of the rescaled slab and \(A_1\) is a sufficiently large dimensional
constant, bounds \(|\nabla\Rm|\) at the final time. Indeed the
curvature derivative inequalities give an upper bound
\((1+C s-2A_1)|\nabla\Rm|^2+C A_1\) for its heat operator
\cite[\S\S7 and 13]{Hamilton1982}. Scaling back yields
\(\|\nabla\Rm(t)\|_\infty\le C K^{3/2}\).
If also \(K\ge\max(1,8K_0)\), then \(U(t)\ge K/2\) and
\(v(t)\ge K/4\) on the \(h_i(t)\)-ball of radius \(cK^{-1/2}\)
about \(x\), for a uniform \(0<c\le1\). Its volume lower bound
gives \(\int v(t)^2dV_t\ge bc^4/16\), contradicting the uniform
smallness. Hence \(t\|U(t)\|_\infty\) is bounded for all existing
times up to the horizon. The compact curvature continuation criterion
extends any flow whose maximal time would be at or below \(t_*\)
\cite[Theorem~8.1]{HamiltonSingularities}.

To propagate the Einstein defect, use an evolving orthonormal frame
to keep the norm fixed under time differentiation. For
\(P=\Ric-\lambda(t)h_i(t)\), curvature contraction and
\(\lambda'=2\lambda^2\) give
\((\partial_t-\Delta)P=2\Rm(P)+2\lambda P\) in this frame, where
\(\Rm(h_i(t))=\Ric\) as a contraction on symmetric forms. Thus
\[
 \frac{d}{dt}\|P\|_2^2+2\|\nabla P\|_2^2
 \le C\|P\|_2^2+C\int v|P|^2\,dV_t.
\]
Use \(|\nabla|P||\le|\nabla P|\), H\"older, and the unweighted
Sobolev bound to absorb the last gradient contribution by taking the
initial excess smaller. The asserted \(L^2\) estimate follows. With
\(V_i(t)=\operatorname{Vol}_{h_i(t)}\), one has
\(V_i\le e^{t_*}\) and
\[
 V_i'+4\lambda V_i=-\int\operatorname{tr}P\,dV_t=o(1)
\]
uniformly in time. The volume conclusion follows by solving this
linear equation with \(V_i(0)=1\).

At any fixed interior positive time, bounded curvature and the
volume lower bound on balls give a uniform injectivity radius
\cite[Theorem~4.7]{CheegerGromovTaylor1982}. A uniform diameter
bound follows by packing disjoint fixed-radius balls along a
minimizing geodesic, since the total volume is bounded. Shi's
positive-time estimates bound all curvature derivatives
\cite{Shi1989Flow}. Smooth compactness in the bounded-diameter case
gives convergence under diffeomorphisms from the closed limit manifold
to the sources \cite{HamiltonCompactness}.
The zero limit of the \(L^2\) defect makes this limit Einstein.
\end{proof}

\begin{theorem}[Exclusion of the unattained endpoint]\label{thm:endpoint-smoothing}
There is \(\eta>0\) such that a closed oriented smooth homotopy
four-sphere admitting a positive Yamabe class with chiral energy
less than \(T+\eta\) is diffeomorphic to the standard sphere. In
particular, \(I^+(M)>T\) on a nonstandard smooth homotopy four-sphere
with a positive Yamabe class. For every fixed \(0<\varepsilon_*<\eta\),
the inclusive bound \(I^+(M)\le T+\varepsilon_*\) recognizes the
standard smooth type, whether or not the infimum is attained.
\end{theorem}

\begin{proof}
If the assertion fails, take nonstandard smooth homotopy spheres
\(M_i\), possibly varying, with positive physical Yamabe metrics
\(g_i\) and \(e_i=E_+(g_i)\downarrow T\) along a subsequence. The
strict deficit \eqref{eq:endpoint-deficit} and the modified Yamabe
calculation in Section~\ref{sec:endpoint-charge}, up to and
including the full-curvature density estimate, apply separately to
each class. They give volume-one auxiliary metrics \(q_i\) with
\(Y_i\to y_*=8\sqrt3\,\pi\), \(R_{q_i}\ge c_i\to0\),
\(\|R_{q_i}-y_*\|_2\to0\), \(\|Z_{q_i}\|_2\to0\), and
\(\||\Rm_{q_i}|_{\rm tens}^2-64\pi^2\|_1\to0\).
Choose smooth positive conformal factors converging in \(C^2\)
separately at each index and rescale back to volume one. Choose their
errors so that all the stated estimates persist, with a uniform
scalar lower bound \(-1\). Call the resulting metrics \(h_i\).
The Yamabe inequality and H\"older give
\[
 (Y_i-\|R_{h_i}-y_*\|_2)\|f\|_4^2
 \le 6\|df\|_2^2+y_*\|f\|_2^2.
\]
Also \((U-8\pi)_+^2\le|U^2-64\pi^2|\), and
\(P_i(0)=\Ric_{h_i}-(y_*/4)h_i\) tends to zero in \(L^2\).
Apply Lemma~\ref{lem:endpoint-smoothing} with \(K_0=8\pi\) and
\(\lambda_0=y_*/4\).

At a fixed positive time the smooth limit is on a manifold
diffeomorphic to \(M_i\) for all sufficiently large indices. Its
scalar curvature is \(4\lambda(t)>0\), and its scalar square integral
is \(16\lambda(t)^2(1-2\lambda_0t)^2=y_*^2\). Its trace-free Ricci
tensor vanishes. The signature and Chern--Gauss--Bonnet identities
in Section~1 therefore give its chiral energy equal to \(T\).
This actual smooth positive class meets the inclusive energy
criterion in Section~1 and identifies the standard smooth sphere,
contradicting the choice of the \(M_i\). If
\(I^+(M)\le T+\varepsilon_*<T+\eta\), some positive class has
energy strictly below \(T+\eta\) by definition of the infimum.
\end{proof}

\section{A uniform isotropic spectral gap}\label{sec:isotropic-gap}

On a smooth closed oriented homotopy four-sphere let \([g]\) be
a positive Yamabe class and take its volume-one Yamabe metric
\(g\). Put \(e=E_+(g)=E_-(g)\), \(Y=Y([g])\),
\(D=\int |Z_g|^2dV_g\), and
\[
 L=\max\{\lambda_{\max}(W^+),\lambda_{\max}(W^-)\},\qquad
 \sigma=\int L^2dV,\qquad \mathcal P([g])=4e+3\sigma .
\]
The integrals defining \(\mathcal P\) are conformally invariant.
Gauss--Bonnet and the trace-free rank-three spectral bounds give
\begin{equation}\label{eq:isotropic-quantity}
 \mathcal P-8T=D+\frac{36\sigma-Y^2}{12},\qquad
 \frac92e\le\mathcal P\le8e .
\end{equation}
In fact, for a block with largest eigenvalue \(z\), its squared
Hilbert--Schmidt norm lies between \(3z^2/2\) and \(6z^2\).
The lower bound uses Cauchy--Schwarz on the two smaller
eigenvalues whose sum is \(-z\). The upper bound uses that each
of those eigenvalues is at most \(z\), with extreme pair \((-2z,z)\).
The bounds imply \(e/6\le\sigma\le4e/3\) by signature.

\begin{theorem}[Recognition with a spectral term]\label{thm:isotropic-gap}
There is a universal \(\kappa>0\) such that every positive Yamabe
class on a nonstandard smooth homotopy four-sphere has
\(\mathcal P\ge8T+\kappa\). For any fixed \(0<\varepsilon<\kappa\),
the bound \(\inf_{Y([g])>0}\mathcal P([g])\le8T+\varepsilon\)
recognizes the standard smooth type, including an unattained
infimum. An empty family has infinite infimum.
\end{theorem}

\begin{proof}
On a nonstandard type \(36\sigma-Y^2>0\) by the strict and
equality classifications \cite[Theorems~1.4 and 1.6]{ChenZhu2014Conformal}.
If the uniform gap failed, choose positive classes on possibly
varying nonstandard types with \(\mathcal P_i\to8T\).
Equation~\eqref{eq:isotropic-quantity} gives
\(D_i\to0\), \(36\sigma_i-Y_i^2\to0\), and
\(\limsup e_i\le16T/9\). The sources are eventually in a fixed
strict band \(e_i\le2T-\delta\) for some \(\delta>0\).
In that band Gauss--Bonnet and the Yamabe upper bound give
\(\sqrt{48\delta}\le Y_i\le8\sqrt6\,\pi\). Pass to a subsequence
with \(Y_i\to y>0\), so that \(6\sqrt{\sigma_i}\to y\).

Take the volume-one modified Yamabe representatives \(q_i\)
as in Section~\ref{sec:endpoint-charge}. They satisfy
\[
 R_{q_i}-6L_{q_i}=c_i\le0,\qquad
 Y_i\le c_i+6\int L_{q_i}dV_{q_i}\le c_i+6\sqrt{\sigma_i}
 \quad\cite[Lemma~2.1 and Theorem~1.7]{ChenZhu2014Conformal}.
\]
Thus \(c_i\to0\) and \(\int L_{q_i}dV_{q_i}-\sqrt{\sigma_i}\to0\).
Since the volume is one, \(L_{q_i}\to y/6\) and \(R_{q_i}\to y\)
in \(L^2(q_i)\). In this gauge Gauss--Bonnet gives
\[
 \int |Z_{q_i}|^2dV_{q_i}
 =D_i+\frac{\int R_{q_i}^2dV_{q_i}-Y_i^2}{12}\longrightarrow0.
\]
Smooth approximations of the conformal factors in \(C^2\),
chosen separately on each source and rescaled to volume one,
give smooth representatives \(h_i\) with these same limits and
\(R_{h_i}-6L_{h_i}=o(1)\) uniformly. In particular \(R_{h_i}\ge-1\)
eventually.

Write \(U_i=|\Rm_{h_i}|_{\rm tens}\) for the full ordered-index
norm. The upper spectral bound and the curvature decomposition give
\[
 U_i\le\sqrt{48}\,L_{h_i}+\sqrt2\,|Z_{h_i}|+|R_{h_i}|/\sqrt6.
\]
For a fixed \(K_0>(\sqrt{48}/6+1/\sqrt6)8\sqrt6\,\pi\) this
implies \(\|(U_i-K_0)_+\|_{2,h_i}\to0\). The Yamabe inequality
gives a uniform Sobolev constant since
\[
 (Y_i-\|R_{h_i}-Y_i\|_2)\|f\|_4^2
 \le6\|df\|_2^2+Y_i\|f\|_2^2.
\]
The initial Einstein defects \(\|\Ric_{h_i}-(y/4)h_i\|_2\)
tend to zero as well. Apply Lemma~\ref{lem:endpoint-smoothing}.
The unnormalized flows exist on a common short interval,
and at any fixed interior positive time a subsequence converges
smoothly to a positive Einstein metric on a closed manifold,
via whole-manifold diffeomorphisms. We use also the lemma's
uniform Sobolev estimate along the flow and the estimate
\(\sup_t\|v_i(t)\|_2\to0\) for its evolving excess
\(v_i=(|\Rm_{h_i(t)}|_{\rm tens}-\phi(t))_+\) with bounded \(\phi\).

For either diagonal chiral curvature-operator block
\(A=W^\pm+RI/12\) let \(a\le b\le d\) be its eigenvalues.
Set \(F_\pm=a+b=R/6-\lambda_{\max}(W^\pm)\). Nonnegative
isotropic curvature is equivalent to \(F_+,F_-\ge0\)
\cite[Section~3, (3.1)--(3.2)]{ChenZhu2014Conformal}.
In moving orthonormal frames Hamilton's block equation is
\((D_t-\Delta)A=2(A^2+2\operatorname{adj}A+BB^*)\), using
the appropriate off-diagonal block or its transpose \(B\)
\cite[Section~2, pp.~157--158]{Hamilton1986}.
The trace on any minimizing two-plane gives the lower heat
bound \(2(a^2+b^2)+4dF_\pm\) for \(F_\pm\) after discarding
the nonnegative off-diagonal term. Hence \(w_\pm=(-F_\pm)_+\)
obeys \((\partial_t-\Delta)w_\pm\le C|\Rm|_{\rm tens}w_\pm\)
weakly. One may smooth \(w_\pm\) as the convex maximum of zero
and the negative two-plane traces in the symmetric endomorphism
fibers. Its nonnegative Hessian has the required Laplacian sign.
The active traces have the asserted bound also at zero, and
the quadratic reaction error tends to zero with the mollifier
scale on each smooth flow.

Initially \(w_\pm=o(1)\) uniformly by the modified scalar
equation. Testing its inequality by \(w_\pm\) using \(R\ge-1\)
and \(|\Rm|_{\rm tens}\le v_i+\phi\) gives
\[
 \frac d{dt}\|w_\pm\|_2^2+2\|dw_\pm\|_2^2
 \le C\|w_\pm\|_2^2+C\|v_i\|_2\|w_\pm\|_4^2.
\]
The uniform Sobolev constant absorbs the last term for large
\(i\), and Gronwall gives \(\sup_t\|w_\pm(t)\|_2\to0\).
Consequently the smooth closed positive-time limit has
nonnegative isotropic curvature.

The positive Einstein limit \(h_\infty\) realizes its class Yamabe
constant. Indeed, if \(h_\infty=u^{-2}\widehat g\) for a Yamabe
minimizer \(\widehat g\), the trace-free Ricci transformation gives
\(u Z_{\widehat g}=-2(\nabla^2_{\widehat g}u)^0\). Integrating its
pairing with the divergence-free \(Z_{\widehat g}\) makes \(\widehat g\)
Einstein. The remaining Hessian is pure trace, so either \(u\) is
constant or Obata's theorem gives the standard round type directly
\cite{Obata1962Sphere}. In the constant case the homothety
preserves the normalized scalar integral. Thus
\[
 Y([h_\infty])^2=\int R_{h_\infty}^2dV_{h_\infty}
 \ge36\int L_{h_\infty}^2dV_{h_\infty}.
\]
At strict inequality Chen--Zhu's classification gives the
standard smooth sphere in the simply connected case. At equality
the extra cases have nonzero second homology, since they are
\(\mathbb{CP}^2\) or have a compact simply connected cover
conformal to a product of two surfaces
\cite[Theorems~1.4 and 1.6]{ChenZhu2014Conformal}. The limit is
therefore standard, as are all sufficiently late sources of the
subsequence. This contradiction gives the universal gap. If
the infimum on any type is at most \(8T+\varepsilon<8T+\kappa\),
an actual class has energy quantity strictly below \(8T+\kappa\)
by the definition of infimum.
\end{proof}

In the notation \(B=4e/3-\sigma\) of Section~\ref{sec:endpoint-charge},
the theorem gives \(8(e-T)\ge3B+\kappa\) on a nonstandard type.
The theorem supplies a non-effective lower bound for
\(D+(36\sigma-Y^2)/12\) throughout the strict band and also for
arbitrary positive classes above the band. The upper bound
\(\mathcal P\le8e\) recovers a uniform recognition interval for
\(e\) alone.

\begin{corollary}[Where a vanishing spectral deficit can persist]
\label{cor:isotropic-ricci-tail}
Suppose \(e_i\le2T-\delta\) with \(\delta>0\) for positive Yamabe
classes on nonstandard types and \(36\sigma_i-Y_i^2\to0\).
Write \(\mu_i=dV_{q_i}\) for the volume measure of the volume-one
modified metric. On a subsequence there are measurable sets \(H_i\) with
\[
 \mu_i(H_i)\to0,\qquad
 \liminf_i\int_{H_i}|Z_{q_i}|^2d\mu_i>0,\qquad
 \int_{H_i}(R_{q_i}^2+|W_{q_i}|_{\rm tens}^2)d\mu_i\to0 .
\]
\end{corollary}
\begin{proof}
Suppose instead the high tails of \(|Z_{q_i}|\) are uniformly
integrable in squared norm, that is,
\(\lim_{K\to\infty}\limsup_i\|(|Z_{q_i}|-K)_+\|_2=0\).
As in the theorem, the spectral deficit gives \(c_i\to0\) as real
numbers, and both \(R_{q_i}-Y_i\) and \(L_{q_i}-Y_i/6\) tend to zero
in \(L^2(q_i)\).
The curvature decomposition therefore gives the same tail condition
for the full curvature norm of the smooth approximations.
We explain why one sufficiently small tail at a sufficiently large
fixed cutoff suffices here. In the proof of
Lemma~\ref{lem:endpoint-smoothing} keep the two coefficients of the
ordinary Sobolev inequality separate. With \(K\ge1\) and
\(\phi\le2K\), absorption in \eqref{eq:smoothing-weighted-sob}
gives a gradient coefficient bounded independently of \(K\) and
a lower-order coefficient at most \(C(1+K)\), as long as
\(\|v_i\|_2\) is below a fixed small constant. In the energy
estimates for \(v_i\) and \(w_\pm\) the gradient absorption
threshold is thus independent of \(K\). The remaining Gronwall
coefficient is at most \(C(1+K)\). On a horizon \(c/(1+K)\)
the growth factor is uniform. The ball cutoff iteration on radii
\(r\le(1+K)^{-1/2}\) gives a uniform volume ratio floor since
\((1+K)r^2\le1\). A curvature record above a sufficiently large
multiple of \(1+K\) uses only these radii, so its lower bound on
the squared excess is also independent of \(K\).
Choose \(K\) once by the tail condition and take late indices.
The common flow and smooth positive-time compactness follow,
and the negative chiral sums still tend to zero in squared norm.
The limit has nonnegative isotropic curvature.

For completeness, a smooth closed homotopy sphere with this curvature
has a positive-isotropic metric. Restart its Ricci flow; the lower
pair-sum inequality preserves both nonnegative sums. The scalar
curvature is nonnegative and somewhere positive, since if it vanished
identically both Weyl blocks would vanish and Gauss--Bonnet would
contradict positive Euler characteristic. At positive time scalar
curvature is positive. If a lower pair sum vanishes there, the
strong minimum principle gives vanishing on all earlier positive
slices. The reaction bound forces its two eigenvalues to be zero.
The block has rank one with positive trace. For a local section of
its kernel, differentiate its annihilation by the block, and test
the block evolution on that section. The sum of the off-diagonal
square and the squared norm of the component of its derivative in
the positive line is zero. Thus the kernel and its orthogonal line
are parallel. Simple connectivity gives a nonzero parallel two-form,
contrary to \(b_2=0\). The strict positive-isotropic classification
now gives the standard sphere \cite[Main Theorem]{ChenTangZhu2012}.
This contradicts the smooth identification of the limit with
late sources.

The tail condition must therefore fail. Taking increasing cutoffs
and a subsequence gives \(H_i=\{|Z_{q_i}|>K_i\}\) with definite
squared integral and \(K_i\to\infty\). The uniformly bounded total
Ricci energy, by Gauss--Bonnet in \(q_i\), gives \(\mu_i(H_i)\to0\).
The scalar and spectral squared-integral limits and
\(|W|_{\rm tens}^2\le48L^2\) prove the last assertion.
\end{proof}

\section{The modified Ricci concentration regime}\label{sec:modified-ricci-concentration}

Suppose \(g_i\) are unit-volume positive Yamabe metrics on possibly
varying nonstandard smooth homotopy four-spheres \(M_i\), with
\(e_i=E_+(g_i)=E_-(g_i)\le2T-\delta\) for fixed \(\delta>0\).
Use \(\sigma_i=\int L_{g_i}^2dV_{g_i}\) as in
Section~\ref{sec:isotropic-gap}. Consider the residual regime
\(36\sigma_i-Y_i^2\to0\) and the actual modified Yamabe metrics
\(q_i\) of volume one in these classes. The following statement
identifies which shrinking comparisons can carry Ricci energy.

\begin{theorem}[Modified shrinking limits]\label{thm:modified-shrinking-limits}
Let \(a_i\to0\). If \(a_i^{-2}q_i\) converge locally uniformly and
strongly in \(W^{1,4}\) to a complete connected smooth Riemannian
four-manifold under oriented smooth embeddings of neighborhoods
of a compact exhaustion, then the limit is Euclidean. The full
curvature integrals on the images of fixed compacts tend to zero.

The same conclusion holds for the complete pointed harmonic
limit obtained if the curvature of \(a_i^{-2}q_i\) is uniformly
bounded on each fixed-radius ball at the chosen centers, with
bounds depending on the radius. In particular, for
\(K_i=\max_{M_i}|\Rm_{q_i}|_{\rm tens}\), one has \(K_i\to\infty\)
and for each fixed \(A>0\),
\begin{equation}\label{eq:modified-peak-energy}
 \sup_{x\in M_i}\int_{B_{q_i}(x,A/\sqrt{K_i})}
              |\Rm_{q_i}|_{\rm tens}^2dV_{q_i}\longrightarrow0.
\end{equation}
\end{theorem}

\begin{proof}
Pass to a subsequence with \(e_i\to e\) and \(Y_i\to y>0\), using
the strict band for the positive lower bound. As in the proof of
Theorem~\ref{thm:isotropic-gap}, \(R_{q_i}=6L_{q_i}+c_i\),
\(c_i\to0\), \(L_{q_i}\to y/6\), and \(R_{q_i}\to y\) in
the auxiliary \(L^2\) norms for the last two statements. Since
\(|W^\pm|_{\rm HS}^2\le6L^2\), the Weyl and scalar squared
integrals tend to zero on sets of vanishing \(q_i\)-volume.
Gauss--Bonnet gives
\begin{equation}\label{eq:modified-total-ricci}
 \int |Z_{q_i}|^2dV_{q_i}\longrightarrow4e+y^2/12-8T.
\end{equation}
There is also an inhomogeneous Sobolev inequality with uniform
constant, by the ordinary Yamabe inequality and the scalar
\(L^2\) estimate. These facts persist on smooth volume-one
conformal approximations of \(q_i\), chosen separately at each
index. We may require relative distance errors and pointwise
errors in curvature norms tending to zero, and \(o(1)\) total
variation error for the squared curvature measures.

For a supplied marked comparison \((X,h)\), the unrescaled
volume of a fixed compact image is \(O(a_i^4)\). Hence both
rescaled Weyl blocks and the scalar curvature tend to zero
in local \(L^2\). The all-lowered curvature formula passes
distributionally under the stated first-order convergence,
since its quadratic first derivatives converge strongly in
\(L^2\). The full curvature is uniformly bounded in \(L^2\) by
\eqref{eq:modified-total-ricci} and the other curvature budgets.
Consequently \(R_h=W_h=0\), and by exhaustion
\(d_X:=\int_X|\Ric_h|^2dV_h\le4e+y^2/12-8T\).
The ordinary Yamabe inequality passes to compactly supported
tests on \(X\), since the compact scalar term tends to zero by
H\"older. Thus \(y\|f\|_4^2\le6\|df\|_2^2\). This excludes
compactness and gives \(\operatorname{Vol}_h B_s(x)\ge bs^4\)
at all radii by the dyadic cutoff estimate.

The exact scalar-flat conformally flat metric has harmonic
curvature and finite \(\int|\Rm|_{\rm tens}^2=2d_X\). With the finite homogeneous Sobolev constant already proved,
the harmonic-curvature case of Tian--Viaclovsky gives finitely many
ends, each asymptotically locally Euclidean (ALE) of every order
less than two \cite[Theorem~5.1]{TianViaclovsky2008Volume}.

Let \(k\) be the number of ends, \(m_j\) their link-group orders,
and \(m=\max_j m_j\). A whole compact truncation embeds in an
actual source. Its first and second rational homology vanish by
Mayer--Vietoris since the end links are rational homology spheres.
The relative fundamental class and Poincar\'e--Lefschetz duality
give third Betti number \(k-1\), hence \(\chi(X)=2-k\).
The ALE Gauss--Bonnet identity in the scalar-flat conformally
flat case now gives the exact cost
\begin{equation}\label{eq:modified-ale-cost}
 d_X=4T\left(k-2+\sum_{j=1}^k m_j^{-1}\right)
 \quad\cite[Section~7, (7.4)]{TianViaclovsky2005ALE}.
\end{equation}
The homogeneous Yamabe inequality on each end gives
\(y\le S/\sqrt{m_j}\), where \(S=8\sqrt6\pi\). Cut the Euclidean
radial sharp test to an annulus and dilate its support out the
end at fixed cutoffs. On the quotient the gradient
integral is divided by \(m_j\) and the squared \(L^4\) norm by
\(\sqrt{m_j}\). Remove the cutoffs after the metric limit. The
Ricci upper budget and \eqref{eq:modified-ale-cost} now give
\[
 \delta+T\left(k-2+\sum_j m_j^{-1}\right)\le 2T/m.
\]
For \(k\ge2\) the left side is strictly larger than \(2T/m\).
Thus \(k=1\). The nonnegativity of \(d_X\) in
\eqref{eq:modified-ale-cost} forces the single end group to be
trivial and \(d_X=0\). The metric is flat with asymptotic volume
ratio one, and Bishop--Gromov rigidity makes it Euclidean.

To get compact energy convergence, put the two Levi-Civita
connections on the same oriented Euclidean bundle using metric
square roots. Their difference \(\Theta_i\) tends to zero in
\(L^4\). The affine path of connections has curvature forms
\(\Omega_{i,t}=(1-t)\Omega+t\Omega_i-t(1-t)\Theta_i\wedge\Theta_i\)
uniformly bounded in \(L^2\). Differentiating the Pfaffian along
the path and using covariant Bianchi gives the difference of
Euler forms as the derivative of a transgression involving
\(2\int_0^1\operatorname{Pf}(\Theta_i,\Omega_{i,t})dt\), with
the common Euler normalization. Its \(L^{4/3}\) norm tends to
zero. Thus the Euler densities converge against compact smooth
cutoffs. Since the limit is flat and the scalar and Weyl
terms vanish in \(L^1\), the Euler identity with coefficient
\(-1/2\) on \(|Z|^2\) gives vanishing trace-free Ricci integrals
on fixed compacts. The full curvature then also tends to zero
in \(L^2\) there.

For the bounded-curvature extraction, the original Sobolev
estimate rescales with a zero-order coefficient \(O(a_i^2)\).
It gives uniform lower volumes on fixed small balls by cutoff
iteration. Together with curvature bounds on enlarged balls
this gives local injectivity and harmonic-radius estimates
\cite[Theorem~4.7]{CheegerGromovTaylor1982}%
\cite[Theorem~0.3]{AndersonCheeger1992Harmonic}.
The harmonic coordinate Ricci equation and the curvature bound
give local \(W^{2,p}\) bounds for finite \(p>4\), by first gaining
bounded first derivatives from a large-exponent harmonic-radius
bound. A diagonal extraction on buffered windows gives a
complete marked limit in \(C^{1,\beta}\) and weak \(W^{2,p}\),
with the images covering each fixed smaller source ball
\cite[Theorem~3.1 and proof]{BelegradekKapovitch2000Finiteness}.
The scalar and Weyl limits vanish by the previous distributional
argument. For this limit with bounded curvature the weak
curvature identities give \(\Delta_h\Ric_h=\Rm_h*\Ric_h\).
Local tensor elliptic regularity and the harmonic metric equation
bootstrap to smoothness in the same charts
\cite[Lemma~5.1 and local form of Theorem~2.2(ii)]{AvalosCogoRoyoAbrego2026Bach}.
The resulting smooth exhaustion embeddings retain the first-order
convergence. The preceding limit and energy conclusions apply.

If \(K_i\) were bounded on a subsequence, the fixed-cutoff
Ricci tail would vanish and the tail-only isotropic recognition
argument of Section~\ref{sec:isotropic-gap} would identify the
standard type. Thus \(K_i\to\infty\). At scale \(K_i^{-1/2}\)
the full rescaled curvature is uniformly bounded globally.
Any sequence of centers violating \eqref{eq:modified-peak-energy}
would have a buffered harmonic limit with vanishing curvature
integral on the indicated window, a contradiction. The same
statement holds for the unsmoothed \(q_i\) by the relative
distance and squared-curvature measure bounds above.
\end{proof}

If a sequence of balls of radii \(r_i\to0\) carries a fixed
positive modified trace-free Ricci charge, then
\(r_i\sqrt{K_i}\to\infty\) by \eqref{eq:modified-peak-energy}.
Along that sequence the rescaled full curvature must be unbounded
on some fixed enlargement, since uniform bounds on every such
enlargement would give the second extraction in
Theorem~\ref{thm:modified-shrinking-limits}. This is a restriction
on actual charged balls. The loss of uniform integrability alone
gives no such marked ball or bounded-curvature enlargement.

\subsection{A charged ball in the modified metric}
\label{sec:modified-ball-charge}

This comparison tests a particular degeneration of the modified Yamabe
representatives when the spectral deficit tends to zero and a uniform
upper ball-volume ratio is available. The negative Euler trace below
restricts possible initial ends of that Ricci flow. The regularized
flow used for fixed-type attainment has its own compactness argument
\cite{Baker2026BourAttainment}.

\begin{proposition}\label{prop:modified-ball-charge}
In the fixed subcritical band of Section~\ref{sec:modified-ricci-concentration}, let the modified metrics have vanishing spectral deficit on nonstandard types and
suppose additionally that
\[
 \operatorname{Vol}_{q_i} B_{q_i}(x,r)\le V r^4
 \qquad(x\in M_i,\ 0<r\le r_0)
\]
for fixed \(r_0,V>0\). There are \(c>0\), centers \(x_i\), and
radii \(a_i\to0\)
for which
\[
 \int_{B_{q_i}(x_i,a_i)} |Z_{q_i}|^2\,dV_{q_i}\ge c.
\]
On the same balls the scalar and both Weyl square integrals vanish
in the limit, and \(a_i\sqrt{\max|\Rm_{q_i}|}\to\infty\).
\end{proposition}

\begin{proof}
Use the smooth conformal approximations in
Theorem~\ref{thm:modified-shrinking-limits}. They preserve the
volume ratio upper bound with fixed enlarged constants on a smaller
interval, and distort distances by \(1+o(1)\). They have scalar
curvature at least \(-1\), a common inhomogeneous Sobolev bound,
and negative parts of both isotropic pair sums tending uniformly
to zero. We first prove that for a fixed sufficiently small
\(\epsilon>0\) and each fixed \(0<r<\min(1,r_0/2)\), at all late
indices some radius-\(r\) ball of the smoothed metric has full
curvature \(L^2\) norm greater than \(\epsilon\).

Otherwise fix one such radius and rescale by \(r^{-2}\) on the
contrary subsequence. The initial ball volume ratios through unit
scale and Sobolev constants have upper bounds independent of \(r\).
The rescaled zero-order Sobolev coefficient is at most its original
value. Apply Lee--Lee's local smoothing theorem with center arbitrary
and domain parameter larger than the diameter of each rescaled
closed manifold \cite[Theorem~1.1]{LeeLee2026Local}.
The theorem's constants are independent of that domain parameter.
It supplies a closed flow up to a common time \(t_*>0\), which we
take at most \(\min(1,\epsilon^2)\). In dimension four the evolving local bound of Theorem~1.1(3) is \(\Lambda_1(\epsilon+\sqrt{t})\). Curvature is bounded on each positive-time
subinterval and
\[
 \sup_{x,\ 0<t\le t_*}
 \|\Rm(t)\|_{L^2(B_t(x,1/2))}\le C\epsilon .
\]
Here and in choosing \(\epsilon\), the constants in this estimate
are independent of the chosen radius.

For clarity, the potential in the isotropic pair-sum evolution needs
control down to time zero. The weighted Sobolev estimate
\eqref{eq:smoothing-weighted-sob} follows from the initial Sobolev
bound and scalar lower bound exactly as in its proof. Those steps
use normalized test functions and hold at the rescaled volume
\(r^{-4}\), with both constants bounded independently of \(r\).
For a test supported on an evolving ball of radius \(1/2\),
H\"older bounds the scalar potential term by
\(C\epsilon\|f\|_4^2\). Absorption gives an ordinary local
Sobolev inequality with uniform coefficients. Iterating a ball
cutoff gives a uniform lower volume on every evolving ball of
radius \(1/16\). The total rescaled volume stays at most
\(e^{t_*}r^{-4}\), since the scalar curvature stays above \(-1\).

At each time take a maximal \(1/8\)-separated net. Disjoint small
balls bound its cardinality by \(N(r)\), uniformly in the index
and time. Use Lipschitz cutoffs on radius-\(1/4\) balls, equal
to one on the radius-\(1/8\) balls, and normalize by the square
root of their squared sum. The resulting functions satisfy
\(\sum_j\chi_j^2=1\) and \(\sum_j|d\chi_j|^2\le C(r)\).
Apply the local Sobolev estimate and the local curvature bound to
each \(\chi_j f\). The cross terms sum to zero, giving
\[
 \int |\Rm| f^2\,dV_t
 \le C\epsilon\int |df|^2\,dV_t
     +C(r)\int f^2\,dV_t .
\]
Only the zero-order coefficient depends on the covering number.
The cutoffs enter a spatial inequality at a fixed time.

For either negative isotropic pair sum \(w\) from
Section~\ref{sec:isotropic-gap}, test its weak evolution inequality
against \(2w\). Choose \(\epsilon\) once so the gradient term on
the right is absorbed. Since the scalar curvature is bounded
below, Gronwall gives
\(\|w(t)\|_2^2\le e^{C(r)t}\|w(0)\|_2^2=o_i(1)\) at this fixed
radius. The initial squared integral is scale invariant. At any
interior positive time the curvature and derivative bounds, injectivity
bound and total volume upper bound give a uniform diameter bound by packing along minimizing geodesics. Smooth compactness then gives a closed limit under
whole-manifold diffeomorphisms. Both pair sums are nonnegative there.
The positive-time strong maximum principle and \(b_2=0\) argument
in Corollary~\ref{cor:isotropic-ricci-tail} recognize its smooth
homotopy sphere as standard. This contradicts the nonstandard
types on the subsequence.

Now take a slow diagonal of fixed radii tending to zero. Transfer
the charged balls back to \(q_i\) by the distance and curvature
measure comparison, enlarging their radii by a factor tending to
one if necessary. Their volumes tend to zero by the assumed upper
bound. The modified scalar and spectral estimates imply vanishing
scalar and Weyl squared integrals there. The full tensor
decomposition then gives a definite trace-free Ricci integral.
Equation~\eqref{eq:modified-peak-energy} places their radii above
the maximum-curvature scale as claimed.
\end{proof}

Thus a sequence without shrinking balls carrying a fixed modified
Ricci charge cannot have uniformly bounded small-ball upper volume
ratios on any subsequence. The proposition uses neither a physical
chiral quantile nor an embedding of a compact exhaustion into a
charged ball.

\begin{theorem}[A nonflat flow at the modified charge scale]
\label{thm:charged-modified-flow}
Under the same upper volume bound, the radii can be chosen as the
first radii carrying a fixed small full curvature integral on the
smooth modified approximations. The corresponding pointed rescaled
Ricci flows converge on a subsequence, at positive times, to a
complete noncompact four-dimensional flow with nonnegative isotropic
curvature. On a ball of fixed radius about the marked center the
trace-free Ricci energy has a fixed positive lower bound at every
sufficiently small positive time. In particular the flow is nonflat.
\end{theorem}

\begin{proof}
Take the small norm threshold \(\epsilon\) in the preceding proof,
decreasing it below all the dimensional absorption thresholds below.
For each source the maximal ball integral is continuous in the radius
and attained in the center: distance spheres have zero volume, and
the source is smooth and compact. Its value at zero is zero. Let
\(a_i\) be the first radius at which its square root equals
\(\epsilon\), and choose an attaining center \(x_i\). The preceding
fixed-radius assertion gives \(a_i\to0\). Rescale by \(a_i^{-2}\)
and run the same closed flows, writing \(U=|\Rm|_{\rm tens}\).
Every initial unit ball now has curvature norm at most \(\epsilon\);
at \(x_i\) it is exactly \(\epsilon\).
The initial scalar curvature is at least \(-a_i^2\).

We need bounds uniform in this varying scale. Write
\(\|\Ric(t)\|_\infty\le b/t\) with \(b<1/2\), by the
curvature estimate in Lee--Lee Theorem~1.1(1), choosing \(\epsilon\)
and the common horizon sufficiently small. The shrinking balls
estimate gives \(B_t(z,1/2)\subset B_0(z,1)\) on that horizon
\cite[Lemma~3.2]{LeeLee2026Local}. Thus every evolving radius-\(1/2\)
ball has volume at most \(e^{t}V'\), from the original upper bound
and scalar lower bound, for a fixed \(V'\). Together with the
radius-\(1/16\) lower bound from the preceding proof, this bounds
the overlap of the radius-\(1/4\) cutoffs uniformly. For example,
the disjoint small balls about net points within \(1/4\) of a
point are all in its radius-\(1/2\) ball. The previous calculation
therefore gives uniform constants in the form bound
\[
 \int U f^2\,dV_t\le C\epsilon\!
       \left(\int |df|^2\,dV_t+\int f^2\,dV_t\right).
\]
The curvature tensor evolution, Kato's inequality and this bound
with \(f=U\) absorb the cubic term in the derivative of its
squared integral. The initial global full curvature budget then gives
\[
 \sup_{0\le t\le t_*}\int U^2\,dV_t+
 \int_0^{t_*}\!\int |\nabla\Rm|^2\,dV_t\,dt\le C_0 .
\]
The same form bound propagates vanishing of both negative pair sums
with a uniform Gronwall constant. With \(f=1\) it also gives
\(\int |R|\le C\operatorname{Vol}_t\); hence at every fixed time
the rescaled volumes tend to infinity. Curvature and injectivity
bounds, derivative estimates and pointed smooth compactness give
a complete limit flow on positive times. A compact smooth pointed
limit would identify its entire connected source along the
subsequence, contrary to the diverging volume. The propagated
pair-sum estimates give nonnegative isotropic curvature on the limit.

To locate the charge, fix a small terminal time \(\tau\).
All initial unit-ball points lie in \(B_\tau(x_i,D)\) for a
uniform \(D\). Indeed apply the local two-point distance estimate
of Lee--Lee Theorem~1.1 at a chain of centers along an initial
minimizing geodesic of length at most one. Its steps have length
less than the fixed comparison radius \(r_1/2\); each has uniformly
bounded length in the terminal distance. Centered versions of the
theorem concern the same smooth flow by compact uniqueness. Take
a terminal Lipschitz cutoff \(0\le\chi\le1\) equal to one on
this radius-\(D\) ball, zero outside radius \(D+1\), with
\(|d\chi|_{\tau}\le2\). The shrinking balls estimate puts its
entire support inside \(B_0(x_i,D+2)\). This fixed multiple of
the original \(a_i\) has vanishing unscaled modified volume.
Both original Weyl and scalar square integrals there tend to zero.

Normalize the Euler form as
\[
 \mathcal E=\left(|W^+|_{\rm HS}^2+|W^-|_{\rm HS}^2
                 +R^2/24-|Z|^2/2\right)dV .
\]
The decomposition \(U^2=4(|W^+|^2+|W^-|^2)+2|Z|^2+R^2/6\)
on the attaining ball now gives
\(\int\chi\mathcal E(0)\le-\epsilon^2/4+o_i(1)\).
In orthonormal bundles identified by the time-dependent isometries
of the Ricci flow, the time derivative of its metric connection
has norm at most \(C|\nabla\Ric|\): the differentiated Koszul
formula and the derivative of the moving frame each involve one
covariant derivative of \(\Ric\). Differentiating the Euler
polynomial gives an exact form with primitive bounded by
\(C|\nabla\Ric|U\) by the covariant Bianchi identity
\cite[\S3]{ChernSimons1974Characteristic}.
Keep the terminal cutoff fixed as a function of the source point.
Since \(|d\chi|_s\le2(\tau/s)^b\) for \(0<s\le\tau\), integration
by parts and the global curvature and derivative estimates give
\[
 \left|\int\chi\mathcal E(\tau)-\int\chi\mathcal E(0)\right|
 \le C C_0\left(\int_0^\tau(\tau/s)^{2b}\,ds\right)^{1/2}
 \le C'\sqrt\tau .
\]
One can first start at positive time and then use smooth continuity
at zero for each fixed source. Lipschitz tests are admissible in
the spatial integration by parts. The constants and support bound
are uniform in the source and in \(\tau\). Take \(C'\sqrt\tau\le
\epsilon^2/16\). The nonnegative Weyl and scalar terms in the terminal
Euler form force a definite trace-free Ricci integral on
\(B_\tau(x_i,D+1)\). Under complete pointed smooth convergence
every such ball is covered by a slightly larger compact comparison
at this fixed time. This transfers the lower bound to the limit on
a radius-\(D+2\) ball at every sufficiently small positive time.
The scalar curvature is nonnegative by the initial rescaled lower
bound. An interior zero would force it to vanish on all earlier
positive slices by the strong minimum principle. Its evolution
would give zero Ricci tensor there, contradicting the charge
just obtained. Thus the scalar curvature is positive on the flow.
\end{proof}

The model carries curvature after parabolic smoothing. The excluding
Euclidean-limit theorem uses first-order smooth markings for the
original rescalings at time zero. The distance comparison along the
flow alone does not furnish those markings.

\begin{corollary}[The initial end of the charged flow]
\label{cor:charged-flow-initial-end}
On the fixed manifold of the positive-time limit there is no sequence
\(t_j\downarrow0\) along which the metrics converge locally uniformly
and strongly in \(W^{1,4}_{\rm loc}\) to a smooth complete metric.
In particular the limiting flow has no smooth complete time-zero
extension. This assertion uses the markings at positive time only.
\end{corollary}

\begin{proof}
Pass to \(e_i\to e\), \(Y_i\to y>0\) and write
\(d=4e+y^2/12-8T<y^2/12\) for the limiting total modified
trace-free Ricci energy. Fix a common terminal time \(\tau_0\) small
enough for Theorem~\ref{thm:charged-modified-flow}. Write \(h(t)\)
for its limit flow on \(X\). Smooth pointed comparison uses the same
embeddings on closed positive-time subintervals in the flow gauge.
Transfer a fixed compactly supported smooth function \(f\) on \(X\)
to the sources via these embeddings, extending by zero. At \(\tau_0\)
its size, gradient and support radius are bounded uniformly in \(i\).
Its support lies in a fixed-radius initial ball in the rescaling by
the shrinking balls estimate. The rescaled volume there is uniformly
bounded at zero by the upper volume ratio and at all subsequent times
by the scalar lower bound. The rescaled initial scalar and Weyl
\(L^2\) norms on that ball tend to zero. For the transferred test
\(f_i\), \(|df_i|_s\le C_f (\tau_0/s)^b\).

Three estimates transport the initial data through smoothing. First,
\(\partial_t(R\,dV)=(\Delta R+2|\Ric|^2-R^2)dV\).
Testing against \(f_i\), the reaction terms have bounded spatial
integral, and integration by parts in the Laplacian costs at most
\(C_f(\tau_0/s)^b\|\nabla\Rm\|_2\). The initial scalar integral
against the test tends to zero by Cauchy--Schwarz. Sending \(i\) to
infinity at fixed time gives
\[
 \left|\int_X fR_{h(t)}\,dV_{h(t)}\right|
       \le C_f(t+t^{1/2-b}).
\]
Since \(h(t)\) has nonnegative isotropic curvature, both top Weyl
eigenvalues are bounded by \(R/6\). Trace freeness thus bounds both
Weyl norms by a multiple of \(R\). All three curvature terms vanish
locally in \(L^1\) as \(t\downarrow0\), using nonnegative cutoff
tests equal to one on the compact in question.

Second, the derivative of the Dirichlet energy of \(f_i\) is
the integral of \((2\Ric-Rg)(\nabla f_i,\nabla f_i)\).
Its absolute value is at most \(C_f s^{-2b}\) by the support-volume
and curvature bounds. The fourth-power volume integral changes
by at most \(C_ft\). At time zero the original rescaled Yamabe
inequality has constant \(Y_i\), and the scalar term against
\(f_i^2\) tends to zero. As \(2b<1\), both the quadratic energy
and fourth-power integral at time \(t\) differ from their initial
values by errors tending to zero as \(t\downarrow0\), uniformly
in \(i\). This transfers the full constant \(y\) to any proposed
locally uniform positive-definite time-zero limit.

Third, the Euler transport estimate with the gradient measured at
\(\tau_0\) bounds the change through time \(t\le\tau_0\) by
\(C_f t^{1/2-b}\). In the positive-time embeddings the initial
rescaled measures \(|Z_i|^2dV_i\) admit a local weak subsequential
limit \(\mu\), by compact exhaustion and the global bound.
It has total mass at most \(d\). In these same embeddings the
initial Euler forms tend weakly to \(-\mu/2\), since the initial
scalar and Weyl square integrals vanish on every compact support.
Consequently
\[
  \lim_{t\downarrow0}\int_X f\,\mathcal E(h(t))
       =-\tfrac12\int_X f\,d\mu .
\]
Moreover \(\mu\ne0\). The entire charged initial unit ball lies
in the fixed terminal ball \(B_{\tau_0}(x_i,D)\). A smooth limit
cutoff equal to one on a slightly larger closed terminal ball,
transferred to the sources, equals one there for large \(i\).
The initial trace-free Ricci integral there is at least
\(\epsilon^2/2-o_i(1)\). This argument fixes the terminal
marking before taking the time-zero trace.

Now suppose a sequence of times has the asserted limit \(h_0\).
The global curvature energy stays bounded. In fixed charts the
all-lowered curvature identity and strong first-order convergence
pass the curvatures to the smooth limit distributionally. The
local \(L^1\) bounds above give \(R_{h_0}=W_{h_0}=0\).
The transferred Yamabe test becomes
\(y\|f\|_{4,h_0}^2\le6\|df\|_{2,h_0}^2\).
Metric square roots identify the orthonormal bundles, with
connection differences tending to zero in \(L^4\). The Euler
transgression between each slice and \(h_0\) tends to zero on
compact smooth tests by the \(L^2\) curvature bound, as in
Theorem~\ref{thm:modified-shrinking-limits}. Thus the tested
trace gives \(|\Ric_{h_0}|^2dV_{h_0}=\mu\).

A complete such \(h_0\) has harmonic curvature, finite total
curvature and the homogeneous Sobolev inequality. Its finitely
many ALE ends and the end formula of
Theorem~\ref{thm:modified-shrinking-limits} apply. Indeed every
fixed smooth truncation of \(X\) already embeds in the original
homotopy spheres by a positive-time marking, which is all that
the topological part requires. The budget
\(\int|\Ric_{h_0}|^2\le d<y^2/12\) and the radial Yamabe test on
each end exclude two or more ends. The one-end Euler formula
then forces zero Ricci energy, contradicting \(\mu(X)>0\).
\end{proof}

\section{Chiral charge on disks above the recognition threshold}
\label{sec:band-charge}

Here the physical metrics \(g_i\) are volume-one positive Yamabe
minimizers on one nonstandard smooth homotopy sphere. Write
\(e_i=E_+(g_i)=E_-(g_i)\to I<2T\). Let \(q_i\) denote the
volume-one modified Yamabe representative used in
Section~\ref{sec:endpoint-charge}, and let
\(\mu_i=dV_{q_i}\), \(\ell_i^2=\int L^2dV\). The measures
\(\sigma_i=L^2dV\) and \(d\lambda_i^\pm\) are conformally invariant.
In particular \(0<Y_i/6\le\ell_i\le\sqrt{4e_i/3}\).

\begin{theorem}[A partial chiral residual bound]\label{thm:band-charge}
Suppose \(\liminf_i(Y_i-6\ell_i/\sqrt{2})>0\). After passing to
a subsequence, for every \(\epsilon,\zeta,r_*>0\) there are finitely
many fixed radii \(r_j<r_*\) and, at all sufficiently large indices,
disjoint smoothly parametrized closed standard disks with union
\(U_i=\bigcup_j D_{ij}\). Each disk lies in a ball
\(B_{k_i}(x_{ij},r_j)\) of a positive-Ricci metric \(k_i\in[g_i]\)
on which the integral of \(|\Rm_{g_i}|_{\rm tens}^2dV_{g_i}\)
is less than \(\zeta\). On these same disks, for both chiralities,
\begin{equation}\label{eq:band-charge-residual}
 \lambda_i^\pm(U_i)\ge Y_i^2/24-e_i-\epsilon,\qquad
 \lambda_i^\pm(M\setminus U_i)\le 2e_i-Y_i^2/24+\epsilon
                                  \le 4(e_i-T)+\epsilon .
\end{equation}
In particular the hypothesis holds whenever \(I<4T/3\), and each
sign carries at least \(4T-3e_i-\epsilon\) on the disks. If
\(I<5T/4\), their raw complement has energy strictly below \(T\)
for both signs when the tolerance is sufficiently small.
\end{theorem}

\begin{proof}
The inequalities \(Y_i^2\ge48(2T-e_i)\) and
\(18\ell_i^2\le24e_i\) imply the asserted uniform margin for
\(I<4T/3\). Choose the smooth positive-Ricci metric in the class
given in Section~\ref{sec:physical-volume} and normalize its scale
so the least radius enclosing half of \(\sigma_i(M)\) is one.
Denote an attaining center by \(p_i\). The measure is absolutely
continuous on the smooth source, so the ball-mass continuity argument
there applies. Every unit ball and the exterior of \(B_1(p_i)\)
have \(\sigma_i\)-mass at most \(\ell_i^2/2\). For a cutoff
supported in either kind of set the ordinary Yamabe inequality in
\(q_i\), its scalar bound \(R_{q_i}=6L_{q_i}+c_i\le6L_{q_i}\),
and H\"older give
\[
 (Y_i-6\ell_i/\sqrt 2)\|f\|_{4,q_i}^2\le6\|df\|_{2,q_i}^2.
\]
The conformal fourth-power gradient bound of
\eqref{eq:physical-small-support} therefore applies with this floor
and with \(\mu_i\) on its right. Inward logarithmic cutoffs on
unit balls and outward cutoffs based at \(p_i\) give
\eqref{eq:physical-power} for \(\mu_i\). Choose a fixed \(D\) large
enough that \(\mu_i(B_D(p_i))>1/2\). Packing this ball by a bounded
number of quarter-balls finds one of uniformly positive
\(\mu_i\)-mass. Applying the displayed inequality on the
corresponding half-ball with a linear cutoff bounds the
\(k_i\)-volume of \(B_{D+1}(p_i)\) below. Relative volume comparison
now gives \eqref{eq:physical-noncollapse} on bounded windows.

Take a common pointed limit and local weak limits
\(\mu,\sigma,\nu\) for \(\mu_i,\sigma_i\) and
\(\nu_i=|\Rm_{g_i}|^2_{\rm tens}dV_{g_i}\), respectively.
The power estimates give tightness of \(\mu_i\) and
\(\mu(\{x\})=0\) everywhere. For \(q_i=w_i^2k_i\), positivity
of \(R_{k_i}\) gives
\(-\Delta_{k_i}w_i\le V_iw_i\) with \(V_i=L_{q_i}w_i^2\).
On every source ball \(H\), \(\|V_i\|_{2,k_i;H}^2=\sigma_i(H)\).
At each point with \(\sigma(\{x\})=0\) choose a fixed small
continuity radius in a bounded window. On the corresponding balls
the potential term in \eqref{eq:physical-moser} at \(p=4\) is
absorbed. The local \(L^8(k_i)\) bound on \(w_i\) follows from
\(\int w_i^4dV_{k_i}=1\). Thus \(\mu\) is absolutely continuous
relative to \(\mathcal H^4\) on a neighborhood of each such point
by the local \(L^2\) bound on the densities and noncollapsed volume
convergence. A countable subcover and the zero \(\mu\)-mass of the
countably many atoms of \(\sigma\) give \(\mu\ll\mathcal H^4\).
The regular points outside the atoms of \(\nu\) have full
\(\mu\)-measure. The finite disjoint ball selection and smooth
embeddings in Theorem~\ref{thm:volume-balls} apply there by
differentiation of this density. For every fixed \(d>0\) they
give disks with the specified radii and physical curvature ceilings,
and \(\mu_i(M\setminus U_i)<d\) for all large indices.

Finally the constant test in the ordinary Yamabe quotient gives
\(\int L_{q_i}d\mu_i\ge Y_i/6\). On the omitted set this integral
is at most \(\ell_i\sqrt d\). Since the total volume is one,
Cauchy--Schwarz on \(U_i\) and the pointwise spectral bound yield
\[
 (\lambda_i^++\lambda_i^-)(U_i)
 \ge \tfrac32\sigma_i(U_i)
 \ge \tfrac32(Y_i/6-\ell_i\sqrt d)_+^2 .
\]
Either opposite chirality has total mass \(e_i\). Subtract it and
choose \(d\) small using the uniform upper bounds for \(Y_i,\ell_i\).
This proves \eqref{eq:band-charge-residual}.
\end{proof}

These estimates bound the unmodified physical charge of the
complement. A fixed tolerance selects a fixed finite family size.
The strict path budget, the measured allowances of the actual
collars, and the positive sign of a replacement are still needed
when applying the comparison on this family.

\section{A comparison scale carrying diffuse Weyl curvature}
\label{sec:spectral-windows}

The median construction in Theorem~\ref{thm:band-charge} controls
the entire modified volume when its spectral margin is positive.
For descent it suffices to find a single region carrying a fixed
positive chiral mass. We now obtain one throughout the strict band.

\begin{theorem}\label{thm:spectral-windows}
Let \(g_i\) be unit-volume positive Yamabe metrics on a fixed
nonstandard smooth oriented homotopy four-sphere with
\(e_i\to I=I^+(M)<2T\). After passing to a subsequence let
\(Y_i\to y>0\). There are constant rescalings \(k_i\) of smooth
positive-Ricci metrics in these classes and basepoints for which
the pointed limit \(X\) is noncollapsed. Write \(m=\mathcal H^4\)
in Euclidean volume normalization. In this comparison let
\(\sigma,\lambda^+,\lambda^-\) be the local weak limits of
\(L^2dV,|W^+|^2_{\rm HS}dV,|W^-|^2_{\rm HS}dV\) on the sources.
The absolutely continuous parts relative to \(m\) satisfy
\[
 \sigma_{\rm ac}(X)\ge y^2/36,\qquad
 (\lambda^+_{\rm ac}+\lambda^-_{\rm ac})(X)\ge y^2/24 .
\]
For every \(\epsilon,\zeta,r_*>0\) a fixed finite number of
smoothly parametrized disjoint standard disks on each late source,
with union \(U_i\), can be chosen in comparison balls of fixed radii
\(r_j<r_*\), each having original full curvature integral below
\(\zeta\). One sign \(s\), fixed along the subsequence, satisfies
\[
 \lambda_i^s(U_i)\ge Y_i^2/48-\epsilon
                  \ge 2T-e_i-\epsilon.
\]
On the same disks both signs satisfy
\(\lambda_i^\pm(M\setminus U_i)\le 2e_i-Y_i^2/24+\epsilon\).
The second bound is informative for each sign when \(Y_i^2/24>e_i\).
\end{theorem}

\begin{proof}
Use the volume-one modified representatives \(q_i\), with
\(\mu_i=dV_{q_i}\), \(R_{q_i}=6L_{q_i}+c_i\), \(c_i\le0\),
from Section~\ref{sec:endpoint-charge}. The positive-Ricci comparisons
exist by Section~\ref{sec:physical-volume}. Put
\(\sigma_i=L_{q_i}^2\mu_i\), a conformal measure of bounded total
mass. Discard finitely many indices and fix
\(0<6\sqrt\beta<Y_0\le Y_i\) using the strict subcritical budget.
Start with the function \(h_i=1\). At each step it is Lipschitz,
\(0\le h_i\le1\), \(\int h_i^4d\mu_i\to1\), and
\(\int|dh_i|^4dV=o(1)\), the latter integral being conformally
invariant. The restricted mass \(\sigma_i(\{h_i>0\})\) must
exceed \(\beta\) eventually: otherwise the ordinary Yamabe
inequality tested on \(h_i\) gives
\[
 (Y_0-6\sqrt\beta)\|h_i\|_{4,q_i}^2
   \le 6\|dh_i\|_{2,q_i}^2=o(1)
\]
on a subsequence. Here and below we use \(\mu_i(M)=1\).

Scale the positive-Ricci comparison so that the least radius
enclosing \(\beta\) of this restricted measure is one, witnessed
at \(p_i\). The restricted measure is absolutely continuous on
each smooth source, so continuity of ball mass and compactness of
the center set give the normalization. If some bounded centered
window carries a positive lower bound of \(h_i^4\mu_i\) on a
subsequence, keep this scale. Otherwise every fixed window has
weighted volume tending to zero. Choose \(R_i\uparrow\infty\)
sufficiently slowly that \(\int_{B_{R_i^2}(p_i)}h_i^4d\mu_i=o(1)\).
Multiply \(h_i\) by a distance cutoff zero through \(R_i\), rising
linearly in log radius to one at \(R_i^2\). Its fourth-power
gradient integral is at most \(2\pi^2/(\log R_i)^3\) by the
Bishop area bound. The new functions preserve both limits and
their positive sets lose at least \(\beta\) of spectral mass.
Uniform boundedness of \(\sigma_i(M)\) forces termination after
finitely many such deletions.

On the retained scale put \(\theta_i=h_i^4\mu_i\). For a Lipschitz
cutoff \(0\le\chi\le1\) in a unit ball the Yamabe test gives
\[
 (Y_0-6\sqrt\beta)\Big(\int\chi^4d\theta_i\Big)^{1/2}
       \le6\|d(h_i\chi)\|_{2,q_i}^2 .
\]
The norm on the right before squaring is at most
\(o(1)+\|d\chi\|_{4,k_i}\). Inward logarithmic cutoffs give
a uniform upper bound tending to zero for limiting mass on balls
of radius tending to zero. On the other hand a fixed bounded
window carries positive weighted mass. The relative Bishop packing
bound covers it by boundedly many quarter-balls. A linear cutoff on
the half-ball about a quarter-ball of definite mass now gives a
positive lower bound on the comparison volume of that half-ball.
Relative volume comparison gives noncollapse on all bounded
windows. Pass to a pointed limit and common local weak limits of
the measures, including \(\theta_i\), the unrestricted \(\mu_i\)
and the original full curvature measures \(\nu_i\). The preceding
estimates give \(v:=\theta(X)>0\) and \(\theta(\{x\})=0\).

Off the countable atoms of the unrestricted \(\sigma\), choose a
small continuity ball at each point. For \(q_i=w_i^2 k_i\),
\(-\Delta_{k_i}w_i\le L_{q_i}w_i^3\), and the squared local
\(L^2(k_i)\) norm of the potential \(L_{q_i}w_i^2\) is the
original spectral mass on the ball. The \(p=4\) estimate
\eqref{eq:physical-moser} therefore absorbs this term on a sufficiently
small fixed ball and bounds \(w_i^4\) in \(L^2(k_i)\) on its half.
Noncollapsed volume convergence and Cauchy--Schwarz show that the
limit of \(\mu_i\) is absolutely continuous there. Since
\(\theta_i\le\mu_i\), countable covering and the absence of
\(\theta\)-atoms give \(\theta\ll m\) on all of \(X\).
This uses only local convergence.

Set \(\alpha_i=L_{q_i} h_i^2\mu_i\) and take its local weak
limit \(\alpha\). The matrix of measures with diagonal
\((\theta_i,\sigma_i)\) and off-diagonal \(\alpha_i\) is
positive semidefinite: every constant quadratic test has a square
as its density. The property passes to local weak limits. Taking
densities relative to \(\theta+\sigma\) (Cauchy--Schwarz also
gives \(\alpha\ll\theta+\sigma\)) and testing rational pairs yields
\(\alpha=f\theta\) and \(\sigma\ge f^2\theta\).
For \(R>1\), use the centered cutoff \(\chi_R\) equal to one
through \(R\), falling linearly in log radius to zero at \(R^2\).
Testing the Yamabe inequality on \(h_i\chi_R\) gives
\[
 Y_i\Big(\int\chi_R^4d\theta_i\Big)^{1/2}
 \le 6\|d(h_i\chi_R)\|_{2,q_i}^2 +6\int\chi_R^2d\alpha_i .
\]
At fixed \(R\) the radial tests converge on the common pointed
realization. The gradient norm is at most
\(o_i(1)+(2\pi^2)^{1/4}(\log R)^{-3/4}\).
Taking the source limit and then \(R\to\infty\) gives
\(y\sqrt v\le6\alpha(X)\). Cauchy--Schwarz and \(\theta\ll m\)
now give \(\sigma_{\rm ac}(X)\ge\int f^2d\theta
\ge\alpha(X)^2/v\ge y^2/36\). The pointwise algebraic inequality
\(\lambda_i^++\lambda_i^-\ge(3/2)\sigma_i\) passes to measures
and to their absolutely continuous parts.

Take \(s\) maximizing \(\lambda^s_{\rm ac}(X)\).
Both absolutely continuous parts give full measure to regular
points outside atoms of \(\nu\). Apply the finite ball selection
in Section~\ref{sec:physical-volume} to their sum together with
\(\sigma_{\rm ac}\) on a bounded
continuity window containing all but a prescribed tolerance of its
mass. Use differentiation of this density for the loss on shrinking
the balls. The same loss estimate bounds the omitted part of
either absolutely continuous measure. On the sources use the
smooth embedded disks containing the fixed inner comparison balls
and contained in the fixed outer ones. Lower bounds transfer by
weak convergence of the full chiral measures at common continuity
radii; the curvature ceiling uses \(\nu_i\) of \(g_i\).
Subtracting the opposite total mass \(e_i\) from the combined
lower bound proves the two individual residual estimates. Adjust
the tolerance using \(Y_i\to y\) and the budget
\(Y_i^2\ge48(2T-e_i)\).
\end{proof}

There is a second estimate for the full residual. Take the same fixed
disks using the absolutely continuous spectral measure as above.
With \(\Delta_i=36\ell_i^2-Y_i^2\),
\[
   \sigma_i(M\setminus U_i)\le \Delta_i/36+\epsilon,\qquad
   \lambda_i^\pm(M\setminus U_i)\le \Delta_i/6+6\epsilon .
\]
Indeed the transfer captures at least \(Y_i^2/36-\epsilon\) of
\(\sigma_i\), and each Weyl block satisfies \(|W^\pm|^2\le6L^2\).
In particular if \(\Delta_i\to0\), both chiral measures can be
captured to arbitrary tolerance in this one frame. This controls
escape, singular mass and the excluded atoms together. The selected
finite disks and their radii are fixed before the source limit.

For the selected sign, escape together with the entire singular part
in this frame costs at most \(I-y^2/48\). Its raw complement on the
finite family is at most \(2(e_i-T)+\epsilon\), hence is strictly
below \(T\) when \(I<3T/2\) and the tolerance is small. The chiral
density need not be positive in both signs at the same regular point.
Taking a curvature ceiling \(\zeta<2I\) makes every individual disk
carry less than \(I/2\) in either sign. At fixed tolerance the number
of disks is fixed, so a further subsequence has one label carrying a
fixed positive mass of sign \(s\), with a fixed positive mass of the
same sign on its other side. This marked sphere is a candidate for
one repair on the original type.

\begin{corollary}[A marked bicollar with small incoming curvature]
\label{cor:thin-raw-collar}
For every \(\eta,\zeta,r_*>0\) the selected sign admits a fixed
comparison radius \(r<r_*\) and a number \(b>0\) such that on all
late sources a parametrized disk \(D_i\) and a two-sided smooth
bicollar \(A_i\) of its boundary lie in a single comparison ball
\(B_i(r)\). Both complete sides carry at least \(b\) of sign \(s\),
and, writing \(\nu_i=|\Rm_{g_i}|_{\rm tens}^2dV_{g_i}\),
\[
 \nu_i(B_i(r))<\zeta,\qquad
 (\nu_i+\lambda_i^s)(A_i)<\eta\,\lambda_i^s(D_i).
\]
\end{corollary}
\begin{proof}
Choose a regular point of differentiation of \(m,\lambda^s,\nu\)
where \(d\lambda^s_{\rm ac}/dm\) is finite and positive and
\(d\nu_{\rm ac}/dm\) is finite. There the singular parts have zero
four-dimensional density. Choose \(0<\tau<1/6\) so the limiting
density ratio of the mass of \(\nu+\lambda^s\) on
\(B_r\setminus B_{(1-3\tau)r}\) to that of \(\lambda^s\) on
\(B_{(1-3\tau)r}\) is below \(\eta/2\). Then choose a sufficiently
small fixed joint continuity radius meeting the comparison volume
threshold at \((1-\tau)r\) and a strict outer curvature ceiling
below \(\min(\zeta,2I)\). The finite-radius transfers and smooth
open model ball of Section~\ref{sec:physical-volume} apply on
every late source. Its interior contains the closed inner ball
compactly. In its parameter domain choose three nested compact
concentric disks, the smallest containing that compact preimage
in its interior. Use the middle disk for \(D_i\) and the annulus
between the other two for its bicollar. The displayed estimate
follows by inclusion. The fixed inner ball supplies a positive
lower mass, while the entire outer ball carries less than \(I/2\)
of sign \(s\). Taking \(b\) smaller than both margins proves the claim.
\end{proof}

The marking is smooth on both faces. Its normalized distortion and
differentiated coordinate residual can depend on its actual width.
Even a collar of arbitrarily small raw curvature can require a large
change in these quantities. If the entire protected core lies in
the image of one fixed compact of a smooth locally conformally flat
comparison satisfying Theorem~\ref{prop:local-min}, strong
local chiral convergence would give core energy tending to zero.
This excludes a round restart covering the charged core in such
a marking. The evacuated graft of Proposition~\ref{prop:graft}
can still occur on another region.

\subsection{Retaining the marked core inside a larger disk}
\label{sec:balanced-exterior}

\begin{proposition}\label{prop:balanced-exterior}
Let \(Q\) be a parametrized closed four-ball in a connected closed
oriented four-manifold, and let \(m_1,\dots,m_4\) be finite measures
absolutely continuous with respect to smooth volume. Given \(t>0\),
a parametrized standard disk \(B\) containing \(Q\) in its interior
satisfies
\[
 \left|m_j(M\setminus B)-\tfrac12m_j(M\setminus Q)\right|<t
 \qquad (1\le j\le4).
\]
A smooth bicollar of its boundary disjoint from \(Q\) can have
total incoming measure below \(t\) for the sum of the four measures.
\end{proposition}

\begin{proof}
The interior of the exterior of \(Q\) is connected, since paths
can be diverted along the connected boundary of a slightly larger
collared ball. Small coordinate balls there form a Vitali cover.
Select finitely many pairwise disjoint closed balls missing arbitrarily
little exterior mass for the sum of the measures. One may shrink
finitely many open balls slightly since the smooth faces are null.
Join them by a finite tree of smoothly embedded arcs with distinct
transverse endpoints. The complement of the ball interiors is
connected by the same path argument, and general position makes the
arc interiors disjoint. A thin regular neighborhood, rounded at the
feet, is a smooth standard ball \(K\) inside the exterior. Attaching
a tube between two distinct ball components is their boundary
connected sum. Removing leaves of the tree gives a parametrization.
Rounding can retain slightly smaller concentric balls.

Pull the four measures on \(K\) to the closed Euclidean unit ball
and extend them by zero. For \((n,c)\in S^4\subset\mathbb R^4\times
\mathbb R\), integrate \(\operatorname{sign}(n\cdot x+c)\) against
each measure. All affine hyperplanes are null and the signs are
constant near the two poles, so this map to \(\mathbb R^4\) is
continuous and odd. It has a zero by Borsuk--Ulam: a zero-free
normalization would induce a map of real projective spaces from
dimension four to three pulling back the degree-one mod-two generator
to the generator, contrary to their fourth powers. For a nonzero
restricted measure the zero is away from the poles. The resulting
hyperplane bisects all four masses in the chart
\cite{StoneTukey1942Sandwich}. If all masses vanish, any open
half-ball serves. One convex cap with nonempty interior can be
exhausted in measure by smooth closed standard balls inside it:
shrink toward an interior point and approximate the positive radial
function smoothly from below. Choose one such ball \(K'\).

Enlarge \(Q\) slightly in a smooth exterior collar missing \(K'\),
at arbitrarily small added mass. Join the enlargement to \(K'\) by
one embedded arc and round the attaching feet away from \(Q\).
At the fixed source, its tube and the roundings change all four
masses by as little as prescribed, since the arc and faces are null.
This boundary connected sum gives \(B\) and preserves the entire
marking on \(Q\). Shrink a smooth collar of the new boundary using
absolute continuity.
\end{proof}

Use the outer disk of Corollary~\ref{cor:thin-raw-collar} for \(Q_i\),
with its original inner disk and bicollar. Take the two chiral
measures \(\lambda_i^\pm\), the conformal spectral measure
\(\sigma_i=L_i^2dV\), and the full physical measure
\(\nu_i=|\Rm_{g_i}|_{\rm tens}^2dV_{g_i}\). Since \(\nu_i(Q_i)<\zeta\),
its individual chiral masses are below \(\zeta/4\).
Write \(V_i=M\setminus\operatorname{int} B_i\) and
\(p_i=2(\lambda_i^++\lambda_i^-)+3\sigma_i\).
On the closed source \(p_i(M)=\mathcal P_i\) by signature.
For any prescribed \(t_i>0\),
\[
 \left|\lambda_i^s(V_i)-\frac{e_i-\lambda_i^s(Q_i)}2\right|<t_i,
 \qquad
 \left|p_i(V_i)-\frac{\mathcal P_i-p_i(Q_i)}2\right|<7t_i .
\]
With a small fixed \(\zeta\) and \(t_i\to0\), both complete sides
carry a fixed positive amount of each chirality. For \(e_i\to I<2T\)
the raw exterior falls below \(T\) in both signs. Also
\(\mathcal P_i\le8e_i\), so \(p_i(V_i)<8T\) with a fixed margin.
The larger disk can extend beyond the original comparison balls.
Its new collar has its own marking, width and coordinate estimates.

For a positive capped exterior \((\widehat V_i,\widehat g_i)\) use
the disk parametrization to identify it smoothly with \(M\). Define
\(c_i^s=E_s(\widehat g_i)-\lambda_i^s(V_i)\) and
\(c_i^p=\mathcal P([\widehat g_i])-p_i(V_i)\), including every
change on the retained region and the entire filling. Either of
the following inclusive tests then gives smooth recognition:
\[
 c_i^s+\tfrac12(e_i-\lambda_i^s(Q_i))+t_i\le T+\varepsilon_*,
 \qquad
 c_i^p+\tfrac12(\mathcal P_i-p_i(Q_i))+7t_i\le8T+\eta .
\]
When positivity is obtained by the joined path and cylindrical
compactification, its energy ceiling is strictly below \(2T\)
and its comparison includes the matched ends. Separate verification
of the capped ordinary Yamabe sign also suffices.

\paragraph{Spectral cost on the retained marking.}
Use the common density frames of the first-order splice and put
\(L(F)=\max_s\lambda_{\max}(F_s)\). For a unit vector \(u\),
\(\|uu^{\mathsf T}-\mathrm{Id}/3\|_{\rm HS}=\sqrt{2/3}\).
Thus the variational formula for the largest eigenvalue gives
\[
 |L(F')-L(F)|\le\sqrt{2/3}\max_s\|F'_s-F_s\|_{\rm HS}.
\]
Suppose on the actual exterior band \(U\) that the repaired
densities satisfy \(\|F'_s-\theta F_s\|_{2,U}\le b_s\) with
the chosen cutoff \(0\le\theta\le1\). The vector
\((\sqrt2 F_+,\sqrt2 F_-,\sqrt3 L(F))\) has squared norm equal
to the density of \(p\). Positive homogeneity and the preceding
inequality bound its \(L^2\) error relative to the cutoff vector by
\(2\sqrt{b_+^2+b_-^2}\). Thus the signed increment on this band
is at most
\[
 \left(\sqrt{\int_U\theta^2\,dp}
             +2\sqrt{b_+^2+b_-^2}\right)^2-p(U).
\]
Sum over disjoint actual bands and add the costs of all other
changed regions and the new cap. An exactly conformally flat added
region has zero spectral cost. The first-order bounds at the
selected widths and the capped scalar form require separate control.

\begin{corollary}[Energy that must be removed before a large disk can be capped]
\label{cor:dyadic-repair}
For every fixed integer \(k\ge1\), the marked enlargement can be
repeated \(k\) times with nested disks. The final exterior satisfies
\(m(V_i)=2^{-k}m(M\setminus Q_i)+o(1)\) for all four original
measures. Suppose a smooth joined path from \(g_i\) ends at \(h_i\)
with a parametrized round product seam on this final boundary, up to
a smooth conformal multiplier. If its whole chiral energy stays
strictly below \(2T\), then for each chirality
\[
 E_s(g_i;B_i)-E_s(h_i;B_i)
   >(1-2^{-k})e_i+2^{-k}\lambda_i^s(Q_i)+I-2T-o(1).
\]
For \(k=1\) the limiting lower bound is positive in both signs when
\(I>4T/3\). At \(I=4T/3\) the sign selected on the original marked
core still has a fixed positive lower bound from
\(\lambda_i^s(Q_i)\).
One universal finite \(k\) gives a bound above \(\varepsilon_*\)
in the limit on every surviving nonstandard type.
\end{corollary}
\begin{proof}
Use tolerance \(t_i\to0\) at each of the fixed number of steps.
The new root is the preceding disk. Its old collar can also be
protected by taking an outer parallel disk with \(o(1)\) additional
raw mass. The errors satisfy the recursion \(d_j\le d_{j-1}/2+O(t_i)\).
Only the initial root uses the comparison curvature ceiling.

At the joined endpoint, stretching the actual round cylinder into
the disk and compactifying the retained end gives a positive metric
on the original smooth type as in Theorem~\ref{thm:simultaneous}.
Indeed the path and stretching stay in the positive component,
since stretching adds zero Weyl energy. The strict energy bound
for this source gives a positive lower Yamabe bound independent of
stretch length. Compactly supported tests on the infinite end
inherit it. The conformally flat cap adds zero chiral mass.
Consequently \(E_s(h_i;V_i)\ge I\), so \(E_s(h_i;B_i)<2T-I\).
More generally, if a supplied positive same-type cap has energy at
most \(E_s(h_i;V_i)+d_i\), and the actual division has absolute
error at most \(\tau_i\), the finite lower bound on the saving is
\((1-2^{-k})e_i+2^{-k}\lambda_i^s(Q_i)+I-2T-\tau_i-d_i\),
with strict inequality from the endpoint ceiling. On the matched
round cap \(d_i=0\). This gives the asserted estimate.
Take \(2^{-k}(T+\varepsilon_*)\le\varepsilon_*\) and use
\(I>T+\varepsilon_*\) from spectral recognition. Then
\((2-2^{-k})I-2T>\varepsilon_*\).

In particular a path supported up to conformal changes on a final
collar of raw full curvature \(o(1)\) can lower the original disk
energy by only \(o(1)\). Every such path to the stipulated round
seam must reach the \(2T\) wall at some parameter for all large
indices. The assertion concerns the enlarged disk with the specified
charges. Other widths can remove positive disk-side mass, and the
original small disk retains its independent repair question.
\end{proof}


\section{Simultaneous replacements on the source}

Write \(\mu_i=dV_{g_i}\), \(\lambda_i=|W^+_{g_i}|_{\rm HS}^2\mu_i\)
and \(\nu_i=|\Rm_{g_i}|_{\rm tens}^2\mu_i\). We use one subsequence
and its actual markings. Physical volume, chiral energy and total
curvature measure distinct features of this sequence. In particular,
the positive-Ricci distance used to select disks in the preceding
section does not supply a first-order estimate in the physical metric.
An estimate computed in a conformal comparison must state the metric
used for its total curvature and derivative norms.

\subsection{One exterior competitor from all marked collars}

\begin{theorem}[Simultaneous removal of marked standard disks]\label{thm:simultaneous}
Let \(g\) be a metric of positive Yamabe class on \(M\), with energy
\(e=E_+(g)<8\pi^2\). Take finitely many disjoint parametrized smooth
closed standard four-disks \(B_j\) with disjoint oriented bicollar
charts on \(\exp(-1)<|x|<\exp(1)\). In the given markings
\(\{|x|=1\}=\partial B_j\), and \(|x|>1\) points towards the retained
piece \(P=M\setminus\bigcup_j\operatorname{int}B_j\).
Write the pulled-back source metric as \(u_j^2m_j\) with \(u_j>0\)
smooth. Put
\[
 \begin{gathered}
 a_j=\|m_j-\Eucl\|_\infty,\qquad b_j=\|Dm_j\|_{L^4(dx)},\qquad
 K_j=\|\Rm_{m_j}\|_{L^2(m_j)},\\
 \eta^2=\sum_j[C(a_jK_j+a_j+b_j+b_j^2)]^2 ,
 \end{gathered}
\]
where \(a_j\) is within a fixed small ellipticity threshold and \(C\)
depends only on the fixed annulus and cutoff. Define
\begin{equation}\label{eq:discount}
 c=2\sqrt{e}\,\eta+\eta^2,\qquad m=\sum_j E_+(g;B_j).
\end{equation}
If \(e+c<8\pi^2\), there is a smooth metric \(k\) on the original
oriented \(M\) with
\begin{equation}\label{eq:disc-bound}
 E_+(k)\le e-m+c,\qquad
 Y([k])\ge\sqrt{48(8\pi^2-e-c)}>0 .
\end{equation}
In particular \(e-m+c<I\) contradicts the definition of \(I\).
On the chosen minimizing sequence, the theorem gives
\begin{equation}\label{eq:removed}
 m_i\le s_i+c_i
\end{equation}
at each index where its hypotheses hold, even if the finite number
of disks changes with \(i\).
\end{theorem}

\begin{proof}
Write \(t=\log|x|\). Choose a fixed smooth \(0\le\psi\le1\) which
equals one around \(t=0\) and vanishes for \(|t|\ge3/4\).
In chart \(j\) use
\[
 n_{j,s}=\Eucl+(1-s\psi)(m_j-\Eucl),\qquad 0\le s\le1.
\]
Restore the multiplier \(u_j^2\) and keep \(g\) elsewhere. This
makes one smooth path of metrics on the actual source. To estimate
it, express the all-lowered curvature as a constant-coefficient
linear function of \(D^2m\) plus terms quadratic in \(Dm\) with
smooth metric coefficients. In an oriented coordinate frame
orthonormalized by \(m^{-1/2}\), write \(F_\pm(m)\) for the Weyl
matrix times \((\det m)^{1/4}\). Its squared Euclidean \(L^2\) norm
is the local energy. The two-derivative principal part cancels
from \(\Rm_{\rm low}(n_{j,s})-(1-s\psi)\Rm_{\rm low}(m_j)\),
except for cutoff commutators of orders one and zero in \(m_j-\Eucl\).
The commutators have \(L^2\) cost \(C(a_j+b_j)\). The terms quadratic
in first derivatives cost \(C(b_j^2+a_j+b_j)\). Changing the
projection and volume density multiplies the old curvature by an
algebraic map of norm at most \(Ca_j\). Consequently
\[
 \|F_\pm(n_{j,s})-(1-s\psi)F_\pm(m_j)\|_2
 \le C(a_jK_j+a_j+b_j+b_j^2)
\]
uniformly in \(s\). This calculation is entirely in the collar
comparison gauge. Chiral energies after restoring \(u_j^2\)
are the same as their values in that gauge.

Take the Hilbert direct sum over the disjoint collars and their
unchanged complement. The positive-chiral leading term has norm
at most \(\sqrt e\), and the combined error has norm at most
\(\eta\). This requires both sides of every cut while the path is
joined. Thus the path energy is bounded by \(e+c\). Restricting
the same estimate to \(P\) at the endpoint gives the sharper
bound \((\sqrt{e-m}+\eta)^2\le e-m+c\). The incoming charge is the
energy of each entire raw disk in \(g\). No derivative or curvature
hypothesis inside the discarded disks away from the collars enters
this comparison.

At \(s=1\) the collar comparison is Euclidean near \(t=0\). In a
smaller collar choose the conformal multiplier \(|x|^{-1}/u_j\)
with logarithm smoothly cut off further out. The actual metric is
then \(dt^2+\gamma\) there, with \(\gamma\) the unit round metric.
The conformal change preserves all chiral energies. At this joined
endpoint the class is positive by Lemma~\ref{lem:wall}. Stretch
each exact product into the discarded side by replacing \(dt^2\)
there by \((1+L\kappa(t))^2dt^2\), with \(\kappa\) a nonnegative
bump supported strictly between the seam and the disk-side edge of
the exact product region, of integral one. Arclength increases the
product length by \(L\) and keeps \(P\) fixed. For every finite
\(L\ge0\) this is a smooth metric on \(M\), and the stretching path
has unchanged chiral energy. Hence all the stretched sources have
Yamabe constant at least \(\sqrt{48(8\pi^2-e-c)}\).

Append one half-infinite round product to each boundary of the
straightened \(P\). Any compactly supported test there extends
by zero to a joined source stretched past all its support. The
same lower Yamabe bound holds for these tests. On an end at
arclength \(s\), a positive conformal factor equal to one near
the seam and to \(\exp(-s)\) near infinity fills in one smooth
point in the radial coordinate \(r=\exp(-s)\). The entire added
piece is conformally flat, and the given angular collar
identifies it with the actual removed standard disk. The result
is a metric on the original oriented \(M\). A finite set of
points has zero \(H^1\) capacity in four dimensions: for smooth
tests a cutoff at scale \(r\) has squared gradient cost \(O(r^2)\)
per point, and density extends this to all tests. Conformal
covariance therefore gives the lower bound on the compact
metric. Its Weyl energy is precisely the energy of the retained
piece. This proves \eqref{eq:disc-bound} and then \eqref{eq:removed}.
\end{proof}

\paragraph{Cost measured on the actual collars.}
Let \(s\) and \(s_{\rm out}\) denote the original positive chiral
energy on all the collars and on their exterior halves, respectively.
Let \(r=e-m\) be the raw retained charge. Take the direct-sum
\(L^2\) splice estimate only on the collars, adding back the squares
on the unchanged regions. The joined path has energy at most
\[
 e-s+(\sqrt s+\eta)^2=e+2\sqrt{s}\,\eta+\eta^2 .
\]
If this bound is strictly below \(2T\), the same stretching and
compactification give a positive same-type metric of energy at most
\[
 r-s_{\rm out}+(\sqrt{s_{\rm out}}+\eta_{\rm out})^2
   =r+2\sqrt{s_{\rm out}}\,\eta_{\rm out}+\eta_{\rm out}^2 ,
\]
where \(\eta_{\rm out}\le\eta\) is any certified direct-sum allowance
on the exterior halves obtained by restriction of the collar estimates.
In particular choose the disks of Theorem~\ref{thm:endpoint-disks} on
a slow diagonal giving \(r_i\to0\). Smooth disjoint tubular collars
can be narrowed so that \(s_i\to0\). Whenever their normalized charts
meet the comparison hypotheses, \(\liminf\eta_i\ge\sqrt T\):
a subsequence bounded strictly below \(\sqrt T\) gives a joined
energy strictly below \(2T\) and a retained energy strictly below
\(T\). This is a bound on the certified allowances for these charts.
Narrowing a collar changes its derivative and ellipticity tests.

\subsection{A conformally invariant coordinate test on a selected collar}
\label{sec:disk-collar-comparison}

Take a marked bicollar of one of the actual disks in
Theorem~\ref{thm:endpoint-disks} or Theorem~\ref{thm:band-charge}, contained in its open model
ball. Use the disk's angular parametrization and send its
boundary to the unit sphere of \(A=\{e^{-1}<|x|<e\}\) by an
increasing radial change of parameter. For the resulting
embedding \(\Psi:A\hookrightarrow M\), put
\[
 m=(\det p)^{-1/4}p=\exp h,\qquad p=\Psi^*k_i,\qquad
 G_\alpha(m)=\partial_j(m^{\alpha\alpha}m^{j\alpha})
 \quad\text{(no sum on \(\alpha\))}.
\]
The same \(m\) arises from the physical pullback \(\Psi^*g_i\).
Before normalization the coordinate flux has the expression
\(\sqrt{\det p}\,p^{\alpha\alpha}p^{j\alpha}\), which is
unchanged by the conformal multiplier. Thus \(G=0\) is the
four-harmonic coordinate condition in either gauge. The
determinant-one and \(n\)-harmonic gauge is the one of
Liimatainen--Salo \cite[Theorem~1.1]{LiimatainenSalo2016ConformalGauge}.

\begin{proposition}[Physical Weyl bound and differentiated coordinate residual]
\label{prop:disk-collar-comparison}
Fix \(A_0\Subset A\) around the seam and a straightening cutoff
\(0\le\psi\le1\) supported in \(A_0\), equal to one near the
seam. Set \(a=\|h\|_{\infty,A}\), \(\tau=\|DG(m)\|_{2,A}\),
and \(w^2=4(E_+(g_i;\Psi(A))+E_-(g_i;\Psi(A)))\). Derivative
norms use the Euclidean components and measure on the indicated
annulus. For \(a\) within a fixed small threshold, write
\(b=w+\tau+a\). Then
\[
 \|D^2h\|_{2,A_0}\le Cb,\qquad
 \|Dh\|_{4,A_0}^2\le Cab .
\]
With the source's scalar multiplier restored, the path
\(n_s=\exp((1-s\psi)h)\) gives the joined and exterior
splice estimates of Theorem~\ref{thm:simultaneous} using the
per-collar error
\begin{equation}\label{eq:disk-collar-allowance}
 \beta=C(a+\sqrt{ab}+ab).
\end{equation}
The constants depend on the fixed trim and profiles.
\end{proposition}

\begin{proof}
The matrix \(h\) is trace-free. At the Euclidean metric let
\(\mathcal W_0\) denote linearized covariant Weyl and put
\(\mathcal G_{0,\alpha}h=-\partial_jh_{j\alpha}
-\partial_\alpha h_{\alpha\alpha}\), with no sum on \(\alpha\).
At \(\xi\ne0\), the Weyl symbol kernel consists of
\(\ell\Eucl+\xi\otimes u+u\otimes\xi\). Indeed, in a frame with
first direction parallel to \(\xi\), the Riemann symbol has only
the components generated by \(R'_{1a1b}\), \(a,b>1\). If its
Weyl part vanishes, it is the Kulkarni--Nomizu product of the
Euclidean metric with a Schouten tensor. Restriction to the
three-dimensional transverse space forces the transverse
Schouten block to vanish. The mixed components force its
\(1a\) entries to vanish, leaving \(R'_{1a1b}\) proportional to
\(\delta_{ab}\). The converse follows from conformal and
diffeomorphism invariance. For the trace-free kernel,
\(4\ell+2\xi\cdot u=0\), and
\[
 \xi_j h_{j\alpha}+\xi_\alpha h_{\alpha\alpha}
 =(|\xi|^2+2\xi_\alpha^2)u_\alpha .
\]
Thus \((\mathcal W_0,D\mathcal G_0)\) is jointly elliptic on
trace-free tensors. Plancherel bounds \(\|D^2f\|_2\) by the
\(L^2\) norms of this pair for compactly supported trace-free \(f\).

Use this bound on \(\chi^2h\), where \(\chi\) is a fixed
interior cutoff equal to one on \(A_0\), and set
\(X=\|\chi^2D^2h\|_2\) and \(J=\|\chi Dh\|_4^2\).
The nonlinear errors in covariant Weyl and \(DG\) are bounded
by \(C(a|D^2h|+|Dh|^2)\), by smooth dependence of their
principal coefficients on \(h\). The covariant Weyl norm in
coordinates is at most \(Cw\) by conformal invariance and
uniform ellipticity. The first-order cutoff terms carry one
factor of \(\chi\). It follows that
\[
 X\le C(b+\|\chi Dh\|_2+aX+J).
\]
Integration by parts in \(\int\chi^2|Dh|^2\) gives
\(\|\chi Dh\|_2\le C\sqrt{aX}+Ca\). Applying the same method
to \(\int\chi^4|Dh|^4\) gives
\(J^2\le Ca(JX+J^{3/2})\), hence \(J\le Ca(X+a)\).
Young's inequality and small fixed \(a\) give the two
interior bounds.

In the Euclidean density representation of chiral Weyl
used in Theorem~\ref{thm:simultaneous}, the matrix
\(F_\pm(\exp f)\) is a constant-coefficient linear function
of \(D^2f\) up to \(O(|f|\,|D^2f|+|Df|^2)\). In the
difference \(F_\pm(n_s)-(1-s\psi)F_\pm(m)\) the linear
second derivatives cancel except for cutoff terms of order
at most one in \(h\). The remaining terms are bounded in
\(L^2\) by \(C(a\|D^2h\|_2+\|Dh\|_4^2+\|Dh\|_4+a)\) on
the fixed support. This gives \eqref{eq:disk-collar-allowance}
on the whole support and its exterior half. Restoring the
scalar multiplier preserves the chiral energies. At \(s=1\)
the metric is Euclidean in comparison gauge near the seam,
so the same round-product insertion and the same original
disk marking give the stretched and capped comparisons.
\end{proof}

For disjoint selected collars, \(w_j^2<\zeta\) and
\(\sum_jw_j^2\le8e_i\). Put \(A_2=(\sum_j a_j^2)^{1/2}\)
and \(J_1=\sum_j a_j(w_j+\tau_j+a_j)\). Taking the direct
sum in \eqref{eq:disk-collar-allowance} gives
\(\eta\le C(A_2+\sqrt{J_1}+J_1)\). The physical contribution
\(\sum_j a_jw_j\) is bounded by \(A_2\sqrt{8e_i}\).
In particular \(A_2\to0\) and \(\sum_j a_j\tau_j\to0\)
are sufficient for the joined energy bound \(e_i+o(1)\)
and the retained bound \(e_i-\sum_j E_\pm(g_i;B_{ij})+o(1)\).
The first bound keeps the joined path strictly below \(2T\),
so the capped class is positive. For \(I<4T/3\) choose a tolerance smaller than
\((4T-3I)/2\) in Theorem~\ref{thm:band-charge}.
A fixed finite family then carries a uniformly positive amount
of either specified chirality. By passing to a subsequence some
fixed disk index has charge bounded away from zero for that sign.
A vanishing allowance on this charged disk alone gives a positive
same-type competitor with energy below \(I\) at large indices.
More generally it suffices that the measured savings exceed
\(e_i-I\) at one index and that the joined bound is strictly
below \(2T\). The global equality of the two chiral energies
allows the same argument using either sign.

The disk sandwich does not control these coordinate tests
for arbitrary parametrizations. For example, in a round
stereographic chart compose with the increasing radial map
\(r\mapsto r+N^{-3/2}\chi_1(r)\sin(N(r-1))\), with fixed
\(\chi_1=1\) near one and supported between both collar
faces. It fixes the seam and both faces, preserving every
source ball inclusion and curvature integral. If \(f_N\)
is this radial map and \(\rho_N=f_N'/(f_N/r)\), the
determinant-one metric has radial eigenvalue
\(\rho_N^{3/2}\) and tangential eigenvalue \(\rho_N^{-1/2}\).
The distortion tends to zero as \(N^{-1/2}\), but
\(\|Dh_N\|_{4,A_0}\ge cN^{1/2}\) from the radial derivative
of \(\log\rho_N\). Since Weyl vanishes, the interior estimate
implies \(a_N\tau_N\to\infty\). The ordinary round chart
straightens with zero cost. This example concerns the
prescribed coordinates. Finding a controlled annulus
retaining an actual charged disk remains an additional
geometric requirement.

\subsection{Capacity and measured repair of buffered disks}
\label{sec:collar-capacity}

A logarithmic buffer outside a protected core consumes auxiliary
volume. Fix one source of Theorem~\ref{thm:band-charge}, write
\(\mu_q=dV_q\), and suppose \(Y-6\ell/\sqrt{2}\ge\kappa>0\).
For a compact core \(C_j\) take nested exterior marked annuli
\(A_{j\ell}\), with disjoint interiors and smoothly extending faces.
Assume that the whole region through the outermost face lies in a
single comparison ball of radius at most one in the spectral-median
gauge. Let \(b_{j\ell}>0\) be each annulus's Euclidean logarithmic
width and \(\mathcal L_j=\sum_\ell b_{j\ell}\). In its marking the
determinant-one metric has smallest eigenvalue at least
\(\Lambda^{-1}\). Then
\begin{equation}\label{eq:marked-collar-capacity}
 \sum_j\mathcal L_j^3\mu_q(C_j)
 \le \frac{72\pi^2\Lambda^2}{\kappa^2}
          \sum_{j,\ell}\mu_q(A_{j\ell}).
\end{equation}
If all annuli are disjoint and avoid all the cores, the last sum
is at most \(1-\mu_q(\bigcup_j C_j)\), for pairwise disjoint cores.

Indeed, use one cutoff on each stack, equal to one through the inner
face and zero beyond the outer face. Across a layer let it drop by
\(b_{j\ell}/\mathcal L_j\) linearly in \(\log |x|\), keeping it constant
in the gaps. Its fourth-power Dirichlet integral is conformally
invariant. In determinant-one coordinates its radial derivative
has Euclidean size \((\mathcal L_j |x|)^{-1}\) on the layer, and
the metric costs at most \(\Lambda^2\) in the fourth power.
The summed integral is therefore at most
\(2\pi^2\Lambda^2/\mathcal L_j^3\). The cutoff's nonzero set is
in a comparison ball of at most half the total spectral mass.
The small-support Yamabe inequality in
Theorem~\ref{thm:band-charge} and H\"older give
\[
 \kappa\,\mu_q(C_j)^{1/2}
 \le 6\left(\int |d\chi_j|_q^4\,d\mu_q\right)^{1/2}
          \left(\sum_\ell\mu_q(A_{j\ell})\right)^{1/2}.
\]
Squaring and adding proves \eqref{eq:marked-collar-capacity}.
On a Euclidean annulus the log-radial cutoff attains the minimum
fourth-power energy for a unit drop. H\"older on each ray gives
the same lower bound \(2\pi^2/\mathcal L^3\) after integration over
the unit sphere.

In particular, suppose the protected cores are all the original E58
disks and cover auxiliary volume greater than \(1-d_0\). Disjoint
stacks each containing \(N\) annuli of width two would require
\[
 N^3(1-d_0)<\frac{9\pi^2\Lambda^2}{\kappa^2}\,d_0 .
\]
Thus even one outward annulus of fixed width on every original disk
is excluded at sufficiently small \(d_0\) and fixed \(\Lambda\),
when the cutoffs stay in the stated comparison balls. The estimate
also holds when separate markings supply the distinct layers.
A smaller protected core needs its own measured charge.

\paragraph{A refined allowance in the marked chart.}
On one fixed-width annulus of
Proposition~\ref{prop:disk-collar-comparison}, let
\(d=\|h\|_{2,A}\), retaining \(a=\|h\|_{\infty,A}\),
\(w^2=4(E_++E_-)(g;A)\) and \(\tau=\|DG(m)\|_{2,A}\).
Put \(b=w+\tau+d\). For the same fixed cutoff \(\psi\), now set
\[
 \beta=C(ab+\sqrt{db}+d),\qquad
 n_s=\exp((1-s\psi)h).
\]
For small universal \(a\), in the chiral density representation of
Theorem~\ref{thm:simultaneous} one has
\begin{equation}\label{eq:weighted-collar-allowance}
 \|F_\pm(n_s)-(1-s\psi)F_\pm(m)\|_{2,A}\le s\beta .
\end{equation}
The estimate also holds on either half. Let \(\chi\) be a fixed
interior cutoff equal to one on the support of \(\psi\), and write
\(X=\|\chi^2 D^2h\|_2\), \(Z=\|\chi Dh\|_2\) and
\(J=\|\chi Dh\|_4^2\). The joint elliptic estimate of
Proposition~\ref{prop:disk-collar-comparison} gives
\(X\le C(b+Z+aX+J)\). Integration by parts gives
\(Z^2\le C(dX+dZ)\) and
\(J^2\le C(aXJ+\|h\|_4 J^{3/2})\). For the last term one uses
H\"older on \(h\,D\chi\) and \(|\chi Dh|^3\).
Since \(\|h\|_4^2\le ad\), it follows that
\(Z\le C(\sqrt{dX}+d)\) and \(J\le Ca(X+d)\).
Absorption gives \(X\le Cb\).

Write \(F_\pm(\exp h)=P_\pm(h)D^2h+Q_\pm(h)[Dh,Dh]\) with
smooth coefficients, and set \(\vartheta=1-s\psi\).
The coefficient changes in the density difference are
\(P_\pm(\vartheta h)-P_\pm(h)=O(sa)\) and
\(\vartheta^2Q_\pm(\vartheta h)-\vartheta Q_\pm(h)=O(s)\).
Derivatives landing on \(\vartheta\) cost \(Cs(Z+d)\) in
\(L^2\). This gives \(Cs(aX+J+Z+d)\), proving
\eqref{eq:weighted-collar-allowance}. Restoring the source multiplier
preserves the chiral energies.

\paragraph{Signed path and exterior costs.}
For each selected disjoint patch, let \(B_j\) be the entire raw
disk, \(A_j\) the annulus and \(A_j^{\rm out}\) its exterior half.
For either chirality write
\[
 \begin{gathered}
 S_j^\pm=\lambda^\pm(A_j),\qquad
 S_{j,{\rm out}}^\pm=\lambda^\pm(A_j^{\rm out}),\\
 H_j^\pm=\int_{A_j}(1-\psi)^2d\lambda^\pm,\qquad
 H_{j,{\rm out}}^\pm=\int_{A_j^{\rm out}}(1-\psi)^2d\lambda^\pm .
 \end{gathered}
\]
Using \eqref{eq:weighted-collar-allowance}, put
\begin{align}
 U_j^\pm&=\max\{0,(\sqrt{H_j^\pm}+\beta_j)^2-S_j^\pm\},
       \label{eq:weighted-collar-path}\\
 V_j^\pm&=(\sqrt{H_{j,{\rm out}}^\pm}+\beta_{j,{\rm out}})^2
                                      -S_{j,{\rm out}}^\pm .
       \label{eq:weighted-collar-exterior}
\end{align}
Here \(\beta_{j,{\rm out}}\) may use the restricted error, or the
full \(\beta_j\). The \(L^2\) norm of the leading joined-path density
is the norm of an affine function of \(s\). Adding the allowance
\(s\beta_j\) gives a nonnegative convex function. Its square is
bounded by the larger endpoint square, which proves
\eqref{eq:weighted-collar-path}. The same triangle inequality
restricted to the exterior at \(s=1\) gives
\eqref{eq:weighted-collar-exterior}. In particular \(V_j^\pm\)
can be negative.

At each path parameter a single-patch change alone is a smooth
metric on the same closed oriented \(M\). Its two total chiral
increments agree by signature, and agree with the corresponding
increments on that patch. The same reasoning applies after the
single patch is capped with its marked conformally flat disk.
Thus if
\(e+\sum_j\min_\pm U_j^\pm<2T\), the simultaneous joined path,
product stretching and compactification of
Theorem~\ref{thm:simultaneous} yield a positive class on the
original \(M\). Its energy is at most
\begin{equation}\label{eq:weighted-collar-gain}
 e-\sum_j\max_\pm\{\lambda^\pm(B_j)-V_j^\pm\}.
\end{equation}
In particular, the sign used to estimate the gain may vary by
disk. An actual single-patch exterior change \(o_j^\pm\)
satisfies \(o_j^+-o_j^-=\lambda^+(B_j)-\lambda^-(B_j)\)
by the same signature identity.

\paragraph{Finite measured selection.}
For each protected core suppose there are \(N_j\) candidate
annuli of width two, with the preceding nesting and disjointness.
Normalize each to the annulus of
Proposition~\ref{prop:disk-collar-comparison}. Define
\(W_j^2=\sum_\ell w_{j\ell}^2\),
\(L_j^2=\sum_\ell\tau_{j\ell}^2\),
\(P_j^2=\sum_\ell d_{j\ell}^2\), and \(a_j=\max_\ell a_{j\ell}\).
Then \eqref{eq:weighted-collar-allowance} and Cauchy--Schwarz give
\[
 \sum_\ell(w_{j\ell}^2+\beta_{j\ell}^2)
 \le \mathcal D_j
 =W_j^2+C\bigl[a_j^2(W_j^2+L_j^2)
                +P_j(W_j+L_j)+P_j^2\bigr].
\]
The physical curvature ceiling on the containing ball controls
\(W_j^2\). In a common parent chart the \(L_j^2\) term is the
integral of \(|DG|^2\) on the union of candidate annuli, by
critical scaling, and \(P_j^2\) is bounded by a fixed multiple
of the integral of \(|h|^2|x|^{-4}\) there. Neither of these
coordinate estimates follows from the physical curvature ceiling.

Choose the candidate minimizing \(w_{j\ell}^2+\beta_{j\ell}^2\)
on each core. For each sign and on either region the surcharge
is at most \(2\sqrt S\,\beta+\beta^2\le S+2\beta^2\), with
\(S\le w^2/4\). The summed joined and exterior allowances
are therefore each at most \(\Gamma=2\sum_j\mathcal D_j/N_j\).
All chosen disks contain the corresponding \(C_j\).
If \(e+\Gamma<2T\) and
\(\sum_j\max_\pm\lambda^\pm(C_j)>\Gamma+(e-I)\),
then \eqref{eq:weighted-collar-gain} yields a strict same-type
comparison. A known charge bound on the original disks transfers
to their smaller cores after subtracting the actual lost charge.
The cost and the protected charge are independent measurements on
the same finite family.

\subsection{A first-order restriction on a connected repair}
\label{sec:connected-repair}

A connector between two charged disks could permit their opposite
chiral imbalances to cancel in one clamped comparison. Its role can
be tested on the source by a connection flux. Here a clamped
metric agrees with the original metric near and outside the
boundary of the changed region.

\begin{proposition}[Chiral transfer through a marked tube]
\label{prop:chiral-tube-transfer}
Let \(g,h\) be smooth metrics on the oriented homotopy four-sphere,
and let \(h\) be clamped on a connected patch \(D\) made of two lobes
\(P_1,P_2\) with a shaft \(V=(0,L)\times B_\rho^3\) between them, in
actual smooth coordinates. The shaft's lateral wall lies in
\(\partial D\) and the two transverse end faces meet the lobes.
Write \(\lambda^\pm\) for the original chiral measures of \(g\),
\(A=\nabla^h-\nabla^g\), and \(F_g,F_h\) for the two curvature
forms on \(TM\). Component norms on \(V\) use the product coordinate
frame and measure. Then the common closed chiral gain satisfies
\begin{equation}\label{eq:chiral-tube-transfer}
\begin{split}
 E_\pm(g)-E_\pm(h)\le {}&
 \sum_{j=1}^2\min_\pm\lambda^\pm(P_j)
     +(\lambda^++\lambda^-)(V) \\
 &+C(\rho/L)^{3/4}
  \left[\|A\|_{4,V}(\|F_g\|_{2,V}+\|F_h\|_{2,V})
                      +\|A\|_{4,V}^3\right].
\end{split}
\end{equation}
\end{proposition}

\begin{proof}
Use the matrix trace on the real tangent bundle. The signed chiral
density is
\(\Theta_g=-\tfrac12\operatorname{tr}(F_g\wedge F_g)\).
In an oriented orthonormal frame its coefficient relative to
\(dV_g\) is the trace of \(\mathcal R_g^{*}\circ *\circ\mathcal R_g\) on
two-forms. The squared norms of the two off-diagonal blocks of
\(\mathcal R_g\) agree and cancel. The scalar parts of the
diagonal blocks cancel too, leaving
\(|W_g^+|_{\rm HS}^2-|W_g^-|_{\rm HS}^2\).
For the interpolation of connections \(\nabla_t=\nabla^g+tA\),
the Bianchi identity gives
\(\partial_t\operatorname{tr}(F_t\wedge F_t)
 =2d\operatorname{tr}(A\wedge F_t)\).
Integrating and using the quadratic curvature formula yields
\[
 \Theta_h-\Theta_g=dQ,\qquad
 Q=-\tfrac12\operatorname{tr}\bigl(
     A\wedge(F_h+F_g)-\tfrac13 A\wedge A\wedge A\bigr).
\]
The relative transgression is on the common bundle, even though
the two metrics differ. See also the transgression calculation of
Chern and Simons \cite[Proposition~3.2]{ChernSimons1974Characteristic}.

For \(S_t=\{t\}\times B_\rho^3\), Stokes gives opposite increments
of the signed chiral measure on the two sides, given with their
boundary signs by \(\int_{S_t}Q\). The other boundary contributions
vanish by clamping. Integrate the absolute slice flux in \(t\).
H\"older applied to
\(|A|(|F_h|+|F_g|)+|A|^3\) on \(V\)
gives a factor \(|V|^{1/4}\). Dividing by \(L\) bounds the flux on
some slice by the last line of \eqref{eq:chiral-tube-transfer}.
If the two lobes' weaker signs match, estimate the common gain in
that sign on both sides of this slice, dropping nonnegative new
energies. If the weaker signs differ, change signs on one side
using its signed increment, which costs at most the absolute
slice flux. The original two shaft portions together cost at most
\((\lambda^++\lambda^-)(V)\). This proves the estimate.
\end{proof}

At a fixed source, two disjoint marked disks can be joined by
an embedded arc through their connected exterior. Choose standard
attaching three-balls in collars of the disks, take a product
normal tube, and round its feet relative to the marked cores.
The resulting boundary connected sum is a standard smooth ball.
The rounded extensions and the shaft's original curvature
integrals can be made arbitrarily small at that index. In
particular, if the disk of Corollary~\ref{cor:thin-raw-collar}
has nonzero chiral imbalance, the exterior has the opposite
imbalance. To find a second marked disk with almost as much
opposite imbalance, take finitely many disjoint small balls in
the open region of opposite density and enclose them using a
tree of such tubes. The original signed mass lost or added on
the tubes and roundings can be made arbitrarily small. The
second marking has its own metric and derivative costs.

The shaft can itself be flattened at small joined cost.
On any fixed smooth source use physical normal coordinates
along the arc with an orthonormal normal frame. The axis metric
in these coordinates is the product Euclidean metric, so
\(v=g-\Eucl\) satisfies \(\|v\|_\infty\le C_i\rho\) and
\(\|Dv\|_\infty+\|D^2v\|_\infty\le C_i\) on the shaft.
With a cutoff \(\chi\) equal to one on a smaller shaft and zero
near all faces, take \(|D^j\chi|\le C\rho^{-j}\) for \(j=1,2\).
On the path \(\Eucl+(1-s\chi)v\) the principal second-derivative
terms in the all-lowered curvature cancel after subtracting
\((1-s\chi)\Rm_{\rm low}(g)\). The remaining terms cost at
most \(C_i s/\rho\) pointwise. The change in projection of the
raw curvature has smaller order. The chiral density error on
the shaft thus has norm at most \(C_i s\sqrt{L\rho}\).
Using the raw chiral shaft energy \(O_i(L\rho^3)\) in the cross
term, the whole joined path has energy at most \(e_i+C_iL\rho\).
For small \(\rho\) this is below \(2T\), so
Lemma~\ref{lem:wall} preserves positivity. The connection difference
on this path is \(O_i(1)\) and its curvature form is \(O_i(1/\rho)\)
pointwise in the chart. Its slice flux is therefore \(O_i(\rho^2)\).
The constants here can depend on the source and arc. This gives a
flat interior shaft with the tips and cores unchanged.
In particular it does not produce a round collar on the
boundary of the joined disk.

A supplied connected chart permits a signed comparison using the
same clamped cutoff on the whole patch. For a source metric
\(p=\Eucl+v\) in that chart and a fixed small uniform bound on \(v\),
let \(\chi\) be one on the protected cores and compactly supported
in a region \(U\) of the chart. On \(U\) use the allowance
\[
 \begin{aligned}
 B=C\big(&\|\chi |v|\,|\Rm_{p,\rm low}|\|_2
       +\|\chi |Dv|^2\|_2+\||D\chi|\,|Dv|\|_2\\
       &+\||D^2\chi|\,|v|\|_2+\||D\chi|^2 |v|^2\|_2\big).
 \end{aligned}
\]
The principal curvature cancellation, with the derivatives of
\(\chi\) retained, bounds the \(L^2\) chiral density error by \(sB\)
along the path. Restore the common smooth conformal multiplier
if the chart used one. With
\(S^\pm=\lambda^\pm(U)\) and
\(H^\pm=\int_U(1-\chi)^2\,d\lambda^\pm\), this gives a joined
upper bound \(e+\max\{0,(\sqrt{H^\pm}+B)^2-S^\pm\}\) and an
endpoint upper bound \(e-S^\pm+(\sqrt{H^\pm}+B)^2\).
The joined bound follows by convexity of the Hilbert norm of the
leading path plus \(sB\). If the endpoint bound is below \(I\)
in either sign, the joined bound is \(e<2T\), and the path
yields a strict positive same-type comparison.
The cutoff terms include the shaft, both rounded tips and all
changed strips in the same chart. The incoming curvature
estimate of Corollary~\ref{cor:thin-raw-collar} gives neither
these width-dependent terms nor the candidate connection norms
in \eqref{eq:chiral-tube-transfer}. In particular a thin shaft
with bounded critical norms and small raw energy cannot transmit
a fixed gain from the opposite excesses of the two lobes.

\subsection{Variations on the charged interior}
\label{sec:core-variation}

Use the actual marked disk of Corollary~\ref{cor:thin-raw-collar},
or its marked enlargement in Corollary~\ref{cor:dyadic-repair}.
Take an open patch \(U\) in its interior. For a smooth
\(g\)-self-adjoint trace-free endomorphism \(A\) compactly supported
in \(U\), put
\[
 N_g(A;U)=\|A\|_{\infty,U}+\|\nabla A\|_{4,U}+\|\nabla^2A\|_{2,U},
 \qquad K=\|\Rm_g\|_{2,U}.
\]
All four quantities use the actual physical metric \(g\) and are
invariant under constant rescaling. The curvature norm is the full
tensor norm. The variation \(g_t(X,Y)=g(e^{tA}X,Y)\) preserves the
volume form and gives a smooth metric on the same oriented manifold
with the entire exterior fixed. By signature the two total chiral
increments agree, so the local increments on \(U\) also agree.

\begin{proposition}[Quantified localized variations]\label{prop:core-variation}
Let \(g\) have positive Yamabe class, \(e=E_+(g)<2T\), and
\(d=e-I\). Suppose \(N_g(A;U)\le1\).
There is a dimension-only constant \(C\) such that with
\[
 D=C(1+K)^2,\qquad
 t_* =\min\{1,(2T-e)/(2D)\},\qquad f(t)=E_+(g_t),
\]
both signs of the path through every \(0<t\le t_*\) remain
positive. On one of the two signs the energy saving is at least
\(\lvert f'(0)\rvert t-Dt^2/2\). On one of the two signs it is
at least \(-f''(0)t^2/2-Dt^4/24\). Consequently
\begin{equation}\label{eq:core-variation-bounds}
 |f'(0)|\le d/t+Dt/2,
 \qquad [-f''(0)]_+\le2d/t^2+Dt^2/12 .
\end{equation}
Either saving strictly exceeding \(d\) would give a positive
same-type competitor of energy below \(I\) at that source index.
\end{proposition}

\begin{proof}
Identify the oriented orthonormal bundles via \(e^{-tA/2}\) and use
the common volume form. The connection difference
\(S_t=\nabla^{g_t}-\nabla^g\) is obtained from the Koszul formula as
a smooth coefficient in \(t,A\) acting on \(\nabla A\).
Then \(R^{g_t}=R^g+d^\nabla S_t+S_t\wedge S_t\) as mixed curvature.
In the common frames its terms are smooth coefficients of
\(t,A\) times \(R^g,\nabla^2A,(\nabla A)^2\), respectively.
The frame correction and chiral Weyl projection preserve this form.
For \(|t|\le1\) and \(\|A\|_\infty\le1\), the coefficients and
their first four time derivatives are bounded. The Weyl blocks
and their first four time derivatives thus have \(L^2(U,g)\) norm
at most \(C(K+\|\nabla^2A\|_2+\|\nabla A\|_4^2)\).
Differentiating the integrated squared norm bounds
\(\lvert f^{(j)}(t)\rvert\le D\) for \(1\le j\le4\).
The path for \(|t|\le t_*\) stays below \(2T\) by the first
derivative bound and is positive by Lemma~\ref{lem:wall}.
The first saving follows from Taylor's formula in the descending
sign. The average of \(f(t)\) and \(f(-t)\) cancels the linear
and cubic terms, giving the second saving on at least one sign.
The admissibility of both endpoints yields
\eqref{eq:core-variation-bounds}.
\end{proof}

On the E61 disk \(K^2\) is bounded by the measured curvature in
the containing comparison ball. On an E65 enlargement one can use
the full physical curvature budget. The wall margin is uniform on
the late sources. Minimizing the right sides of
\eqref{eq:core-variation-bounds} in \(t\) gives bounds
\(C\sqrt{e_i-I}\) for the absolute first derivative and the negative
part of the second derivative on the critical unit ball.
In the Bach convention of \cite[\S1]{BahuaudGuentherIsenbergMazzeo2025Bach},
\(f'(0)=-\tfrac12\int_U\langle B_g,gA\rangle_g\,dV_g\), since
the tensor Weyl energy equals \(8E_+\) on the closed source.
The affine metric Hessian differs from \(f''(0)\) by the first
variation in direction \(gA^2\). Conformal invariance removes its
trace and the product rule bounds the critical norm of its
trace-free part by \(CN_g(A;U)^2\). The same lower bound, times
the squared norm of the direction, therefore holds for the
affine Hessian. These estimates concern variational norms on the
actual disk. A total Bach norm would require a further estimate.

In a measured cutoff layer \(C_0\) with collar parameter \(\rho\),
take a smooth fixed profile \(0\le\chi\le1\) of \(\rho/w\)
matching constants at the two edges. The product rule gives
\begin{align*}
 N_g(\chi A;U)\le N_g(A;U)+C_\chi\big[&
 w^{-1}\|d\rho\|_{4,C_0}
       (\|A\|_{\infty,C_0}+\|\nabla A\|_{4,C_0})\\
 &+\big(w^{-2}\|d\rho\|_{4,C_0}^2
           +w^{-1}\|\nabla^2\rho\|_{2,C_0}\big)
             \|A\|_{\infty,C_0}\big].
\end{align*}
For a smooth physical signed-distance layer the first norm is its
volume to the power \(1/4\), and the second is the squared-integrable
second fundamental form of the parallel faces. Apply the proposition
to the normalized cutoff direction, using its actual norm. Raw
collar curvature does not estimate these width terms. A conformal
perturbation alone preserves each chiral density on all compared
sets, and a compactly supported pullback preserves their integrated
energies. A successful interior comparison requires a signed
saving on the chosen supported variation.

\paragraph{A locally minimizing charged germ.}
The charge and raw curvature data alone allow a germ on which
clamped improvement requires crossing the wall. Take the unit
round product \(p\) on the product-oriented \(P=S^2\times S^2\).
Each chiral block has spectrum \((2/3,-1/3,-1/3)\), since the
Riemann operator is one on the two factor area forms, zero on
the mixed forms and \(R_p/12=1/3\). Thus \(E_\pm(p)=32\pi^2/3\).

Here is the positive-class comparison on \(P\).
For any smooth metric on it choose a nonzero harmonic self-dual
two-form \(\omega\), using \(b^+(P)=1\). Put \(v=|\omega|\) and
\(L_+=\lambda_{\max}(W^+)\). Where \(v>0\), write
\(\nabla\omega=\sum_{j=1}^3\alpha_j\otimes\omega_j\) in a unit
self-dual frame with \(\omega_1=\omega/v\). The equation
\(d\omega=0\) and the common length multiplier \(1/\sqrt2\)
of \(\alpha\mapsto*(\alpha\wedge\omega_j)\) give
\(|dv|^2=|\alpha_1|^2\le2(|\alpha_2|^2+|\alpha_3|^2)\).
With \(u_\epsilon=(v^2+\epsilon^2)^{1/4}\), the harmonic
two-form Weitzenb\"ock identity and the chain rule then give
\[
 (-6\Delta+R)u_\epsilon
   \le 6L_+u_\epsilon+
          R_+\frac{\epsilon^2 u_\epsilon}{v^2+\epsilon^2}.
\]
Multiplying by \(u_\epsilon\) and integrating gives an error
at most \(\epsilon\int_P R_+\,dV\). The Yamabe inequality and
\(L_+^2\le(2/3)|W^+|^2_{\rm HS}\) yield \(Y^2\le24E_+\)
when the class is positive. Gauss--Bonnet in a Yamabe
representative gives \(32\pi^2\le2E_++Y^2/24\le3E_+\),
using \(\chi(P)=4\) and \(\tau(P)=0\).
This is the product case of Gursky's harmonic-form inequality
\cite{Gursky1998WeylFunctional}.
The zero-Yamabe wall on \(P\) is \(16\pi^2\) by the same
Gauss--Bonnet identity in a scalar-flat representative.

Suppose a source is conformal to this product on an open
marked region containing the clamped support. Transplant its
smooth clamped path to \(P\), dividing by the fixed source
conformal multiplier in the marked region and using \(p\)
outside. Increments in each chiral integral are identical to
those on the source. A strictly lowering endpoint has negative
Yamabe sign on \(P\) by the positive-class estimate and the
zero wall. Its path reaches at least \(16\pi^2\) on \(P\),
so the source path reaches at least \(e+16\pi^2/3\).
In particular a joined path staying below \(2T\) cannot lower
the energy when \(e\ge8\pi^2/3\).

For an illustration in the full source band, paste a small
product normal ball into a round normal chart of \(S^4\).
The normal-metric mismatch is \(O(r^2)\) with derivative
bounds \(O(r)\) and \(O(1)\) through order two. Cutoff on
scale \(r\) costs \(O(r^4)\) in energy. In a disjoint
conformally Euclidean round chart a compactly supported warp
in one coordinate direction can raise the energy to any
given \(e\in[T,2T)\). For example, on an interior box change
\(dx_2^2\) to \((1+t\epsilon\cos(Nx_1))^2dx_2^2\) with a
cutoff outside it, \(0<\epsilon<1/2\). At \(t=1\) the inner
box is a surface times a flat plane and each block has
squared norm \((a''/a)^2/6\), \(a=1+\epsilon\cos(Nx_1)\).
The energy there diverges as \(N\to\infty\). Stop at the
first parameter giving the prescribed energy, so the source
class stays positive by Lemma~\ref{lem:wall}. After passing
to a smooth physical Yamabe representative, take a sufficiently
small disk strictly inside the conformal product germ. Both
complete sides are charged, full physical curvature is
arbitrarily small on a containing ball, and shrinking the
bicollar at that disk gives the relative raw curvature bound.
This example tests the local charge and width premises on a
standard source. It gives no realization on a hypothetical
nonstandard minimizing sequence.

\section{Energy bookkeeping and measured complements}

For arbitrary smooth replacement geometry, partition the source
into disjoint regions \(A_i\) (the marked disks), \(C_i\) (the
remaining changed connectors or splices) and \(X_i\) (the
exterior). Measure incoming energy on each of these regions in
\(g_i\). If \(k_i\) is identified with the original oriented \(M\),
define \(d_i(U)=E_+(k_i;U)-\lambda_i(U)\). A lower bound \(m_i\) on
\(\lambda_i(A_i)\), an upper bound \(b_i\) on \(E_+(k_i;A_i)\), and
signed upper bounds \(\Gamma_i,\Xi_i\) on \(d_i(C_i),d_i(X_i)\)
give
\[
 e_i-E_+(k_i)\ge m_i-b_i-\Gamma_i-\Xi_i .
\]
The exterior contributes zero when conformally retained and otherwise
may be conservatively charged its full new energy. Every changed
strip is scored once. Signature gives the equality of the sums of
the two signed chiral changes. In particular fixing the exterior
conformally fixes the chiral imbalance on the changed region.
If the cap uses a cutting construction, the oriented same-type
identification is an extra hypothesis until it is proved.

There is also an independent one-disk estimate. Suppose a
positive same-type replacement of just disk \(B_{ij}\) costs
at most \(\gamma_{ij}\ge0\) relative to the full raw exterior
energy. Then
\begin{equation}\label{eq:independent}
 \sum_j\lambda_i(B_{ij})\le N_i s_i+\sum_j\gamma_{ij}.
\end{equation}
Indeed, for each \(j\) the individual competitor has energy at
least \(I\) and at most \(e_i-\lambda_i(B_{ij})+\gamma_{ij}\).
This argument tests \(N_i\) separate competitors and does not
assert positivity of their simultaneous assembly.

If \(r_i^{\rm ext}=\lambda_i(M\setminus\bigcup_j B_{ij})\)
and \(N_i\ge1\), \eqref{eq:independent} equivalently gives
\begin{equation}\label{eq:residual-test}
 I\le r_i^{\rm ext}+\sum_j\gamma_{ij}+(N_i-1)s_i .
\end{equation}
For simultaneous removal \eqref{eq:removed} instead gives
\(I\le r_i^{\rm ext}+c_i\). A lower limit at most \(T\) on the
right of either estimate suffices by Theorem~\ref{thm:endpoint-smoothing},
including when the infimum bound comes without a stated minimizer.
At a fixed index an actual
independent competitor has energy at most \(T\) if
\(r_i^{\rm ext}+\sum_j\gamma_{ij}+(N_i-1)(e_i-T)\le T\), because
the sum of the \(N_i\) individual energy upper bounds is then
at most \(N_i T\).

Proposition~\ref{prop:joint-vol} gives finite disjoint disks on
the same source subsequence at each fixed tolerance \(\varepsilon\)
and curvature ceiling \(\zeta\). With its notation, set
\(G_\zeta=\{x\in\mathcal R:\nu(\{x\})<\zeta\}\), where \(\mathcal R\)
is the Euclidean regular set of the positive-Ricci location limit.
The bound on the complement of the disks is
\begin{equation}\label{eq:complement}
\begin{split}
 \limsup_i\lambda_i\!\left(M\setminus\bigcup_j B_{ij}\right)
 \le{}& [I-\lambda(X)]+\lambda(X\setminus\mathcal R)\\
 &+\lambda(\{x\in\mathcal R:\nu(\{x\})\ge\zeta\})+\varepsilon .
\end{split}
\end{equation}
The four quantities account for escape, nonregular points,
excluded atoms and selection loss. The source limit comes first
on each continuity window. Tightness, regular support and absence
of atoms of \(\lambda\) suffice to make this bound tend to zero.
The selection allows energy singular to \(\mu\) at eligible
points. In particular small aggregate repair cost on the selected
family would already rule out positive mass on \(G_\zeta\),
without any bound on the first three residuals. Either
\eqref{eq:removed} or \eqref{eq:independent} gives this implication
at fixed cardinality before refining the tolerances.

Alternatively suppose \(\lambda_i\) is uniformly absolutely
continuous with respect to physical volume in the sense that
\[
 \omega(a)=\limsup_{i\to\infty}\sup_{\mu_i(U)\le a}\lambda_i(U)
 \longrightarrow 0\quad\text{as }a\downarrow0 .
\]
A uniform physical \(L^{2+\alpha}\) bound on \(W^+\), \(\alpha>0\),
gives \(\omega(a)\le C a^{\alpha/(2+\alpha)}\) by H\"older.
Use the finite volume witnesses at fixed tolerance, with complement
of mass at most \(\varepsilon\), and assume independent repair
costs with \(\limsup_i\sum_j\gamma_{ij}\le b(\varepsilon)\to0\).
Equation~\eqref{eq:independent} then gives
\(I\le b(\varepsilon)+\omega(2\varepsilon)\), using the enlarged
volume tolerance eventually if the coverage is stated as a
lower limit. Thus \(I=0\), which identifies the standard type.
The same conclusion follows when the simultaneous error \(c_i\)
has the corresponding upper bound and the path is subcritical.
Tolerances and marked families are fixed before taking each
source limit. The critical bound \eqref{eq:budgets} alone gives
neither uniform absolute continuity nor control of the three
geometric residuals in \eqref{eq:complement}.

\section{Positivity of a prescribed exterior}

\subsection{Exact round ends}

\begin{theorem}[Round boundary form]\label{thm:round-form}
Let \(P\) be the complement of the interiors of finitely many
marked standard balls of \(M\). Suppose a smooth metric \(h\)
on \(P\) equals \(dt^2+\gamma\) on inner collars of its boundary,
where \(\gamma\) is the unit round metric and \(t\) increases out
of \(P\). Attach exact round half-cylinders using the given marks
and conformally compactify them to obtain \(\widehat h\) on \(M\).
Then \(E_+(\widehat h)=E_+(h;P)\). The Yamabe class is positive
if and only if the form
\begin{equation}\label{eq:round-form}
 \mathcal F_h(v)=
 \int_P(6|dv|_h^2+R_hv^2)\,dV_h+
 6\sum_b\big\langle \tau_b v,\sqrt{1-\Delta_\gamma}\,\tau_b v\big\rangle
\end{equation}
is strictly positive for every nonzero \(v\in H^1(P)\). The pairing
is the \(H^{1/2}\)--\(H^{-1/2}\) pairing on each unit round face.
\end{theorem}

\begin{proof}
On the joined complete manifold \(X\), choose the compactifying
factor \(f=1\) on \(P\) and near the seams and \(f=\exp(-t)\)
far along each end. Then \(\widehat h=f^2g_X\), and the radial
coordinate \(\exp(-t)\) gives the smooth marked filling and
conformal flatness of each end. The energy equality follows.
Under \(u=fw\), the quadratic form and \(L^4\) norm of the
compact test \(w\) are those of \(u\) on \(X\). To extend this
identity from smooth tests supported away from the added points,
add \(C\|w\|_{2,\widehat h}^2\) to the compact form for large \(C\).
Its transfer is \(Q_{g_X}(u)+C\int_X f^2u^2\,dV_{g_X}\).
The latter is equivalent to the ordinary unweighted \(H^1(X)\)
norm squared. Its scalar potential is bounded and positive on
\(P\) by the shift and at least six on the cylinders. Point
capacity on the compact side and compact-support approximation
on the cylinders give the corresponding \(H^1\) isomorphism.

The modes of \(-\Delta_\gamma\) have eigenvalues \(k(k+2)\),
for \(k=0,1,\ldots\), by restriction of homogeneous harmonic
polynomials on \(\mathbb R^4\). The scalar curvature on the
cylinder is six. For a mode with \(a=k+1\), the half-line form
is \(6\int_0^\infty(|u'|^2+a^2|u|^2)\,dt\). Completing the
square gives \(6a|u(0)|^2+6\int_0^\infty|u'+au|^2\,dt\).
The minimizing \(H^1\) mode is \(\exp(-at)u(0)\).
Summing in the eigenbasis gives the boundary form in
\eqref{eq:round-form}. The decaying extension is in \(H^1\)
for every \(H^{1/2}\) trace. For a fixed restriction \(v\) to
\(P\), gluing those extensions realizes the form value on \(X\).
Other extensions contribute a nonnegative surplus. A nonzero
global test restricting to zero on \(P\) has strictly positive
cylinder energy.

It follows that strict positivity in \eqref{eq:round-form}
is equivalent to strict positivity of the compact quadratic
form on nonzero tests. By attainment of the lowest compact
eigenvalue, this means the eigenvalue is positive. Its form
then bounds the \(H^1\) norm and hence the \(L^4\) denominator.
This proves the equivalence.
\end{proof}

With the outward normal of the cylinder at the seam equal to
\(-\partial_t\), its conormal boundary term is
\(6\sqrt{1-\Delta_\gamma}\). The corresponding zero equation
on \(P\) has boundary condition
\(6\partial_{\nu_P}v+6\sqrt{1-\Delta_\gamma}\,\tau_bv=0\).
The constant mode costs six times its squared trace norm.
The estimate \(\sqrt{1-\Delta_\gamma}\ge1\) gives a sufficient
scalar Robin condition by replacing the boundary term by
\(6\sum_b\|\tau_bv\|_2^2\). For connected \(P\) with nonempty
boundary this condition holds if \(R_h\ge0\). More generally,
if \(\lambda_{\rm Rob}>0\) is the first eigenvalue of the form
\(\int_P|dv|^2+\sum_b\|\tau_bv\|_2^2\), the sufficient bound is
\(\|(R_h)_-\|_\infty<6\lambda_{\rm Rob}\). Positivity of this
Robin eigenvalue follows from compactness and the fact that a
constant with zero trace vanishes. A negative test for
\(\mathcal F_h\) gives negative sign, and a nonnegative form
with a nonzero zero test gives zero sign.

Theorem~\ref{thm:round-form} applies before a positive closed
source has been found. For a supplied marked metric \(h\),
it tests entry below \(8\pi^2\), or direct recognition at
or below \(4\pi^2\), using the retained energy. Constructing
the metric on \(P\) and straightening its collars are separate
geometric steps whose Weyl cost must be included. Conversely
any fixed positive class has such presentations with energy
arbitrarily close to its own. Flatten sufficiently small
normal balls with cutoffs of comparable width and use the
endpoint estimate below to preserve positivity. On smaller
exact flat balls the inverse-radius conformal factor gives
the round product collars. Filling in the ends recovers the
flattened class because the radial compactifying factors
differ by smooth positive factors constant near each center.

\paragraph{A recognition test with no exterior energy ceiling.}
In the form of Theorem~\ref{thm:round-form}, replace \(R_h\) by
\(R_h-6L_h\), using the largest Weyl eigenvalue of either chirality
from Section~\ref{sec:endpoint-charge}. Positivity of this modified
boundary form on every nonzero test recognizes the capped domain.
The Weyl eigenvalue \(L_h\) has conformal weight minus two, so its
term preserves covariance of the Yamabe form. It vanishes on the
conformally flat caps. Minimizing the cylindrical
extensions gives the same boundary operator. The closed modified
operator has a positive lowest eigenvalue; combining the spectral gap
with boundedness of its potential gives an \(H^1\) lower bound.
Sobolev embedding yields a positive modified Yamabe constant, so
Chen and Zhu's Theorem 1.7 applies \cite{ChenZhu2014Conformal}.

For explicit exteriors take \([0,\ell+2]\times S^3\) with its product
markings. In the middle length \(\ell\) use \(dt^2+h_s\), where
\(h_s=e^{2s\phi}\gamma\) and the smooth function \(\phi\) has a
critical point with nonzero trace-free Hessian. On each unit end use
a fixed smooth cutoff of \(s\) equal to zero near the boundary and
equal to \(s\) near the middle. At the critical point the conformal
Ricci formula in dimension three gives nonzero trace-free Ricci for
\(s\ne0\). The product thus has nonzero Weyl curvature since
\(W_{0a0b}=-\tfrac12(\Ric_{h_s})_{0,ab}\) in an orthonormal frame,
up to curvature sign. Time reflection interchanges chirality, so
the positive chiral energy per unit middle length is strictly positive.
For fixed sufficiently small \(s\ne0\), uniform \(C^2\) comparison
with the round product gives \(R-6L\ge3\) on the whole exterior,
independently of \(\ell\). The modified form is positive while the
exterior energy grows linearly with \(\ell\). Capping the two
marked ends gives the standard smooth sphere in this construction.
A qualifying exterior on a further smooth type remains to be found.

\subsection{Endpoint comparison and assembly across general interfaces}

For a fixed source \(g\) of positive Yamabe sign and a candidate
metric \(k\) on the same marked manifold, put \(J=dV_k/dV_g\).
If \(Jk^{-1}\ge\alpha g^{-1}\) as quadratic forms on covectors,
where \(0<\alpha\le1\), define
\[
 B=\|(\alpha R_g-JR_k)_+\|_{2,g}.
\]
H\"older gives
\begin{equation}\label{eq:scalar-budget}
 Q_k(w)\ge\alpha Q_g(w)-B\|w\|_{4,g}^2,\qquad
 Y([k])\ge\frac{\alpha Y([g])-B}{\sqrt{\sup_M J}}
 \quad\text{if }B<\alpha Y([g]).
\end{equation}
This estimate uses one fixed background representative. In a
common conformal gauge it still applies with the scalar densities
of that gauge and the same \(Y([g])\). Where the source scalar
curvature is nonnegative and the metrics coincide, there is no
scalar loss. For the physical source the square of \(B\) is the
sum of squared losses over disjoint changed regions. If the
total changed set is \(D\), an upper bound is
\(\alpha Y_i\sqrt{\mu_i(D)}+\|J(R_k)_-\|_{L^2(D,g_i)}\).
This also shows why controlling just the negative part of the
new scalar curvature can miss the loss of old positive scalar
curvature.

For arbitrary smooth matched pieces \(P_a\), suppose each
conformal Laplacian has strictly positive Dirichlet first
eigenvalue. The Dirichlet form on \(H^1_0(P_a)\) is then coercive:
bounded scalar potential controls the gradient in terms of the
form and the \(L^2\) norm, and the eigenvalue bounds the latter.
Thus every \(H^{1/2}\) boundary trace \(\tau_a\) has a unique
zero-energy-equation extension \(H_a\tau_a\), defined weakly by
solving with zero Dirichlet data after choosing a trace extension.
Write \(\Lambda_a\) for its outward conormal form, including
coefficient six. Then
\[
 Q_{k;P_a}(H_a\tau_a+z_a)=
 \langle\tau_a,\Lambda_a\tau_a\rangle+Q_{k;P_a}(z_a),
 \qquad z_a\in H^1_0(P_a).
\]
The cross term vanishes by the weak equation. A piece with
several faces has one form on its entire trace vector. Global
traces are exactly the compatible collections under the actual
boundary identifications. The glued first eigenvalue is positive
if and only if all interior Dirichlet eigenvalues are positive
and \(\sum_a\langle\tau_a,\Lambda_a\tau_a\rangle>0\) for each
nonzero compatible collection. Necessity follows by extending
zero-trace tests by zero, and then by gluing the harmonic extensions.
Sufficiency follows by the decomposition and compact eigenvalue
attainment. Thus interior Dirichlet positivity alone leaves the
boundary form to be checked.

For example, starting from a physical source \(g_i\), suppose
each new interior remains Dirichlet positive and the decrease of
the total boundary form is at most \(\sum_j\beta_j\|\tau_j\|_2^2\),
where \(\beta_j\ge0\) and the norms use actual source faces.
For disjoint source collars \(V_j\) of widths \(t_j\), use trace
constants \(T_j\) with
\[
 \|\tau_j w\|_2^2\le T_j\left(t_j\|dw\|_{L^2(V_j)}^2+
                   t_j^{-1}\|w\|_{L^2(V_j)}^2\right).
\]
Such an estimate follows from the one-dimensional trace
inequality in product coordinates with metric comparison.
Disjointness for the gradients and H\"older for the \(L^2\)
terms bound the weighted trace loss by \(\rho Q_{g_i}(w)\), where
\[
 \rho=\frac{\max_j\beta_jT_jt_j}{6}
       +\frac{\big(\sum_j(\beta_jT_jt_j^{-1}\mu_i(V_j)^{1/2})^2\big)^{1/2}}{Y_i}.
\]
Infimize over all extensions of the given global trace. If
\(\rho<1\), the trace loss is strictly below the positive
assembled source form and the new assembly is positive. This
condition requires actual boundary-form bounds and trace
constants. Smooth topological disk marks alone give neither.

\section{A leading-order obstruction for small normal balls}

Fix a single smooth positive class metric \(g\) and a point \(p\).
In oriented normal coordinates rescaled to the unit ball, the
physical metric after constant rescaling is
\[
 h_r=\Eucl+r^2(q+b_r),\qquad
 q_{ij}(y)=-\tfrac13 R_{ikjl}(p)y^ky^l,\qquad
 \|b_r\|_{C^2}=O_g(r).
\]
Consider smooth boundary-fixed replacements \(h_r+r^2v\) with
\(v\) compactly supported in the unit ball and
\(\|v\|_\infty+\|Dv\|_4+\|D^2v\|_2\le K\) for fixed \(K\).
Use \(\mathcal W_\pm\) for the flat linearized chiral Weyl
operator. It is a constant-coefficient expression in second
derivatives, and \(\mathcal W_\pm q=A_\pm=W^\pm_g(p)\)
is constant. The curvature decomposition used in
Theorem~\ref{thm:simultaneous} gives, uniformly in this class,
\[
 F_\pm(h_r+r^2v)
 =r^2(A_\pm+\mathcal W_\pm v)+O_{L^2,g,K}(r^3).
\]
Indeed, after the linear term the bound is the product of the
metric error with the second derivative norm, plus the square
of the first derivative norm. Since
\(\int_{B_1}\mathcal W_\pm v\,dy=0\), the boundary-fixed global
energy difference is
\begin{equation}\label{eq:normal}
 E_\pm(k_r)-E_\pm(g)=r^4\|\mathcal W_\pm v\|_2^2+O_{g,K}(r^5),
 \qquad
 E_\pm(g;B_r(p))=r^4|B_1|\,|A_\pm|^2+O_g(r^5).
\end{equation}
For \(A_\pm\ne0\) there is no fixed-factor leading reduction in
this controlled class. A flat core of fixed relative radius
forces \(\mathcal W_\pm v=-A_\pm+O_{L^2,g,K}(r)\) there and has a strictly
positive leading cost for that block. Signature also gives the exact
identity
\[
 E_+(k_r;B_r)-E_-(k_r;B_r)
   =E_+(g;B_r)-E_-(g;B_r)
\]
for every smooth replacement fixing the entire exterior.
Equation~\eqref{eq:normal} concerns the specified scaled jets
on a fixed smooth metric. It gives no corresponding expansion
for arbitrary fillings on a degenerating sequence or on a
coupled family with changed intervening geometry.

\section{Whole marked limits and round restarts}

\subsection{Short spherical cuts and strict comparison}

The interpolation and stretching in Theorem~\ref{thm:simultaneous}
can be applied on a short marked spherical collar without assuming
either raw side is a standard ball. For each side, append an
infinite cylinder into the other side, compactify, and test in
the same joined stretched source. Thus both caps are positive
provided that the joined-path energy ceiling is strictly below
\(8\pi^2\). Both raw side energies use the entire physical
source sides. Each cap energy is at most its corresponding
raw charge plus the correction obtained by restricting the
collar error to that side.

Here it is necessary to identify the smooth type after cutting.
In a homotopy sphere every embedded \(S^3\) separates, since
otherwise a closed loop could have intersection one with that
sphere. For a finite disjoint family of cuts the incidence graph
of pieces is a tree. Capping each piece with balls, van Kampen and
Mayer--Vietoris identify \(\pi_1(M)\) with the free product of
the capped groups and \(H_2(M;\mathbb Z)\) with the sum of the
capped groups in degree two. Therefore every capped piece is a
simply connected integral homology four-sphere and hence a
homotopy sphere. If the sum of the capped energies is below \(2T\),
at most one has energy \(\ge T\), and the remaining positive
caps are standard by recognition. A standard capped leaf with
its parametrized cap deleted is a standard ball: ambient
isotopy carries the embedded standard cap to a coordinate ball
\cite[Theorem B]{Palais1960Extending}. The actual boundary map
extends over the standard ball, since each \(S^3\) diffeomorphism
is isotopic to an orthogonal map \cite[p.~553]{Hatcher1983Smale}.
Eliminating standard leaves in the tree shows that the capped
root has the same oriented smooth type as \(M\). This argument
uses marked standard cap disks in already recognized spheres.

In particular, along a minimizing sequence on a hypothetical
nonstandard \(M\), a marked short collar with errors \(o(1)\)
cannot have a fixed positive raw chiral charge on both sides.
The strict joined path and the fixed margin give positive caps
and their sum \(<2T\) at large indices. The possible nonstandard
cap has energy at most the source energy minus the opposite
raw charge plus \(o(1)\), below \(I\) if that charge is bounded
away from zero. This criterion also tests a physical
separated-quantile comparison when the \emph{whole} inner
charged ball and \emph{whole} outer charged set are on opposite
sides of the same actual collar. A first-order chart around
one point, or enclosing only parts of those sets, supplies no
such separation of the specified charges.

\subsection{Local minimization and smooth filled bases}

\begin{proposition}[Strong local chiral convergence]\label{prop:local-min}
Suppose, on an open oriented four-manifold \(U\), there are
oriented smooth embeddings of neighborhoods of a compact
exhaustion into the physical sources with pullbacks \(u_i^2m_i\).
Assume \(m_i\to h\) locally uniformly and strongly in \(W^{1,4}\),
where \(h\) is smooth, and \(\Rm_{m_i}\) is bounded in \(L^2\)
on each fixed compared compact in that gauge. Then the two chiral
matrix densities converge strongly in \(L^2_{\rm loc}\). The
limit locally minimizes chiral energy against small smooth
compactly supported metric perturbations, and is Bach-flat.
\end{proposition}

\begin{proof}
In a relatively compact coordinate domain write
\(r_i=F_\pm(m_i)-F_\pm(h)\). The constant-coefficient principal
curvature formula, bounded curvature and strong first derivatives
give weak \(L^2\) convergence \(r_i\rightharpoonup0\).
Transplant \(h+(1-s\psi)(m_i-h)\) to the source using \(u_i^2\),
where \(0\le\psi\le1\) is a smooth cutoff supported in the domain.
The product-rule terms in the first-order splice estimate tend
to zero in \(L^2\), as do the terms quadratic in the differences
of first derivatives. The change of the algebraic Weyl projection
uses uniform convergence and the curvature bound. The resulting
matrix density is
\(F_\pm(h)+(1-s\psi)r_i+o_{L^2}(1)\), uniformly in \(s\in[0,1]\).
The cross term with \(\psi F_\pm(h)\) vanishes in the limit by weak
convergence. Thus the change of either total energy is
\[
 -\int(2s\psi-s^2\psi^2)|r_i|^2\,dx+o(1).
\]
The path stays positive by Lemma~\ref{lem:wall}. At \(s=1\)
its energy is at least \(I\), giving strong convergence where
\(\psi=1\). Work on a finite cover for a fixed compact.

The repairs at \(s=1\) themselves approach the infimum and give
exact \(h\) in their comparison gauge on the smaller region.
A supported variation there changes the source energy by exactly
the corresponding local variation of \(h\). For small parameter
the whole path has a fixed subcritical margin, so the change
cannot be negative in the limit. The first variation of the full
Weyl functional is a nonzero multiple of the Bach pairing
\cite[\S2.1]{TianViaclovsky2004Moduli}. Here the total two chiral
variations agree by the closed signature identity on the source,
so this gives Bach-flatness.
\end{proof}

Suppose in Proposition~\ref{prop:local-min} the open manifold
is \(N\setminus P_0\), for a closed oriented homotopy four-sphere
\(N\) and finitely many ordinary points \(P_0\), with smooth
extension of \(h\) to \(N\). Write \(e_N=E_+(h;N)\).
Transfer fixed compactly supported Yamabe tests in the
comparison gauge. The scalar density is locally a divergence
linear in first derivatives plus a term quadratic in them,
with smooth coefficient dependence on the metric. Integration
against a fixed test square passes to the limit under the
first-order convergence. Zero point capacity then yields
\(Y([h])\ge\limsup_iY_i>0\).

At fixed small normal radius \(b\) the whole core of \(N\) with
those balls deleted has a marked embedding into each late
source. Its boundary components meet distinct raw exteriors:
in the Mayer--Vietoris sequence the map from the reduced
degree-zero group of the boundary to that of the complement
is an isomorphism, since the core and \(M\) are connected and
\(H_1(M)=0\). Every exterior touches a face, so there is one
per label. First install \(h\) on a thin neighborhood of each
face using the interpolation just proved. At fixed \(b\) the
change in both chiral integrals tends to zero by strong local
convergence. The smooth normal annuli of \(h\) are then
\(O(b^2)\)-close to Euclidean through second derivatives after
rescaling. Straighten them and cap all the pieces with
vanishing surcharge on a slow diagonal, taking the source
limit for each fixed \(b\) first. This order ensures the second
splice tests the curvature of \(h\) even if the source
curvature bound was only local in a variable conformal gauge.
The capped core has limiting energy \(e_N\), while the raw
exterior energies have sum tending to \(I-e_N\).

If \(e_N>0\) and a raw exterior retained positive limiting
charge, it would have positive charge opposite that of the
core across one cut. The preceding short-cut comparison
excludes this on a nonstandard minimizing type. Therefore all
exterior caps are eventually standard and their removal
identifies \(N\) with \(M\). The metric \(h\) attains \(I=e_N\).
If \(e_N=0\), signature implies both Weyl blocks vanish.
The conformal developing map identifies \(N\) with the round
conformal sphere \cite{Kuiper1949}. For \(I>0\), after passing
to limits of the finitely many exit charges, exactly one carries
\(I\). Two positive charges would again be opposite across
one cut. The round core cap and the other exterior caps have
vanishing energy and are eventually standard, so the surviving
cap is on \(M\). The same energy concentration holds for the
opposite chirality by signature on each cap and the small
two-chirality splices. On the standard type \(I=0\) all of these
limiting charges are zero.

\subsection{Signed restart estimates and exact equality}

For a surviving raw exterior \(V_i\), denote its capped metric
by \(c_i\), and let \(\sigma_i=E_+(c_i)-\lambda_i(V_i)\);
in particular this surcharge includes all changes across the
actual cutting interfaces. The signed comparison is
\begin{equation}\label{eq:signed-restart}
 e_i-E_+(c_i)=(e_i-\lambda_i(V_i))-\sigma_i .
\end{equation}
When both terms on the right are \(o(1)\), the restart preserves
the limiting energy. Strict descent needs their difference to
exceed \(s_i\). There is no sign inference from the separate
vanishing estimates.

For comparison of Yamabe constants, suppose the retained side
and the cap have a long exact round product comparison buffer of
length \(\ell\), in the respective conformal gauges, with the same
markings. Require the exterior cutoff support to compare with the
source, and the full cap-side support, including its splice,
to compare with the round cap model. Choose sine--cosine cutoffs
\(\chi_1,\chi_2\) with
\(\chi_1^2+\chi_2^2=1\), the derivative spread over the buffer.
One localized test extends to the source, the other to the
round model. The sum of their quadratic forms equals the cap
form plus \(6\int\sum|d\chi_a|^2w^2\). The squared cutoff gradients
have \(L^2\) norm \(O(\ell^{-3/2})\) on the product, and
\(\sum\|\chi_aw\|_4^2\ge\|w\|_4^2\). Consequently
\[
 Y([c_i])\ge \min(Y_i,S)-C\ell^{-3/2}
\]
in the exact model. This bound also holds with an additional
\(o_i(1;\ell)\) error on supplied first-order product comparisons
with the fixed-domain quadratic-form control first, at each
fixed \(\ell\). One first takes the source limit at fixed
compared windows and then increases \(\ell\) on a diagonal.
In the smoothly physically saturated round case \(Y_i\to S\),
and the upper bound \(S\) then forces the same cap limit.

In a wholly evacuated comparison, every fixed compared compact
has physical volume tending to zero. Use a cutoff for each
actual exit, equal to one on the deep exterior with gradient
supported on a fixed compared annulus. Its physical \(L^4\) gradient
bound is conformally invariant, so its squared physical \(L^2\)
gradient tends to zero by H\"older on the evacuated annulus.
For each limiting volume fraction \(v_j\), the source Yamabe
inequality on that cutoff yields \(\sqrt{v_j}\le v_j\).
The fractions sum to one, so exactly one exit has limiting
volume one. If its cutoff is on a retained annulus, transporting
the test to the cap gives \(\limsup_iY([c_i])\le\lim_iY_i\) for
that exit. Combined with the long retained-buffer lower bound
this gives equality of the limits there. If \(\lim_iY_i<S\), the
unit-volume exit must be the energy-bearing exit, because every
other cap has energy tending to zero and its own positive
Yamabe constant then tends to \(S\) by \eqref{eq:budgets}.
If \(\lim_iY_i=S\), the cap lower and upper bounds give equality
on the energy-bearing exit regardless of its physical volume.
These assertions use the stated buffer and evacuation hypotheses.

\begin{proposition}[Exact repeatable disk graft]\label{prop:graft}
Suppose \(h\) agrees with \(a(x)^2\Eucl\), \(a>0\) smooth, in an
oriented marked chart over a closed Euclidean ball of radius \(b\)
and its collar. Use a stereographic disk \(|z|\le b/\lambda\),
\(\lambda>0\), with gluing \(x=\lambda z\). One can choose a metric
on it that is round on the inner half, with a smooth conformal
transition. The pulled-back chiral measures and Yamabe constant
then agree exactly with those of \(h\). The graft can be
repeated inside the same marked germ.
\end{proposition}
\begin{proof}
The scaling on the disk with the exterior identity defines an
oriented diffeomorphism to the source. The round conformal factor
in \(x\) coordinates is \(2\lambda/(\lambda^2+|x|^2)\).
Multiply \(h\) by \(f_\lambda^2\), where
\[
 f_\lambda=\exp\!\left[\eta(x)
     \log\!\left(\frac{2\lambda}{a(x)(\lambda^2+|x|^2)}\right)\right],
\]
with a cutoff \(\eta\) one on the inner half and zero near the
boundary, and extend it as one on the exterior. Pull the metric
back by the disk diffeomorphism. This is the round metric near
the center, and smooth matching takes place in the
conformally Euclidean chart. Conformal covariance and invariance
give the asserted exact identities, and the center still admits
the specified round marking.
\end{proof}

Such an evacuated conformal comparison can arise at limiting
minimizing energy on any type admitting the input sequence.
For each fixed \(i\), first flatten one small normal ball of
\(g_i\), with errors tending to zero in energy and Yamabe constant
by choosing the radius sufficiently small at that index.
Indeed the scalar density changes by \(O_i(1)\) on volume
\(O_i(r^4)\), and the gradient and denominator densities by
\(O_i(r^2)\). The estimate \eqref{eq:scalar-budget} and its
upper counterpart give convergence of the Yamabe quotient
uniformly relative to the positive source bound. Choose a smooth
volume-one Yamabe minimizer \(g_i^\flat\) of the flattened class.
In its exact conformal disk write \(g_i^\flat=a_i(x)^2\Eucl\).
Choose \(A_i\to\infty\) and \(\lambda_i>0\) with
\(\lambda_i A_i<b_i/2\) and
\(\lambda_i A_i\sup_{B_{b_i}}a_i\to0\), where \(b_i\) is the
exact disk radius. Mark the comparison on \(|z|\le A_i\)
by \(x=\lambda_i z\). In these markings the physical pullback
is \(u_i(z)^2q_{\rm rd}\) with
\(u_i(z)=\lambda_i a_i(\lambda_i z)(1+|z|^2)/2\).
The auxiliary comparison metric \(q_{\rm rd}\) is the exact
round sphere on all fixed stereographic compacts. The compared
region has vanishing physical volume in \(g_i^\flat\).
It has zero chiral measure at every index, so its
single raw exterior carries all the charge and the exact graft
preserves the Yamabe constant even at finite indices. This
construction starts from the existence of an admissible
minimizing sequence and gives no new entry class on an arbitrary
type.

\section{The closed alternative}

\subsection{Variational tests at an attained minimum}

Suppose a smooth positive class attains \(I\), and choose its
volume-one Yamabe metric \(g\), with scalar curvature \(R>0\).
Positive Yamabe sign is open under smooth variations. Hence \(g\)
is Bach-flat, and the Hessian \(H_g=D^2E_+\) is nonnegative on
all smooth symmetric tensors. Conformal and diffeomorphism
tangents belong to its kernel: differentiate the invariance
of the first variation in these directions at the stationary
metric. A strict instability must be tested in the remaining
directions.

One exact necessary condition comes from the direction \(Z=Z_g\).
At each point use an orthonormal eigenframe for \(Z^\sharp\)
with eigenvalues \(\xi_i\). Set \(T_{ijk}=\nabla_iZ_{jk}\) and
\(C_{ijk}=T_{ijk}-T_{jik}\). In the Weyl convention of this
paper the Ricci-direction test is
\begin{equation}\label{eq:ricci-hess}
\begin{split}
0\le4H_g(Z,Z)=\int_M\Big[
 &\sum_{i,j,k}(\xi_j+\xi_k-\xi_i)T_{ijk}^2
  -\sum_{i,j,k}(\xi_i+\xi_j-\xi_k)C_{ijk}^2\\
 &+\frac R6\operatorname{tr}((Z^\sharp)^3)
   +\frac13|Z|^4\Big]\,dV_g .
\end{split}
\end{equation}
The sums are tensor contractions and remain defined at repeated
eigenvalues. Here are the identities behind the calculation.
Up to a topological constant \(4E_+=\int(|\Ric|^2-R^2/3)\,dV\).
For \((\Rm\cdot Z)_{ij}=R_{ikjl}Z^{kl}\), the gradient of the
latter integral is
\[
 G=-\Delta Z-2\Rm\cdot Z+\tfrac13(\nabla^2R)^0+
          \tfrac R6Z+\tfrac12|Z|^2g .
\]
This follows by subtracting one third of the scalar-square
variation from the Ricci-square variation
\cite[\S2]{GurskyViaclovsky2015}. At a minimum \(G=0\);
\(Z\) is trace-free and divergence-free by constant scalar
curvature. Along the variation with tangent \(Z\), substituting
the stationary identity into the connection variation gives
\(\dot{\Ric}=RZ/6+(Z^2)^0\),
\(\dot R=-|Z|^2\), and
\(\dot Z^\sharp=-RZ^\sharp/12\), with zero volume derivative.
Differentiate the integrated pairing
\[
 \int_M\langle G,Z\rangle\,dV=
 \int_M\left(|\nabla Z|^2-2\langle\Rm\cdot Z,Z\rangle+
                  R|Z|^2/6-|dR|^2/12\right)dV .
\]
To describe the contractions in this derivative, put
\(X=\langle\Rm\cdot Z,Z\rangle\) and \(Z^2=Z\circ Z^\sharp\).
In the eigenframe the gradient term has derivative
\[
 -\tfrac R6|\nabla Z|^2
 -\sum_{i,j,k}\xi_i T_{ijk}^2
 +2\sum_{i,j,k}(\xi_j-\xi_k) T_{jik}T_{ijk}.
\]
The integrated derivative of \(X\) is
\[
 \int_M\left[\sum_{i,j,k}\xi_i(T_{ijk}^2-T_{jik}T_{ijk})
 -\tfrac R6 X-\langle\Rm\cdot Z,Z^2\rangle\right]dV .
\]
In the latter computation the derivative of mixed curvature
is the antisymmetrized covariant derivative of the varied
connection, and the terms with \(\operatorname{div} Z\) vanish
after integration by parts. The multiple \(-R/6\) of the
stationary pairing cancels. Pair \(G=0\) with \(Z^2\) and
integrate to replace \(2\int\langle\Rm\cdot Z,Z^2\rangle\) by
\[
 \int_M\left[\sum_{i,j,k}(\xi_j+\xi_k)T_{ijk}^2+
             \tfrac R6\operatorname{tr}((Z^\sharp)^3)+
             \tfrac12|Z|^4\right]dV .
\]
The other scalar derivative contributes \(-|Z|^4/6\). Finally
\(2T_{ijk}T_{jik}=T_{ijk}^2+T_{jik}^2-C_{ijk}^2\);
swapping indices in the summed terms gives \eqref{eq:ricci-hess}.

This formula permits a genuine same-type competitor if its
right-hand side is strictly negative for a supplied stationary
metric. It has no general algebraic negative sign. For example,
at \(T_{ijk}=0\) and spectrum \((-3a,a,a,a)\) for \(a>0\), the
algebraic part is \(-4Ra^3+48a^4\), with either sign as \(a\)
varies. This is only an algebraic test. Similarly at a stationary
metric all directions in the conformal family \(\widehat g=f^{-2}g\),
\(f>0\), can be compared in the original Hessian. The conformal
Ricci law and gauge kernel give
\[
 H_{\widehat g}(\Ric_{0,\widehat g},\Ric_{0,\widehat g})
 =H_g(f^2Z-2df\otimes df,\,f^2Z-2df\otimes df).
\]
Indeed \(f^2\Ric_{0,\widehat g}\) represents \(f^2Z+
2f(\nabla^2f)^0\), and subtracting \(\nabla^2(f^2)\) removes a
diffeomorphism tangent up to trace. The formula gives further
necessary sign tests. Stability alone gives no negative test
on the remaining non-Einstein class.

\subsection{The reduced Ricci-flow test}\label{sec:cotton-variation}

At a smooth constant-scalar Bach-flat metric let \(g(t)\) start at
\(g\) under Ricci flow. Use the ordered-index tensor
\(C_{ijk}=\nabla_i Z_{jk}-\nabla_j Z_{ik}\), which is twice the
Cotton tensor formed from \((\Ric-Rg/6)/2\) with the same index order.
Write \(s_m=\operatorname{tr}((Z^\sharp)^m)\). The tangent to the
volume-normalized flow is \(-2Z\). The unnormalized tangent differs
by a pure trace, which is in the Hessian kernel at the critical
metric. The second derivative for either flow in its own time
variable is therefore \(4H_g(Z,Z)\). It reduces to
\begin{equation}\label{eq:ricci-cotton-hess}
 \left.\frac{d^2}{dt^2}E_+(g(t))\right|_{t=0}
 =\int_M\left[\frac7{12}s_2^2-s_4
      -3Z^{ia}C_{ijk}C_a{}^{jk}\right]dV_g .
\end{equation}
The quartic term is nonnegative. In a Ricci eigenframe with
\(Z^\sharp\) eigenvalues \(\xi_1,\ldots,\xi_4\), it is
\begin{equation}\label{eq:ricci-quartic}
 \frac7{12}s_2^2-s_4
 =\frac16\sum_{ij|kl}(\xi_i-\xi_j)^2(\xi_k-\xi_l)^2 ,
\end{equation}
where the sum runs over the three partitions into pairs. In fact,
if \(p=\xi_1+\xi_2\), \(q=\xi_1+\xi_3\) and \(r=\xi_1+\xi_4\),
then \(s_2=p^2+q^2+r^2\) and
\(s_4=3s_2^2/4-(p^4+q^4+r^4)/2\). Six times the quartic is
\((p^2-q^2)^2+(q^2-r^2)^2+(r^2-p^2)^2\). It vanishes exactly when
at least three eigenvalues coincide.

To derive \eqref{eq:ricci-cotton-hess}, the constant-scalar Bach
equation and \(\operatorname{div}Z=0\) give
\begin{equation}\label{eq:cotton-div-bach}
 \nabla^i C_{ijk}=-(W\cdot Z)_{jk} .
\end{equation}
Indeed, \(\nabla^i\nabla_j Z_{ik}=(\Ric\circ Z)_{jk}-(\Rm\cdot Z)_{jk}\),
while \(\Rm\cdot Z=W\cdot Z-Z^2-RZ/12+s_2g/2\), with \(Z^2=Z\circ Z^\sharp\).
Subtract the commuted derivative from the stationary formula for
\(\Delta Z\). The Weyl contraction \(\langle W\cdot Z,Z^2\rangle\)
vanishes pointwise in dimension four. To see this, set
\(w_{ij}=W_{ijij}\) in the eigenframe. Weyl trace-freeness gives
\(w_{ij}=w_{kl}\) for complementary pairs and
\(w_{12}+w_{13}+w_{14}=0\). The coefficient of each \(w_{ij}\) in
the contraction is
\[
 \xi_i\xi_j(\xi_i+\xi_j)+\xi_k\xi_l(\xi_k+\xi_l)
 =-\sum_{a<b<c}\xi_a\xi_b\xi_c .
\]
Integration by parts in \eqref{eq:cotton-div-bach} therefore gives
\(\int_M C_{ijk}\nabla_i(Z^2)_{jk}\,dV_g=0\).

For \(T_{ijk}=\nabla_iZ_{jk}\), let \(P_1,P_2,X_1,X_2\) denote
the sums of \(\xi_i T_{ijk}^2\), \(\xi_j T_{ijk}^2\),
\(\xi_i T_{ijk}T_{jik}\), and \(\xi_k T_{ijk}T_{jik}\), respectively.
Let \(J=Z^{ia}C_{ijk}C_a{}^{jk}=P_1+P_2-2X_1\).
The two derivative terms of \eqref{eq:ricci-hess} add to
\(5P_2-3J-2(X_1+X_2)\). The vanishing integrated pairing just
proved is \(\int_M(2P_2-X_1-X_2)\,dV_g=0\). Pairing the stationary
Bach equation with \(Z^2\) gives
\[
 2\int_M\Rm(Z,Z^2)\,dV_g
  =\int_M(2P_2+Rs_3/6+s_2^2/2)\,dV_g .
\]
The Weyl cancellation and the curvature decomposition give
\(\Rm(Z,Z^2)=-s_4-Rs_3/12+s_2^2/2\). Thus the integrated
derivative terms of \eqref{eq:ricci-hess} reduce to the integral
of \(-s_4-Rs_3/6+s_2^2/4-3J\). Its scalar terms then produce
\eqref{eq:ricci-cotton-hess}.

At an attained minimum, three times the integrated Cotton
contraction is at most the integrated quartic in
\eqref{eq:ricci-quartic}. A strict reverse inequality at a
stationary positive metric gives a same-type positive competitor
along short forward flow. The first derivative vanishes, and the
positive Yamabe sign persists by continuity of the first conformal
Laplacian eigenvalue. The contraction can have either sign at
the level of a first jet. For example, entries
\(T_{123}=T_{132}=T_{312}=T_{321}=h\) with all other entries zero
give \(J=(\xi_2-\xi_4)h^2\) and satisfy the trace and divergence
constraints on \(T\). They supply no control of the integrated
sign on an actual stationary metric in the surviving band.

\begin{proposition}[Cotton-flat positive constant-scalar gauge]
\label{prop:cotton-flat-bach}
Let \(g\) be a smooth Bach-flat Riemannian metric on a connected
four-manifold without boundary, of constant scalar curvature
\(R>0\). If \(C=0\), then \(g\) is Einstein or locally
conformally flat. On a closed smooth homotopy four-sphere the
locally conformally flat case is round up to scale and
diffeomorphism.
\end{proposition}

\begin{proof}
Equation~\eqref{eq:cotton-div-bach} gives \(W\cdot Z=0\).
Adding its diagonal components with indices 1 and 2 gives
\(p w_{12}=0\), using the same \(p,q,r\) and \(w_{ij}\) as above.
The other complementary pairs give \(q w_{13}=r w_{14}=0\).
If \(Z\ne0\) and its spectrum is different from \((a,a,-a,-a)\)
up to ordering, at most one of \(p,q,r\) is zero. The trace
identity for \(w_{ij}\) then makes every sectional Weyl coefficient
zero in every Ricci eigenframe.

The condition \(C=0\) says that the tangent-valued one-form
\(Z^\sharp\) is covariantly closed. Apply the covariant exterior
derivative again and use a Ricci eigenframe. With the curvature
indices ordered by \(R_{ijkl}=\langle R(e_i,e_j)e_l,e_k\rangle\),
the resulting identity is
\[
 \xi_kR_{ijkl}+\xi_iR_{jkil}+\xi_jR_{kijl}=0 .
\]
For \(l=i\), ordinary Bianchi gives
\((\xi_k-\xi_j)R_{ijki}=0\). Thus a Weyl coefficient with exactly
one shared index is zero when the other two eigenvalues differ,
since \(\Rm-W\) is diagonal in the Ricci eigenframe. When they
coincide, polarize the zero sectional Weyl coefficient using a
rotation in their common eigenspace. For entries with four
distinct indices, set \(x=R_{1234}\), \(y=R_{2314}\), and
\(z=R_{3124}\). Bianchi gives \(x+y+z=0\). The covariant closure
identity with last indices 4 and 3 gives, respectively,
\(\xi_3x+\xi_1y+\xi_2z=0\) and
\(\xi_4x+\xi_2y+\xi_1z=0\). Adding gives \(px=0\).
Permuting indices gives \(qz=ry=0\), so at most one of \(x,y,z\)
can be nonzero and Bianchi forces all three to vanish. Thus
\(W=0\) at every point where \(Z\ne0\) has spectrum different
from the double pairs.

The open set on which both \(W\) and \(Z\) are nonzero therefore
has Ricci defect eigenvalues \(a,-a\) of multiplicity two, with
\(a>0\) smooth. For either repeated eigenvalue, take two
orthogonal unit vectors in its eigenspace. Pair the Codazzi
equation with one of them to see that the eigenvalue derivative
along the other is zero. Applying this in the two eigenspaces
makes \(a\) locally constant. For inputs from opposite
eigenspaces the two sides of the Codazzi equation lie in
orthogonal summands and must both vanish. This gives parallel
eigenspace distributions and hence a local product of surfaces.
Its Gaussian curvatures are \(R/4+a\) and \(R/4-a\). If the first
two indices span the first factor, the decomposition of curvature
gives \(w_{12}=R/6\) and \(p=2a\), contradicting \(p w_{12}=0\).
Consequently wherever \(Z\ne0\), one has \(W=0\).

If \(Z\) is nonzero anywhere, \(W\) thus vanishes on an open set.
The unique continuation theorem for local conformal flatness of
smooth connected obstruction-flat metrics applies, since
obstruction-flatness is Bach-flatness in four dimensions
\cite[Theorem~3.8]{LassasLiimatainen2019ConformalHarmonic}.
It gives \(W=0\) globally. On a closed simply connected four-sphere
type, Kuiper's theorem supplies a conformal diffeomorphism to the
round sphere \cite{Kuiper1949}. Write the pulled-back round metric
as \(f^{-2}g\) with \(f>0\). The trace-free Ricci transformation
then gives \(f Z=-2(\nabla^2 f)^0\). Pairing with \(Z\) and
integrating using \(\operatorname{div}Z=0\) shows that
\(\int f|Z|^2dV_g=0\). Hence \(Z=0\) and the positive conformally
flat metric has constant sectional curvature.
\end{proof}

The Cotton-flat subclass on the homotopy sphere accordingly
has zero in \eqref{eq:ricci-cotton-hess} and enters the Einstein
comparison below. Cotton flatness remains a separate hypothesis
from Bach stationarity and second-order stability.


\subsection{Einstein and near-Einstein restrictions}

For a positive Einstein metric on the homotopy sphere, set
\(w=|W^\pm|_{\rm HS}\) for one block. A nonround Einstein metric
has nonzero blocks of both signs, by signature and the
positive Einstein space-form condition. Define
\(\phi=36\det W^\pm/w^2\) when \(w>0\), and \(\phi=0\) at its
zeros. The half-harmonic Weitzenb\"ock identity and refined
Kato inequality \cite[\S3]{GurskyMcKeownTyrrell2026} give
\[
 \tfrac12\Delta w^2=|\nabla W^\pm|^2+\tfrac R2w^2-18\det W^\pm,
 \qquad |\nabla W^\pm|^2\ge\tfrac53|dw|^2
 \quad(w>0).
\]
Apply the chain rule to \(u_\varepsilon=(w^2+\varepsilon^2)^{1/6}\).
After multiplication by \(u_\varepsilon\) and integration, the
gradient terms of the same sign are discarded by Kato. The
remaining potential is
\(36\det W^\pm/(w^2+\varepsilon^2)+R\varepsilon^2/(w^2+\varepsilon^2)\).
The determinant bound for a trace-free rank-three matrix,
\(|\det W^\pm|\le w^3/(3\sqrt6)\), bounds the signed term and
gives convergence by dominated convergence. The second term
times \(u_\varepsilon^2\) has integral \(O(\varepsilon^{2/3})\).
Thus \(u=w^{1/3}\) satisfies
\[
  Y([g])\|u\|_4^2\le Q_g(u)
  \le\int_M\phi u^2\,dV_g\le\|\phi_+\|_2\|u\|_4^2 .
\]
Equality would make \(u\) a positive Yamabe minimizer. The
nonround positive Einstein Obata theorem makes \(u\) constant
\cite{Obata1962Sphere,AkutagawaEndoSeshadri2019Gap}.
Equality in H\"older and the signed bound then gives
\(\phi=R\); the Weitzenb\"ock identity makes the nonzero block
parallel. Its simple parallel eigenline gives a nonzero parallel
two-form on the simply connected sphere, contrary to \(b_2=0\).
With \(\mathcal D_\pm=\int(24 w^2-\phi_+^2)\,dV_g\) for each sign,
the strict inequality and the Einstein form of \eqref{eq:budgets}
therefore give
\[
 e>\frac{16\pi^2}{3}+\frac{\max(\mathcal D_+,\mathcal D_-)}{72}
 \ge\frac{16\pi^2}{3}
   +\frac13\max_\pm\int_{\{\det W^\pm\le0\}}
                      |W^\pm|_{\rm HS}^2\,dV_g .
\]
The second inequality uses the full \(24w^2\) contribution on
the nonpositive determinant set.
If \(x\le y\le z\) are the eigenvalues of the block, set
\(\theta=(x-y)^2/z^2\in[0,9]\) on its nonzero set.
Then \(w^2=(3+\theta)z^2/2\) and
\(\phi_+=18 z(1-\theta)_+/(3+\theta)\).
For \(0\le\theta\le1\), the deficit is at least
\((84\theta-36\theta^2)z^2\). For \(1\le\theta\le9\) it
equals \(12(3+\theta)z^2\). In particular it is at least
\(16(x-y)^2\) throughout, giving the lower eigengap
correction \(2\max_\pm\int(x-y)^2/9\) to \(16\pi^2/3\).

There is also a uniform recognition test for almost Einstein
physical metrics. The additional compactness input used here is
the physical harmonic-radius estimate in the master Ricci-scale
argument. For smooth volume-one Yamabe metrics in the band
\(e\le8\pi^2-\delta\) on oriented homotopy spheres and \(p>2\),
a bound \(\|Z\|_p\le L\) gives a uniform harmonic-radius floor,
local \(W^{2,p}\) metric bounds and a whole marked closed
subsequence limit, weakly \(W^{2,p}\) and strongly
\(C^0\cap W^{1,2p}\) in its charts. The diameter bound in
that argument follows from physical noncollapse and unit volume.
Its harmonic-radius proof uses the exclusion of nonflat
Ricci-flat blow-up limits with the marked ends and band in
this topological setting \cite{Baker2026MarkedNecks}.

\begin{proposition}[A near-Einstein recognition test]\label{prop:ein}
Fix \(\delta>0\), \(p>2\) finite, \(L<\infty\), and \(0\le q<1\).
For one real Weyl block with ordered eigenvalues
\(\lambda_1\le\lambda_2\le\lambda_3\), set
\[
 D_q^\pm(g)=\int_M
   (\lambda_2(W^\pm)-q\lambda_3(W^\pm))_+^2\,dV_g .
\]
Under the preceding energy and \(L^p\) bounds, for every
\(\epsilon_0>0\) there is a threshold depending only on
\(\delta,p,L,q,\epsilon_0\) such that
\(\|Z\|_2^2+\min_\pm D_q^\pm(g)\) below that threshold
forces \(E_+(g)<\epsilon_0\). In particular one can take
\(\epsilon_0=T\) for smooth recognition. No minimizing
assumption across conformal classes is needed.
\end{proposition}
\begin{proof}
On a counterexample sequence the defects tend to zero. Pass
to one fixed chirality and a closed marked harmonic limit,
with \(Y_i\to y\ge\sqrt{48\delta}\). Interpolation gives
\(Z_{g_i}\to0\) in \(L^r\) for any \(2<r<p\).
The harmonic-coordinate metric equation passes to the positive
Einstein equation. Its interior regularity gives a smooth limit.
Subtract the limiting harmonic equation on smaller charts.
The quadratic first derivatives converge strongly, the Ricci
right-hand sides converge in \(L^r\), and the coefficient error
from the Hessians is controlled by the uniform convergence
and the \(W^{2,p}\) bounds. Interior elliptic estimates give
strong \(W^{2,r}\) convergence, hence strong \(L^2\) convergence
of the curvature and of the chiral densities on the closed
limit. Ordered eigenvalues are Lipschitz in a real symmetric
matrix, so the chosen defect vanishes on the limit.

The pointwise inequality on that block follows by continuity.
If the block is nonzero, Catino--Dameno's Einstein theorem
makes a cover of degree at most two conformally K\"ahler, with
the orientation reversed when testing \(W^-\)
\cite[Theorem 1.4]{CatinoDameno2026ConformalKahler}.
The cover is trivial on each connected component by simple
connectivity. The positive integral of the square of a K\"ahler
form contradicts \(b_2=0\). Hence the tested block and, by
signature, the other vanish. The limit is round, so the tested
sequence has energies tending to zero by strong convergence.
This contradicts \(E_+\ge\epsilon_0\).
\end{proof}

When the source is a minimizing sequence, bounded physical
\(L^p\) trace-free Ricci for one \(p>2\) likewise gives smooth
positive attainment on the original type. In that case the
strong chiral convergence and local stationarity argument of
Proposition~\ref{prop:local-min} starts from the rough harmonic
limit, uses strong \(W^{2,p}\) approximation by smooth local
comparisons, and passes the Bach identity to the limit.
The constant-scalar Bach equation for the limiting Ricci
tensor then bootstraps the harmonic metric to smoothness
\cite[Lemma 5.1 and Theorem 2.2]{AvalosCogoRoyoAbrego2026Bach}.
This attainment conclusion uses the bounded Ricci subsequence.
A sequence selected only by its chiral energies need not meet
that hypothesis or a prescribed Bach-error bound.

\section{Ricci growth at the recognition endpoint}\label{sec:endpoint-windows}

Let \(I=T=4\pi^2\) and retain the fixed physical minimizing
sequence of this paper. Smooth recognition implies that every
\(d_i=e_i-T\) is positive, since a smooth positive class of energy
at most \(T\) would give \(I=0\). The strict defect inequality
\eqref{eq:endpoint-deficit}, or independently the full-tensor
pinching classification of Fu \cite[Theorem 1.3]{Fu2018FourManifolds},
together with \eqref{eq:budgets} gives, with \(A_i=\|Z_{g_i}\|_{L^2(g_i)}\),
\begin{equation}\label{eq:endpoint-ricci-defect}
 A_i^2\le 8d_i,\qquad
 192\pi^2-48d_i\le Y_i^2\le192\pi^2+48d_i .
\end{equation}
In particular the physical trace-free Ricci energy tends to zero.

For a physical region \(U\), let
\[
 N_{i,q}(U)=\sup_{t>0}t\,\operatorname{Vol}_{g_i}
      (U\cap\{|Z_{g_i}|_{g_i}>t\})^{1/q},\qquad
 \gamma_q=2-4/q .
\]
This is the weak \(L^q\) quasinorm in the physical gauge.

\begin{proposition}[Ricci growth and curvature at the endpoint]
\label{prop:endpoint-ricci-windows}
For every fixed finite \(q>2\) and \(\kappa>0\), a physical
minimizing sequence at \(I=T\) satisfies
\[
 A_i^{2\kappa}N_{i,q}(M)\longrightarrow\infty,\qquad
 d_i^\kappa N_{i,q}(M)\longrightarrow\infty.
\]
In particular \(d_i^\kappa\sup_M|\Rm_{g_i}|_{\rm tens}\to\infty\).
Let \(a_i\to0\), \(x_i\in M\), and
\(U_i(D)=B_{g_i}(x_i,Da_i)\). If for every fixed finite \(D>0\)
there is a bound, eventually on one subsequence, on
\begin{equation}\label{eq:endpoint-protected-weak}
 \|Z_{g_i}\|_{L^2(U_i(D),g_i)}^{2\kappa}
       a_i^{\gamma_q}N_{i,q}(U_i(D)),
\end{equation}
then for each fixed \(A>0\), the full curvature integral on
\(B_{g_i}(x_i,Aa_i)\) tends to zero on that subsequence. One
sufficient condition for \eqref{eq:endpoint-protected-weak} is a
bound with \(d_i^\kappa\) in place of the factor involving the
local \(L^2\) norm.

For \(a_i=(d_i^\kappa N_{i,q}(M))^{-1/\gamma_q}\) the curvature
conclusion holds uniformly in the center. If \(r_i(\vartheta)\)
is the least physical radius at which some ball carries positive
chiral mass \(\vartheta e_i\), then
\[
 r_i(\vartheta)(d_i^\kappa N_{i,q}(M))^{1/\gamma_q}
     \longrightarrow\infty\qquad (0<\vartheta\le1/2).
\]
\end{proposition}

\begin{proof}
On any measure space put \(A=\|f\|_2\), \(B=\|f\|_{q,\infty}\).
The distribution function of \(|f|\) is bounded by
\(\min(A^2t^{-2},B^qt^{-q})\). For \(2<p<q\), split its
layer-cake integral at \(t=(B^q/A^2)^{1/(q-2)}\), obtaining
\[
 \|f\|_p\le C(p,q)A^\alpha B^{1-\alpha},\qquad
 \frac1p=\frac{\alpha}{2}+\frac{1-\alpha}{q} .
\]
If either norm vanishes both sides are zero. Choose
\(p=2+(q-2)/(1+q\kappa)\), so that
\(\alpha/(1-\alpha)=2\kappa\) and
\(\gamma_p=(1-\alpha)\gamma_q\). Therefore
\[
 a_i^{\gamma_p}\|Z_{g_i}\|_{L^p(U,g_i)}
 \le C(p,q)\left(\|Z_{g_i}\|_{L^2(U,g_i)}^{2\kappa}
          a_i^{\gamma_q}N_{i,q}(U)\right)^{1-\alpha} .
\]
On \(M\), a bounded subsequence of \(A_i^{2\kappa}N_{i,q}(M)\)
would imply a bounded subsequence of the global \(L^p\) norm.
The physical Ricci-scale theorem
gives smooth positive attainment on the original type, which
contradicts recognition at \(T\). The bound
\eqref{eq:endpoint-ricci-defect} gives the stated divergence in
terms of \(d_i\). Unit volume and the pointwise
Ricci contraction bound give \(N_{i,q}(M)\le C\sup_M|\Rm|_{\rm tens}\).

On the shrinking balls, the displayed interpolation gives the
strong \(L^p\) bounds on all protected windows required by the
physical Ricci-scale theorem in the companion master manuscript.
Its conclusion is vanishing full curvature mass on every smaller
fixed multiple of the scale. For the specified \(a_i\), the
sufficient bound involving \(d_i^\kappa\) holds uniformly over
the centers for every \(D\). Failure of uniform vanishing would
select a subsequence and centers with a positive integral. If a
bounded multiple of \(a_i\) reached a chiral quantile, its full
curvature integral would be at least \(4\vartheta e_i\), contrary
to this vanishing.
\end{proof}

This condition allows growth of the weak \(L^q\) norm controlled
by any fixed inverse power of the minimizing excess. The
critical estimate \eqref{eq:endpoint-ricci-defect} itself gives
neither an exponent above two nor the coordinate regularity
needed by the harmonic compactness argument. A candidate endpoint
must meet that critical estimate while its supercritical norms
grow faster than every such fixed power.

\section{The spectral deficit of exact critical classes}
\label{sec:exact-spectral-gap}

Put \(\sigma([g])=\int L_g^2dV_g\) and
\(\Delta_{\mathrm{sp}}([g])=36\sigma([g])-Y([g])^2\), using the
larger of the two top chiral Weyl eigenvalues as in
Section~\ref{sec:isotropic-gap}. For fixed \(\delta>0\) consider
the smooth positive Yamabe classes with exactly vanishing Bach
tensor on nonstandard oriented smooth homotopy four-spheres
and with \(e=E_+\le 2T-\delta\). Take
\(\kappa_\delta>0\) from the marked same-type ALE descent theorem
\cite[Same-type descent and Exact compactness]{Baker2026MarkedNecks}.

\begin{theorem}[Exact classes in a fixed band]\label{thm:exact-spectral-gap}
The indicated classes occur on finitely many oriented smooth
types. If the family is nonempty, \(\Delta_{\mathrm{sp}}\) has
an attained, strictly positive minimum \(m_\delta\) on it.
For any fixed \(0\le b<6\kappa_\delta\), the associated volume-one
Yamabe metrics with \(\Delta_{\mathrm{sp}}\le m_\delta+b\) are
sequentially compact modulo oriented diffeomorphisms. The same
attainment and compactness statements hold on each represented
type with its own spectral minimum. There is also a universal
\(\gamma_\delta>0\) such that any class in the band with
\(e-I^+(M)<\gamma_\delta\) attains \(I^+(M)\).
\end{theorem}

\begin{proof}
Choose the volume-one ordinary Yamabe minimizers of the source
classes. Bach-flatness, \(E_+\), and \(\sigma\) are conformally
invariant. The closed gap and Gauss--Bonnet give
\(T<e_i\le2T-\delta\) and \(Y_i\ge\sqrt{48\delta}\). Passing
to a subsequence write \(e_i\to e\), \(Y_i\to y\), and
\(\sigma_i\to s\). The exact compactness theorem allows the
source manifolds \(M_i\) to vary
\cite[Exact compactness]{Baker2026MarkedNecks}.
In its closed alternative the whole oriented diffeomorphisms
identify all late sources of the subsequence with the limit.
The limit has unit volume and scalar curvature \(y\), and has
class Yamabe constant \(y\) by passage of the Yamabe inequality
and the constant test. Smooth convergence of the Weyl blocks
gives convergence of \(\sigma\), so the limit retains \(e\) and
\(\Delta_{\mathrm{sp}}\).

In the ALE alternative the same-type theorem supplies a smooth
one-point conformal completion \((\widehat X,\bar h)\), positive
and exactly Bach-flat, with \(Y([\bar h])\ge y\) and
\(T<e_X=E_+(\bar h)\le e-\kappa_\delta\). It identifies the
completion with each late source of the subsequence by oriented
diffeomorphisms. This identification uses the whole marked core,
the buffer in its smooth puncture chart, and the cap of the actual
exterior with positive Yamabe constant and energy below \(T\)
\cite[Finite ordinary punctures and Same-type descent]{Baker2026MarkedNecks}.

At each fixed compact truncation of \(X\), smooth convergence
passes both \(E_+\) and \(\sigma\) to the limit on the truncation.
The conformal factor of the completion preserves their measures,
and its smooth added point has measure zero. On the entire
complement of that truncation in the actual source, integrate
\(|W^+|_{\rm HS}^2\le6L^2\) and subtract the compared core
integrals from the respective totals. Send the source index to
infinity first and then exhaust \(X\). This gives
\(e-e_X\le6(s-\sigma([\bar h]))\). Consequently
\begin{equation}\label{eq:exact-spectral-drop}
 \Delta_{\mathrm{sp}}([\bar h])
 \le \lim_i\Delta_{\mathrm{sp}}([g_i])-6(e-e_X)
 \le \lim_i\Delta_{\mathrm{sp}}([g_i])-6\kappa_\delta.
\end{equation}

Finiteness also follows directly from Chang--Qing--Yang
\cite[Theorem B]{ChangQingYang2006Gap}: \eqref{eq:schouten} gives
\(\int\sigma_2(g^{-1}A_g)\,dV_g\ge\delta/2\), and the Weyl integral
is uniformly bounded. The following argument also tracks the limiting
spectral deficit. An infinite list of pairwise distinct oriented types would give
a sequence with one class per type. Both branches contradict
this through the whole-source diffeomorphisms. A minimizing
sequence for \(\Delta_{\mathrm{sp}}\) in the band must take the
closed branch, since the completion in the ALE branch belongs
to the same family and \eqref{eq:exact-spectral-drop} gives a
strict decrease. The same reasoning excludes that branch for
\(\Delta_{\mathrm{sp}}\le m_\delta+b\) with \(b<6\kappa_\delta\).
This also works on one represented type because the completion
preserves that type. At an attaining positive class, Chen--Zhu's
strict and equality classifications give
\(\Delta_{\mathrm{sp}}>0\) on a nonstandard type
\cite[Theorems~1.4 and 1.6]{ChenZhu2014Conformal}.

For the near-infimum assertion let \(a_M\) be the first exact
positive chiral value below \(2T\) on a represented type.
The first-value theorem gives its attainment and an
\(\epsilon_M>0\) excluding distinct exact values in
\([a_M,a_M+\epsilon_M)\)
\cite[First fixed-type critical value]{Baker2026MarkedNecks}.
If \(a_M>I^+(M)\), take half their difference as the typewise
threshold. No exact class meets this condition on that type.
If \(a_M=I^+(M)\), take \(\epsilon_M/2\). A class within that
distance has \(e=I^+(M)\). The minimum of the finitely many
typewise thresholds is positive. When the family is empty
the assertions are vacuous.
\end{proof}

Under the whole smooth filled-base hypotheses of the descent
comparison, a minimizing sequence on a nonstandard type with
\(\Delta_{\mathrm{sp}}\to0\) can only have a conformally round
filled base. Indeed, strong local convergence of the chiral
matrix densities in Theorem~\ref{prop:local-min} gives
convergence of the top-eigenvalue maximum density. The largest
eigenvalue is Lipschitz in a self-adjoint matrix with its
Hilbert--Schmidt norm. The same compact exhaustion and Yamabe
test transfer give \(\sigma_N\le s\) and \(Y_N\ge y\), so
\(\Delta_{\mathrm{sp}}(N)\le0\). Chen--Zhu recognizes the filled
base. A positive-energy filled base under the stated puncture
markings would have the original smooth type by the marked
minimizing comparison. Thus the filled base is round and the
entire limiting charge lies on one actual exterior of the
marked cuts.

For a round restart on that exterior, suppose the changed
regions have vanishing discarded and inserted Weyl energy in
both signs, with the kept part conformally unchanged. The bound
\(L^2\le\tfrac23(|W^+|_{\rm HS}^2+|W^-|_{\rm HS}^2)\) gives
\(\sigma([c_i])=\sigma_i+o(1)\) for the resulting capped classes.
If the marked Yamabe comparison gives \(\liminf Y([c_i])\ge y\),
then \(\Delta_{\mathrm{sp}}([c_i])\to0\) by Chen--Zhu on the same
nonstandard type. This sequence cannot contain exact Bach-flat
classes at arbitrarily late indices within a fixed strict band,
by Theorem~\ref{thm:exact-spectral-gap}. General round restarts
still require their signed gain test. The exact threshold alone
does not select a stationary class from a minimizing sequence.

% E68d, 26 November 2026 UTC. Drafted for Edward B. Baker III with
% OpenAI Codex assistance. Actual serving model and effort unavailable.
% Input descent generation d9a6557377462fd3, assessed in the E63--E67
% receipt. Exact-band and supplied marked round alternatives only.

\subsection{Chiral overlap on an exact comparison disk}
\label{sec:exact-overlap}

For the two top Weyl eigenvalues \(z_\pm\), define the conformally
invariant measure \(d\mathcal O=\min(z_+^2,z_-^2)dV\).
Fix a strict energy band \(e\le 2T-\delta\), an integer \(k\ge2\)
and \(\varepsilon>0\). There is \(a>0\), depending on these three
parameters, such that every volume-one Yamabe representative \(g\)
in the nonstandard exact band of Theorem~\ref{thm:exact-spectral-gap}
has an oriented embedding \(F\) of a neighborhood of the Euclidean
closed ball \(\overline B_2\) and a scale \(r>0\) with
\[
 \|r^{-2}F^*g-\delta_{\mathrm{Euc}}\|_{C^k(\overline B_2)}
       <\varepsilon,\qquad
 \mathcal O(F(B_1))\ge a .
\]
For each fixed sublevel \(\Delta_{\mathrm{sp}}\le m_\delta+b\)
with \(0\le b<6\kappa_\delta\), one can also require \(r\) to
have a uniform positive lower bound. The disk in this assertion is
selected by the exact comparison.

To prove the assertion, consider a countersequence along which the
suprema over such charts tend to zero. The exact compactness
alternative used in Theorem~\ref{thm:exact-spectral-gap} gives a
smooth closed limit or a marked smooth ALE limit with a positive
exact same-type completion. In each case the closed limit class
has nonzero energy in both chiralities by the closed gap and
signature. A smooth Bach-flat four-metric has zero obstruction
tensor. It admits local analytic conformal representatives
\cite[Theorem~1.3(c)]{LassasLiimatainen2019ConformalHarmonic}.
In these representatives each chiral Weyl block is analytic,
so vanishing on an open set propagates by the analytic identity
theorem through overlapping charts. Hence each nonzero locus is
open dense, and their intersection contains a small normal chart
with positive overlap capture. For the ALE branch take this chart
away from the filled point and use the metric before completion.
In either branch choose the normal scale with error below
\(\varepsilon/2\) on \(\overline B_2\). Its whole fixed
compact chart transfers to the late sources. For the ALE branch
the source scale includes the constant factor of the pointed
rescaling. Smooth convergence and continuity of the top
eigenvalues give a strictly positive limiting overlap on the
transferred disk, contrary to the countersequence. Halving the
positive infimum of the suprema gives one suitable \(a\).
In the strict deficit sublevel the ALE branch is excluded.
In a closed subsequential limit the full-band supremum is at
least \(2a\). Choose there a chart with capture strictly above
\(a\) and a fixed positive scale. Its transfer rules out a
sequence with no chart of scale \(r\ge i^{-1}\) and capture
at least \(a\), giving the joint lower bounds.

Writing \(D_\pm=\int(\lambda_1^\pm-\lambda_2^\pm)^2dV\) as in
Section~\ref{sec:endpoint-charge}, trace freeness gives
\[
 Y^2=48e-36\mathcal O(M)-12(D_++D_-)-\Delta_{\mathrm{sp}} .
\]
Thus in the exact band \(Y^2\le48e-\Delta_{\mathrm{sp}}-36a\).
In particular no sequence in that band has vanishing chiral
overlap. The bound holds at a class attaining \(m_\delta\).
Extremality there is through exact classes and supplies no
unrestricted variation of the spectral deficit.

\subsection{Volume and curvature at the round exit}
\label{sec:round-joint-exit}

Suppose \(g_i\) are volume-one positive Yamabe metrics on one
nonstandard oriented smooth homotopy four-sphere with
\(e_i\to I^+(M)\), \(e_i\le2T-\delta\), and
\(\Delta_{\mathrm{sp}}([g_i])\to0\). Suppose the whole
marked filled-base comparison in
Theorem~\ref{prop:local-min} and the subsequent cap argument
is given, including its strong first-order convergence and local
full \(L^2\) curvature bound. The spectral transfer makes the
filled base \((N,h)\) conformally round. Denote its finitely many
ordinary marks by \(P\), and by \(s\) the single exterior label
retaining both limiting chiral charges. Let \(q_i\) be the
unit-volume modified Yamabe metric in \([g_i]\).

\begin{proposition}
\label{prop:round-joint-exit}
In these comparisons both \(g_i\) and \(q_i\) have volume tending
to zero on every fixed compared compact. At each fixed small
cutting radius, the actual exterior at \(s\) carries all limiting
volume in both gauges, both chiral integrals and the spectral
integral. If the comparison metrics converge smoothly on
compacts, this same exterior also carries all limiting full
curvature in both gauges.
\end{proposition}

\begin{proof}
On a subsequence put \(Y_i\to y\) and \(S=8\sqrt6\pi\).
The sharp spectral bound \(\sigma_i\le4e_i/3\) and the strict
energy band give
\[
 0<48\delta\le y^2=36\lim_i\sigma_i\le48I^+(M)\le S^2-48\delta .
\]
On the comparison write \(\Phi_i^*g_i=u_i^2m_i\).
Use \(\mathcal L_m=-6\Delta_m+R_m\) and \(Q_m\) for its
quadratic form, with \(\Delta_m=\operatorname{div}_m\nabla_m\).
Strong first-order convergence bounds \(Q_{m_i}-Q_h\)
on each fixed compact by \(o(1)\) times the squared \(H^1\)
norm. Indeed the scalar density is the divergence of a term
linear in the first metric derivatives, plus a term quadratic
in them, with smooth metric coefficients. From the Yamabe
equation and a fixed squared cutoff one obtains
\[
 Q_{m_i}(\chi u_i)
   =Y_i\int \chi^2u_i^4dV_{m_i}
       +6\int u_i^2|d\chi|_{m_i}^2dV_{m_i} .
\]
Hence \(u_i\) is locally bounded in \(H^1\), and any weak
limit \(u\ge0\) solves \(\mathcal L_hu=yu^3\) with
\(\|u\|_{4,h}\le1\). The equation extends across \(P\).
Testing first with squared puncture cutoffs gives finite
energy there, since the \(L^4\) norms of their gradients are
bounded and the \(L^4\) tails of \(u\) vanish. The points
then have zero \(H^1\) capacity. Testing by \(u\) gives
\(S\|u\|_4^2\le y\|u\|_4^4\), which forces \(u=0\).
In the cutoff identity, the nonlinear term is bounded by
\((y+o(1))\|\chi u_i\|_{4,m_i}^2\) by unit volume and
Hölder. On small fixed coordinate balls the sharp Sobolev
inequality has coefficient \(S-o(1)>y\), with the error made
small by the chart radius
\cite[Theorems 3.2 and 3.3]{LeeParker1987Yamabe}.
The \(L^2\) cutoff terms tend to zero by local Rellich convergence.
Absorption gives strong local \(H^1\cap L^4\) convergence
to zero.

On the actual exits, cutoffs which transition on fixed
compared annuli now have physical squared gradient integral
tending to zero by Hölder and conformal invariance of the
fourth-power gradient integral. The source Yamabe inequality
forces \(y\sqrt{\theta_j}\le y\theta_j\) for each limiting
exit fraction \(\theta_j\). Their sum is one, so one exit has
unit limiting volume. Fix the cap radius and buffer, and place
the cutoff transition deeper into that raw exit than the
innermost splice. Its entire support compares conformally on the
unchanged part of the cap, giving a Yamabe upper limit at most
\(y\). If this exit were a label other than \(s\), its positive
cap would have vanishing chiral energy and hence Yamabe constant
tending to \(S\) by Gauss--Bonnet. Thus it is \(s\).
Take the source limit at these fixed choices before the cap
exhaustion. Changing between fixed small cutting radii
only adds or removes an evacuated compared annulus.

For \(q_i\), put \(\beta_i=R_{q_i}-6L_{q_i}\le0\), the modified
Yamabe value
\cite[Lemma~2.1 and Theorem~1.7]{ChenZhu2014Conformal}.
The chain \(Y_i\le\int R_{q_i}dV_{q_i}
   \le\beta_i+6\sqrt{\sigma_i}\) implies
\(\beta_i\to0\) and \(R_{q_i}-Y_i\to0\) in \(L^2(q_i)\).
It also gives \(L_{q_i}\to y/6\) in \(L^2(q_i)\), by
equality in the limiting unit-volume Cauchy--Schwarz inequality.
Consequently \(L^2dV-(y^2/36)dV_{q_i}\to0\) in total variation.
Both chiral masses vanish locally in the round comparison and
on the other exteriors. They bound the invariant spectral
measure, so the modified volume is also evacuated from these
sets and concentrates at \(s\).

Assume now \(m_i\to h\) smoothly on compacts. Write
\(\Phi_i^*q_i=v_i^2m_i\).
The ordinary factor \(u_i\) converges locally uniformly to zero
by local Sobolev iteration in its equation, using the strong
\(L^4\) convergence to absorb the \(L^2\) potential \(Y_i u_i^2\).
The modified equation is
\[
 \mathcal L_{m_i}v_i=(6L_{m_i}+\beta_i v_i^2)v_i .
\]
Since \(\beta_i\le0\), \(v_i\) is a positive subsolution of
\(-6\Delta_{m_i}v_i\le C v_i\) on buffered compacts, with a
uniform local \(C\). Its \(L^4\) convergence to zero then gives
local uniform convergence. At a fixed \(x_0\in N-P\), normalize
the two factors by their values at \(x_0\). Harnack and interior
elliptic estimates give subsequential \(C^2_{\mathrm{loc}}\)
limits, positive and \(\mathcal L_h\)-harmonic
\cite[Sections 6.1 and 8.5--8.8]{GilbargTrudinger2001Elliptic}.
For this regularity the normalized factors first have local
\(W^{2,p}\) bounds for finite \(p\), since their linear potentials
are bounded. The top eigenvalues in smoothly moving frames
are locally Lipschitz with bounded constants. Together with the
\(W^{2,p}\) bounds, the equations then give interior
\(C^{2,\alpha}\) bounds for \(0<\alpha<1\).

To identify the positive limits, take a cutoff equal to one
away from \(s\) and turning off on a fixed compared annulus in
that exit. Its support has vanishing volume in both gauges.
In the modified gauge, the source Yamabe inequality implies
\[
 \left(Y_i-\|R_{q_i}-Y_i\|_{2,q_i}
  -Y_i\sqrt{\operatorname{Vol}_{q_i}(\operatorname{supp}\chi)}\right)
 \|\chi\|_{4,q_i}^2
 \le6\int|d\chi|_{q_i}^2dV_{q_i}=O(v_i(x_0)^2).
\]
The ordinary test replaces \(v_i(x_0)\) by \(u_i(x_0)\)
and has zero scalar error. The positive lower coefficient
bounds the volume after the corresponding constant rescaling
on all the parts where the cutoff is one. After interior
convergence and exhaustion, either limiting factor is locally
\(L^4\) across the other marks. In flat conformal coordinates
it extends there as a harmonic function after multiplication
by the conformal factor of \(h\). To see removability, test
against a smooth function times a shrinking cutoff. The
second-derivative error has \(L^{4/3}\) norm \(O(r)\), so the
weak Laplace equation passes through the mark. Stereographic
projection from \(s\) now gives a positive harmonic function
on \(\mathbb R^4\). It is constant, with its value fixed by
normalization. Both limits therefore give the same flat
rescaled metric on the compared punctured base. The squared
curvature integrals on its fixed compared compacts tend to zero
by convergence and constant-scale invariance.

Finally, at a fixed small face around any \(j\ne s\), the
entire actual exterior has \(\chi=1\), since attaching its
standard ball cap gives a homotopy four-sphere. Its squared
scalar and Weyl integrals vanish by the volume and chiral
accounting above. The Euler boundary term converges to
\(8\pi^2\): the comparison after constant rescaling is \(C^2\)
and flat across the filled mark, and the outward normal from
that exterior points away from the puncture, as for a filled
disk. The boundary term is invariant under constant scaling
and continuous under \(C^2\) convergence on the collar.
In the convention \(8\pi^2\chi\) this is the boundary
integral of the flat filled disk
\cite[Section 7, (7.4)]{TianViaclovsky2005ALE}.
Gauss--Bonnet on the entire exterior now gives
\(\int|Z|^2dV\to0\). The same formula for the
\(C^{2,\alpha}\) modified metric follows by \(C^2\) smooth
approximation on each fixed source. The full curvature
decomposition gives the asserted allocation in both gauges.
\end{proof}

The common exit and its flat round-side model control the limiting
accounting. A strict restart comparison still requires the
finite-source difference in \eqref{eq:signed-restart} to exceed
the minimizing excess. The exact graft of
Proposition~\ref{prop:graft} has equality in that difference.


\section*{Provenance and assistance}

Prepared for Edward B. Baker III on 25 November 2026.
The branch input revision is
\texttt{483ad3cef5a0af4328c5d35cabbc1f5224d26b1e}.
The shared physical-volume analysis is read from the master
source. The independent repair and complement comparisons,
normal-ball calculation, closed variation, restart model,
marked replacement, round-end form and Einstein restrictions
draw on master packets E03, E17, E20, E21, E39, E40 and
E44--E48. GPT-6.2 Sol with ultra effort was requested for
the drafting packet. Actual serving model and effort were
unavailable for independent verification. OpenAI Codex
assisted with the derivations, their checks and the preparation
of this draft. The endpoint analysis uses the working E52 master edition and the E53--E55 derivations and independent checks of 25 November 2026.

\section{Outstanding challenges}

Theorem~\ref{thm:endpoint-smoothing} settles the former endpoint.
Theorem~\ref{thm:spectral-windows} supplies charged disks for every
\(I<2T\). The median construction in Theorem~\ref{thm:band-charge}
also retains essentially all modified volume when its margin holds. A collar on this same charged family satisfying the
determinant-normalized comparison would permit a measured improvement.
The strict comparison \eqref{eq:removed} leaves a geometric
problem on the measured source. One needs actual marked collars
whose aggregate error is smaller than the discarded raw energy
minus the minimizing excess, with the joined path satisfying
\(e_i+c_i<8\pi^2\). The joint selection supplies
marked disks and their chiral measures, while the transition to
first-order collar control in a physical or conformal comparison
is still open. For an absolute recognition bound using
\eqref{eq:residual-test}, the residuals in \eqref{eq:complement}
and the aggregate costs may together have lower limit at most \(T\) by the uniform gap.
For possible concentration at separated physical energy scales,
the whole charged sets must lie on opposite sides of a supplied
short cap collar. A first-order chart without that separating
information does not furnish this comparison.

The whole marked smooth limit is itself an additional premise.
In the closed case, the problem is classification or strict
instability of an attained Bach-flat class beyond
the Einstein and bounded-Ricci tests. In the round case,
\eqref{eq:signed-restart} requires a strict estimate on a
genuinely changed retained class, or a different invariant
detecting a physical charged region. Optimizing only the limiting
Yamabe constant leaves exact equality on the evacuated model of
Proposition~\ref{prop:graft}.

Eliminating a minimizing alternative by fixed-type descent needs
a single positive same-type competitor with energy below \(I\).
Direct recognition or an independent rigidity contradiction can
serve the same purpose.
An iterative procedure instead needs a termination mechanism,
such as a uniform decrement of a bounded-below quantity or a
decrease of a well-founded discrete complexity. A sequence of
strictly smaller energies can converge above the recognition
threshold, and a sequence of same-energy restarts can repeat
exactly. The initial existence of a positive subcritical class
on every smooth homotopy four-sphere also remains a separate
problem. The exact round boundary form tests a supplied
retained metric for that purpose.

This revision was prepared for Edward B. Baker III with OpenAI Codex
assistance on 25 November 2026. It starts from the assessed E63--E67 manuscript generations
recorded at \texttt{5706ef811adc}. The E69--E72 research assessment
records the manuscript generations, independent readers and evidence
for the results used here.

\bibliographystyle{unsrturl}
\bibliography{references}
\end{document}
